\documentclass{unicam_thesis}
\usepackage[utf8]{inputenc}
\usepackage{listings}
\usepackage[italian]{babel}
\usepackage{placeins}
\usepackage{booktabs}
\usepackage{tabularx}
\usepackage{enumitem}
% Elenchi compatti: la classe e' a due facciate e allungherebbe gli spazi fra le voci
\setlist[itemize,enumerate]{itemsep=2pt, parsep=0pt, topsep=4pt}
\usepackage{textcomp}
\usepackage{xurl}
\addto\captionsitalian{\renewcommand{\lstlistingname}{Listato}}

% Posizionamento dei float: figure e tabelle restano nel testo invece di finire da sole
% su una pagina quasi vuota
\renewcommand{\topfraction}{0.85}
\renewcommand{\bottomfraction}{0.7}
\renewcommand{\textfraction}{0.1}
\renewcommand{\floatpagefraction}{0.75}
\setcounter{topnumber}{3}
\setcounter{bottomnumber}{2}
\setcounter{totalnumber}{4}
% Pagine di soli float: allineate in alto, con spazio fisso fra un float e l'altro
\makeatletter
\setlength{\@fptop}{0pt}
\setlength{\@fpsep}{14pt}
\setlength{\@fpbot}{0pt plus 1fil}
\makeatother

\newcommand{\ldb}{HLK-LD2410B}
\newcommand{\ldc}{HLK-LD2420}
\newcommand{\pir}{HC-SR501}
%%%%%%%%%%%%%%%%%%%%%%%%%%%%
% TESI DATI FRONTESPIZIO
%%%%%%%%%%%%%%%%%%%%%%%%%%%%

\title{Analisi sperimentale di sensori radar mmWave e PIR per il rilevamento di persone in scenari post-sisma}

\university{Universit\`a degli Studi di Camerino}%
\school{Scienze e Tecnologie}%
\course{Laurea in Informatica (Classe L-31)}%


\author{\shortstack[l]{{\mdseries\small Laureando}\\ Stefano De Bernardinis}}%
\advisor{\shortstack[l]{{\mdseries\small Relatore}\\ Prof. Massimo Callisto De Donato}}%
\academicyear{2025/2026}%
\studentid{\\ 119777}%

%%%%%%%%%%%%%%%%%%%%%%%%%%%%
% FINE DATI FRONTESPIZIO
%%%%%%%%%%%%%%%%%%%%%%%%%%%%

\theoremstyle{definition} \newtheorem{esempio}{Esempio}[chapter]
\theoremstyle{definition}
\newtheorem{definizione}{Definizione}[chapter] \theoremstyle{plain}
\newtheorem{teorema}{Teorema}[chapter]
\counterwithout{equation}{chapter}% equazioni numerate (1), (2)... in tutta la tesi
\graphicspath{{Screenshot/},{Immagini/},{API/},{Source/},{figures/}}
%\bibliography{biblio.bib}
\addbibresource{biblio.bib}

\begin{document}

\maketitle

\chapter*{Abstract}
\label{chap:Abstract}

La tesi analizza e confronta due tecnologie per il rilevamento della presenza umana in ambienti indoor, il sensore PIR (Passive InfraRed) e il radar mmWave a 24 GHz, per stabilire quale delle due rilevi una persona con maggiore affidabilità in uno scenario post-sisma. Il contesto applicativo in cui il confronto trova riscontro è il progetto DIPME, che prevede l’integrazione di sensori IoT (Internet of Things) negli arredi scolastici per segnalare, in seguito a un evento sismico, l’eventuale presenza di persone rifugiate sotto di essi.

Il nodo di quel progetto affida oggi il rilevamento della presenza a un sensore PIR. Tuttavia, a causa del suo principio di funzionamento, il PIR rileva il movimento ma non una persona immobile, cioè proprio il caso che il progetto deve coprire. Il lavoro verifica quindi se il radar, di tipo FMCW (Frequency Modulated Continuous Wave), possa superare questo limite e con quali compromessi.

Nella prima parte vengono studiati il sensore PIR \pir{} e i due moduli radar \ldb{} e \ldc{}, con particolare attenzione ai loro principi di funzionamento, ai protocolli di comunicazione, ai parametri disponibili e ai limiti riportati nella documentazione dei costruttori. Per il confronto è stata realizzata una piattaforma di acquisizione basata su ESP32, in grado di campionare i sensori in modo sincrono e di accedere, nel caso dei radar, all’energia riflessa per ciascuna cella di distanza. A supporto della campagna è stata progettata e implementata un’interfaccia web, servita direttamente dall’ESP32, per la visualizzazione dei dati in tempo reale e la loro esportazione nello stesso formato usato dalla piattaforma di acquisizione.

La piattaforma è stata utilizzata per una campagna sperimentale di 444 trial che ha preso in esame la persona ferma e in movimento, la distanza di rilevamento, le latenze, i falsi positivi, l’attenuazione del segnale in presenza di ostacoli, la copertura angolare e il rilevamento del respiro. I risultati confermano il limite del PIR. Con la persona immobile, anche sotto il banco, il PIR l’ha rilevata in meno del 2\,\% del tempo, mentre il radar in tutti gli scenari di confronto. Anche il secondo radar ha confermato questo comportamento, con limiti che la tesi attribuisce all’esemplare in prova.

Sulle grandezze graduate fornite dal radar è stato inoltre sviluppato un indice di vitalità calcolato direttamente a bordo, che, oltre alla presenza, stima quanto la persona si muove. Il consumo energetico dei sensori è stato infine valutato sulla base dei datasheet. Poiché il radar consuma circa mille volte più del PIR, la tesi conclude che il radar deve affiancare il PIR e non sostituirlo, restando spento nel funzionamento ordinario e acceso quando il nodo entra nel regime di emergenza.


\tableofcontents
%\lstlistoflistings
%\listoffigures
%\listoftables
\chapter{Introduzione}
\label{chap:intro}

Negli ultimi anni, i sistemi IoT si sono diffusi anche nell'ambito della protezione civile e della gestione delle emergenze, dove la disponibilità di informazioni tempestive e affidabili può incidere direttamente sull'efficacia dei soccorsi.
Il terremoto del 2016 che ha colpito l'Italia centrale, e in particolare l'area di Camerino, ha messo in evidenza questa esigenza. Infatti, nelle prime ore successive a un evento sismico, per i soccorritori è fondamentale sapere dove sono le persone, soprattutto quando potrebbero trovarsi in un luogo difficilmente accessibile o non essere in grado di muoversi. Da qui nasce la domanda alla base della tesi. Quale sensore, tra il PIR a infrarossi passivi e il radar a onde millimetriche (mmWave), rileva la presenza di una persona con maggiore affidabilità, e il radar può sostituire il PIR nei sistemi che oggi lo usano o deve affiancarlo?

Il capitolo è organizzato come segue. La Sezione~\ref{sec:dipme} descrive il progetto DIPME, contesto applicativo del lavoro, la Sezione~\ref{sec:obbiettivo} definisce gli obbiettivi della tesi e la Sezione~\ref{sec:struttura} ne illustra la struttura.


\section{Il progetto DIPME}
\label{sec:dipme}

Il confronto tra i due sensori trova uno scenario di applicazione concreto nel progetto DIPME, che ha come obbiettivo la mitigazione del rischio sismico attraverso l'integrazione di sensoristica IoT negli arredi degli ambienti scolastici \cite{slideDipme}. La tesi ne riprende il problema e ne valuta una possibile estensione. 
L'idea del progetto si articola su tre livelli:

\begin{itemize}
    \item \textbf{Arredi salva-vita.} Banchi e pareti sono progettati dalla Scuola di Architettura e Design di UNICAM per resistere al crollo del solaio e mantenere al di sotto di sé un volume in cui una persona possa trovare riparo (Figura~\ref{fig:arredo-render}). Non si tratta quindi di semplici mobili rinforzati, ma di elementi strutturali pensati per lavorare insieme e vincolati tra loro
    \cite{pietroni2021furniture}.

    \item \textbf{Nodi di sensing.} All'interno dell'arredo è integrata l'elettronica della piattaforma DIPME, sviluppata con partner industriali, fra cui AM Microsystems~\cite{slideDipme}, nell'ambito del progetto VITALITY~\cite{callisto2026dipme}. Ogni nodo, il DIPME-DEVICE, è un'unità autonoma alimentata a batteria e integra una radio LoRa, un sensore di presenza PIR e alcuni sensori ambientali per la misura di temperatura, umidità, pressione e anidride carbonica.

    \item \textbf{Raccolta e visualizzazione.} I nodi comunicano via LoRa con un coordinatore, che inoltra i dati a un gateway locale. LoRa permette di trasmettere a basse potenze anche su distanze dell'ordine di 3--5~km in ambiente urbano~\cite{slideDipme}. Il gateway, alimentato da un gruppo di continuità, coordina il comportamento dei nodi e trasmette i dati a una piattaforma cloud con cruscotti di monitoraggio, conservandoli quando la connessione si interrompe e reinviandoli al ripristino~\cite{callisto2026dipme}. La piattaforma SAFE, di cui DIPME estende l'approccio, prevede inoltre un sottosistema di localizzazione basato su droni e dispositivi mobili a supporto dei soccorritori~\cite{callisto2024safe,callisto2026dipme}.
\end{itemize}

\begin{figure}[htbp]
  \centering
  \includegraphics[width=0.9\textwidth]{Arredo_salva_vita}
  \caption{Il principio dell'arredo salva-vita. Il banco (a sinistra) e la parete attrezzata (a destra) conservano un volume di riparo sotto il carico del crollo. Il nodo di sensing, fissato sotto il piano, osserva l'occupante~\cite{slideDipme}.}
  \label{fig:arredo-render}
\end{figure}

Il sistema opera in due modalità, un tempo di \emph{pace}, dedicato al monitoraggio ordinario, e un tempo di \emph{guerra}, associato alla gestione dell'emergenza sismica. Il passaggio dalla prima alla seconda è comandato dal gateway al segnale di un accelerometro strutturale. Se il gateway non conferma i messaggi per tre ore, il nodo passa da solo al regime di emergenza. È in questa situazione che il rilevamento della presenza assume un'importanza particolare, perché da questo dato dipende la cadenza delle trasmissioni del nodo. Con una persona rilevata il nodo invia un messaggio al minuto, senza presenza uno ogni quindici minuti. Una persona sotto un banco che il sensore non rileva verrebbe quindi segnalata come assente, e il suo banco trasmetterebbe quindici volte meno spesso, proprio durante l'emergenza. DIPME è in esercizio in un dimostratore con 42 nodi installati nei banchi di due aule di una scuola secondaria delle Marche~\cite{callisto2026dipme}.

\subsection{Il limite del sensore adottato}
\label{sec:limite-pir}

Il PIR, la soluzione più diffusa per questa funzione grazie ai consumi e ai costi ridotti, ha però un limite che non dipende dalla qualità del componente, ma dal suo principio di funzionamento, discusso nel Capitolo~\ref{chap:sensori}. Il sensore risponde infatti alla variazione del flusso infrarosso, non al suo valore, e quindi non è in grado di rilevare una persona immobile. 

Nello scenario DIPME questo limite si presenta proprio nella situazione più critica, perché una persona rifugiata sotto un banco dopo un crollo è, tipicamente, ferma. Gli autori del progetto DIPME indicano, tra gli sviluppi futuri, l'estensione del sistema di sensing con radar mmWave a complemento del PIR, da valutare in termini di accuratezza, consumo, privacy e complessità \cite{callisto2026dipme}. Un radar a 24 GHz, sensibile ai micro-movimenti e in grado di misurare la distanza del bersaglio, può in linea di principio superare il limite del PIR.


\section{Obbiettivi}
\label{sec:obbiettivo}

La tesi risponde alla domanda posta in apertura con un confronto sperimentale fra un sensore PIR e due moduli radar mmWave a 24 GHz, l'\ldb\ e l'\ldc, basato su misure ripetibili che stabiliscono anche quanto sia ampio il divario fra le due tecnologie. Il risultato interessa direttamente il progetto DIPME.

Il lavoro è stato organizzato in sei punti:

\begin{itemize}
    \item \textbf{O1. Studiare i sensori}, per capirne il funzionamento, gli impieghi e i dati che producono.

    \item \textbf{O2. Analizzare i dati forniti}, individuando quali grandezze ciascun sensore mette a disposizione e con quale accuratezza.

    \item \textbf{O3. Confrontare sperimentalmente} PIR e mmWave con metriche quantitative.

    \item \textbf{O4. Valutare il consumo energetico}, a titolo informativo e sulla base dei datasheet, non essendo disponibile la strumentazione per misurarlo.

    \item \textbf{O5. Realizzare un'interfaccia web} che mostri i dati in tempo reale, li salvi e li esporti in formato CSV per l'analisi numerica.

    \item \textbf{O6. Definire un indice di vitalità}, calcolato a bordo del nodo, che oltre alla presenza misuri quanto una persona si muove.
\end{itemize}

Gli altri elementi del progetto DIPME (la radio LoRa, il gateway e la piattaforma cloud) e il sottosistema di localizzazione con i droni della piattaforma SAFE non sono oggetto della tesi e compaiono nella Sezione~\ref{sec:dipme} solo come contesto.

\section{Struttura della tesi}
\label{sec:struttura}

La tesi è organizzata come segue.

\begin{itemize}
    \item \textbf{I sensori di presenza.} Il Capitolo~\ref{chap:sensori} spiega come funzionano il PIR e il radar FMCW a 24 GHz e li colloca rispetto alle alternative (O1).
    
    \item \textbf{Analisi dei componenti hardware e dei dataset prodotti.} Il Capitolo~\ref{chap:moduli} presenta i tre sensori impiegati a partire dalla documentazione ufficiale, e la piattaforma di acquisizione su ESP32 con il formato dati comune a tutta la tesi (O2).
    
    \item \textbf{Metodologia di confronto.} Il Capitolo~\ref{chap:metodologia} fissa il protocollo, le metriche e le convenzioni con cui i sensori sono stati confrontati, e presenta l'interfaccia web servita direttamente dall'ESP32, senza alcuna rete esterna, realizzata a supporto della campagna (O5).

    \item \textbf{Confronto sperimentale PIR vs mmWave.} Il Capitolo~\ref{chap:confronto} è il cuore della tesi e raccoglie i risultati della campagna di misura, con ogni risultato accompagnato dalla sua dispersione su ripetizioni indipendenti (O3).

    \item \textbf{L'indice di vitalità.} Il Capitolo~\ref{chap:vitalita} definisce, tara e valida l'indice di vitalità, ne descrive il calcolo a bordo del microcontrollore e chiude con la stima del ritmo respiratorio, verificata contro un ritmo imposto (O6).

    \item \textbf{Consumo energetico.} Il Capitolo~\ref{chap:consumi} stima i consumi dei tre sensori e l'autonomia del nodo a batteria, e ne ricava i vincoli per il nodo del progetto (O4).

    \item \textbf{Conclusioni e sviluppi futuri.} Il Capitolo~\ref{chap:conclusioni} risponde alla domanda di partenza e indica come proseguire il lavoro.
\end{itemize}
\chapter{I sensori di presenza}
\label{chap:sensori}

Questo capitolo introduce i principi fisici alla base delle tecnologie di rilevamento considerate, analizzandone il funzionamento e le caratteristiche principali. L'analisi permette di comprendere come i diversi fenomeni fisici vengano sfruttati per rilevare la presenza di una persona e come questi si riflettano nel comportamento dei sensori.

In particolare, la Sezione~\ref{sec:pir} descrive il sensore PIR (Passive InfraRed), il suo principio di funzionamento e il motivo per cui non rileva un corpo immobile. La Sezione~\ref{sec:mmwave} presenta il radar FMCW a 24 GHz, il concetto di gate di distanza e il modo in cui rileva una persona ferma. La Sezione~\ref{sec:altre} richiama brevemente le altre tecnologie di rilevamento di presenza, non considerate nel confronto sperimentale. Infine, la Sezione~\ref{sec:confronto} mette a confronto PIR e mmWave e indica in quali scenari ciascuno dei due è preferibile.


\section{Il sensore PIR}
\label{sec:pir}

\subsection{Principio fisico}
\label{subsec:pir-principio}

PIR è l'acronimo di Passive InfraRed. Il sensore non emette nulla e si limita a osservare la radiazione infrarossa che lo raggiunge. Il cuore del componente è un cristallo piroelettrico, un materiale che genera una piccola carica elettrica quando la radiazione incidente cambia. Il corpo umano, con una temperatura di circa 36 °C (309 K), emette con un picco intorno a 9,4 µm, e il sensore è filtrato con una finestra ottica di 5–14 µm proprio per privilegiare quella banda \cite{niceraPyro}.

Il PIR risponde quindi alle variazioni della radiazione, non alla sua quantità in un dato istante. Costruttivamente, il sensore contiene due elementi piroelettrici disposti in configurazione differenziale, collegati in serie con polarità opposte. Da questa configurazione seguono due conseguenze fondamentali (Figura~\ref{fig:pir-detection}):

\begin{itemize}
    \item una persona che si muove investe i due elementi in tempi diversi, generando una differenza tra i segnali prodotti e quindi un rilevamento;

    \item una persona ferma illumina entrambi gli elementi con un flusso infrarosso pressoché costante. I due contributi si compensano e il segnale differenziale tende a zero, impedendo il rilevamento.
\end{itemize}

\begin{figure}[!htbp]
  \centering
  \includegraphics[width=0.8\linewidth]{Pir_detection}
  \caption{Sensore piroelettrico a due elementi in configurazione differenziale. Una persona che attraversa le due zone le investe in tempi diversi e produce un segnale prima positivo e poi negativo; una persona ferma in una zona illumina entrambi gli elementi allo stesso modo e il segnale differenziale resta nullo.}
  \label{fig:pir-detection}
\end{figure}

Lo stesso principio spiega la buona immunità del PIR alle variazioni lente e globali della scena, come il sole che riscalda l'ambiente. Queste interessano entrambi gli elementi allo stesso modo e si cancellano nel segnale differenziale. Il rovescio della medaglia è che il segnale utile dipende dal contrasto termico fra corpo e sfondo. Con una temperatura ambiente di 20~°C il contrasto è di circa 16~°C, mentre a 30--32~°C si riduce a 4--6~°C, con segnale più debole e portata ridotta \cite{hcsr501}.

\subsection{La lente di Fresnel}
\label{subsec:fresnel}

Due soli elementi coprirebbero un campo visivo strettissimo. La cupola bianca che si vede su tutti i moduli PIR è una lente di Fresnel multi-segmento. Ogni segmento focalizza sul cristallo una porzione differente della scena, e l'insieme crea decine di zone di rilevamento disposte a ventaglio e separate da zone cieche
\cite{adafruitPIR}, come mostra la Figura~\ref{fig:zone-fresnel}.

\begin{figure}[!htbp]
  \centering
  \includegraphics[width=0.6\linewidth]{Zone_lente_fresnel}
  \caption{Zone di rilevamento create dalla lente di Fresnel, focalizzate alternativamente sui due elementi e separate da spazi ciechi. }
  \label{fig:zone-fresnel}
\end{figure}

Questo dettaglio costruttivo ha due conseguenze dirette sul tipo di movimento che il sensore è in grado di rilevare:

\begin{enumerate}
    \item Il movimento trasversale, che attraversa le zone, genera molte transizioni di flusso ed è rilevato bene. Il movimento radiale, diretto verso il sensore, resta più a lungo nella stessa zona ed è rilevato male.

    \item La portata dichiarata, tipicamente 3--7~m, vale per il movimento trasversale di tutto il corpo. Un movimento sul posto attraversa poche zone, e sempre meno al crescere della distanza, perché le zone si allargano a ventaglio.
\end{enumerate}

Queste due previsioni vengono messe alla prova nel Capitolo~\ref{chap:confronto}.

\subsection{I dati prodotti}
\label{subsec:pir-dati}

L'uscita di un modulo PIR è un solo bit digitale, alto quando il modulo dichiara movimento e basso altrimenti. Il passaggio dal segnale analogico del cristallo a questo bit avviene interamente a bordo del modulo, in un circuito integrato di condizionamento che amplifica il segnale, lo confronta con una soglia e tiene alta l'uscita per un tempo di ritenuta regolabile. Il microcontrollore che legge il sensore riceve quindi una decisione già presa e non ha accesso al segnale originario. I parametri di questo circuito nel modulo impiegato in questa tesi sono descritti nel dettaglio nel Capitolo~\ref{chap:moduli}.

Fa eccezione la modalità di trigger, selezionabile con un ponticello, perché determina il significato stesso del bit in uscita. Le modalità previste sono due (Figura~\ref{fig:pir-lh}) \cite{hcsr501}:

\begin{itemize}
    \item \textbf{Modalità L} (\textit{single trigger}). Il primo movimento alza l'uscita per il tempo di ritenuta e i movimenti successivi, finché l'uscita è alta, vengono ignorati. Allo scadere della ritenuta l'uscita torna bassa anche se la persona si sta ancora muovendo, e serve un nuovo trigger per rialzarla. Il bit si comporta come un monostabile a durata fissa e, di fatto, conta eventi.

    \item \textbf{Modalità H} (\textit{repeat trigger}). Ogni movimento rilevato durante la ritenuta la fa ripartire da zero. Finché il movimento continua l'uscita resta alta, e torna bassa solo quando il movimento cessa per un tempo pari alla ritenuta, per cui il sensore si avvicina di più a un rilevatore di presenza.
\end{itemize}

\begin{figure}[!htbp]
  \centering
  \includegraphics[width=0.9\linewidth]{Pir_L_H}
  \caption{Uscita del PIR nelle due modalità di trigger a parità di movimenti.}
  \label{fig:pir-lh}
\end{figure}

\FloatBarrier
\subsection{Chi lo usa e perché domina}
\label{subsec:pir-usi}

Il PIR è il sensore di presenza più diffuso al mondo e da decenni equipaggia luci automatiche, antifurti, videocitofoni e termostati. La sua diffusione si spiega con diversi fattori:

\begin{itemize}
    \item \textbf{Consumo:} la corrente di riposo è inferiore a 50~$\mu$A, compatibile con anni di funzionamento a batteria \cite{hcsr501}.

    \item \textbf{Costo:} un modulo completo di lente e circuito di condizionamento costa pochi euro.

    \item \textbf{Semplicità:} il modulo fornisce direttamente un bit di rilevamento e non richiede alcuna elaborazione a valle.

    \item \textbf{Natura passiva:} il sensore non emette alcuna radiazione, e questo semplifica sia le certificazioni sia la conformità in materia di privacy.
\end{itemize}


\subsection{Limitazioni}
\label{subsec:pir-limiti}

Gli stessi principi che rendono il PIR semplice ed economico ne fissano anche i limiti.

\begin{itemize}
    \item \textbf{Persona immobile:} il sensore risponde alla variazione del flusso infrarosso e non al suo valore, quindi una persona ferma non produce alcun segnale (Sezione~\ref{subsec:pir-principio}).

    \item \textbf{Tipo di movimento:} la lente divide la scena in zone separate da zone cieche, e il sensore rileva chi le attraversa. Chi si muove sul posto resta nella stessa zona e viene rilevato poco, se non da vicino (Sezione~\ref{subsec:fresnel}).

    \item \textbf{Contrasto termico:} il segnale dipende dalla differenza di temperatura fra corpo e sfondo, e in un ambiente caldo la portata si riduce (Sezione~\ref{subsec:pir-principio}).

    \item \textbf{Ostacoli:} il sensore deve vedere direttamente la scena, perché vetro e pannelli interposti ne bloccano la radiazione, come mostra il Capitolo~\ref{chap:confronto}.

    \item \textbf{Uscita a un bit:} il modulo dice solo se ha visto un movimento, con un impulso di durata fissa, e non dove si trova la persona né quanto si muove (Sezione~\ref{subsec:pir-dati}).
\end{itemize}

\section{Il radar mmWave a 24 GHz}
\label{sec:mmwave}

\subsection{Principio di funzionamento}
\label{subsec:fmcw}

I radar di presenza di fascia consumer sono radar FMCW (\textit{Frequency Modulated Continuous Wave}) che operano nella banda ISM intorno a 24~GHz, un'allocazione libera da licenza larga 250~MHz (24,00--24,25~GHz) e disponibile in Europa e negli Stati Uniti \cite{ituRR}.

Il nome con cui vengono commercializzati, radar mmWave, è in realtà improprio. La banda delle onde millimetriche (EHF) è definita fra 30 e 300~GHz, dove la lunghezza d'onda è compresa fra 10 e 1~mm \cite{ituRR}. A 24~GHz la lunghezza d'onda vale circa 12,4~mm, quindi questi moduli lavorano ancora nella banda SHF delle onde centimetriche. La denominazione deriva dal fatto che la stessa famiglia di sensori esiste anche a 60 e 77~GHz, che sono effettivamente millimetriche, ed è ormai quella adottata dai produttori. Nella tesi si mantiene per coerenza con le fonti citate.

Il principio di funzionamento si può descrivere in quattro passi \cite{tiFmcw}:

\begin{enumerate}
    \item Il trasmettitore emette un segnale continuo la cui frequenza sale in modo lineare nel tempo, per esempio da 24,00 a 24,25~GHz in qualche millisecondo. Questa rampa si chiama \textit{chirp} e si ripete in continuazione.

    \item Il segnale colpisce un oggetto e torna al ricevitore dopo un tempo di volo pari al doppio della distanza diviso la velocità della luce. Per un oggetto a 3~m il ritardo è di circa 20~ns.

    \item In quei 20~ns il trasmettitore ha continuato a salire di frequenza, quindi, nell'istante in cui l'eco arriva, il segnale trasmesso ha una frequenza leggermente più alta. Il ricevitore mescola i due segnali e ne ottiene la differenza, detta frequenza di battimento (Figura~\ref{fig:fmcw}). Poiché la rampa è lineare, questa differenza è proporzionale al ritardo, e dunque alla distanza dell'oggetto.

    \item Se nella scena ci sono più oggetti a distanze diverse, l'eco complessiva contiene una frequenza di battimento per ciascuno. Una trasformata di Fourier del segnale ricevuto le separa e restituisce, per ogni distanza, quanta energia è stata riflessa da lì.
\end{enumerate}

\begin{figure}[!htbp]
  \centering
  \includegraphics[width=0.75\linewidth]{Rilevamento_segnale_radar}
  \caption{Principio FMCW. La frequenza trasmessa cresce linearmente nel tempo, l'eco torna
  con un ritardo proporzionale alla distanza e la differenza di frequenza fra i due
  segnali, costante durante il chirp, misura quel ritardo.}
  \label{fig:fmcw}
\end{figure}

Il risultato di ogni chirp è quindi un profilo che associa a ciascuna distanza un'intensità di eco. La separazione minima fra due oggetti distinguibili dipende dalla larghezza di banda della rampa e, con i 250~MHz disponibili a 24~GHz, è di circa 60~cm \cite{tiFmcw}. È per questo che i moduli di questa classe suddividono l'asse delle distanze in intervalli di 70--75~cm, detti gate.

I moduli consumer non espongono il profilo grezzo. Per ciascun gate riportano a ogni frame un valore di energia riflessa, insieme alla decisione di presenza e alla distanza del bersaglio ricavate da quei valori. Il numero di gate, la scala dell'energia e i canali disponibili variano da modulo a modulo e sono descritti nel Capitolo~\ref{chap:moduli} per i due esemplari in esame.

\FloatBarrier
\subsection{Come si rileva un corpo immobile}
\label{subsec:corpo-immobile}

Un radar FMCW puro misura la distanza e non il movimento, quindi un corpo fermo produce comunque un'eco, come del resto pareti e arredi. In pratica, i moduli commerciali separano due componenti:

\begin{itemize}
    \item \textit{moving}, associata a variazioni rapide della fase dell'eco, cioè allo spostamento del bersaglio;

    \item \textit{stationary} (o \textit{motionless}), associata a un'eco stabile in distanza ma con fluttuazioni residue.
\end{itemize}

Una persona ferma non è mai realmente immobile, perché il torace si muove di alcuni millimetri a ogni atto respiratorio e, a 24~GHz, anche uno spostamento millimetrico produce una variazione di fase misurabile. È questa la ragione fisica per cui il radar può rilevare una persona immobile, a differenza del PIR. Lo stesso segnale, analizzato nel tempo, permette anche di stimare la frequenza respiratoria, come mostrato nel Capitolo~\ref{chap:vitalita}.

\subsection{Applicazioni e diffusione}
\label{subsec:mmwave-usi}

I radar mmWave di fascia consumer si sono diffusi negli ultimi anni in diversi ambiti, per esempio nell'illuminazione automatica di uffici e abitazioni, nell'automazione degli edifici per stimare l'occupazione degli ambienti, nel monitoraggio del sonno e nei sistemi di rilevamento delle cadute per l'assistenza agli anziani \cite{mmwaveOccupancy}. In ambito di ricerca gli stessi radar sono impiegati per il monitoraggio senza contatto dei parametri vitali, in particolare della frequenza respiratoria e cardiaca \cite{upm2024vitals,uts2024vitals,mdpi2024mmwaverm}. La loro diffusione si spiega con diversi fattori:

\begin{itemize}
    \item \textbf{Rilevamento del corpo fermo:} il modulo conferma la presenza anche di una persona immobile (Sezione~\ref{subsec:corpo-immobile}).

    \item \textbf{Informazione graduata:} oltre alla presenza il modulo riporta la distanza del bersaglio e un'energia riflessa per ogni gate, che permettono di sapere non solo se c'è qualcuno ma dove si trova e quanto si muove.

    \item \textbf{Insensibilità alle condizioni ambientali:} la misura non dipende dalla temperatura né dall'illuminazione, e le onde a 24 GHz attraversano vari materiali, per cui il modulo può essere nascosto dietro una copertura.

    \item \textbf{Privacy:} il radar non produce immagini e non permette di riconoscere una persona, e per questo è accettato in ambienti sensibili come scuole o camere da letto. I dati grezzi per gate permettono però di ricavare anche grandezze sullo stato fisico, come il respiro. In questa tesi sono esportati integralmente per l'analisi, mentre in un'installazione reale andrebbero trattati con le cautele del GDPR indicate dal progetto per gli ambienti scolastici~\cite{callisto2026dipme}.

    \item \textbf{Costo:} i moduli integrati, completi di antenne ed elaborazione a bordo, costano nell'ordine della decina di euro.
\end{itemize}

\subsection{Limitazioni}
\label{subsec:mmwave-limiti}

\begin{itemize}
    \item \textbf{Consumo:} un radar è un sensore attivo, che alimenta un trasmettitore e un'elaborazione continua del segnale, e assorbe decine di milliampere. Questo ne condiziona l'impiego nei dispositivi a batteria.

    \item \textbf{Più persone:} il modulo riporta un solo bersaglio per gate, quindi due persone nello stesso gate non si distinguono e viene riportata la più vicina, mentre l'altra resta nascosta (Sezione~\ref{sec:due-persone}).

    \item \textbf{Ostacoli:} le onde a 24~GHz attraversano legno, vetro e plastica, ma sono riflesse dal metallo, che non può stare davanti al modulo (Capitolo~\ref{chap:confronto}).

    \item \textbf{Riflessioni:} in un ambiente chiuso l'eco arriva anche per percorsi indiretti, da pareti e arredi, e può spostare la distanza riportata o mantenere la presenza per qualche secondo dopo che la persona è uscita (Sezione~\ref{sec:c3-5-3}).
\end{itemize}


\FloatBarrier
\section{Altre tecnologie di rilevamento di presenza}
\label{sec:altre}

Oltre al PIR e al radar mmWave esistono altre tecnologie in grado di rilevare la presenza di una persona. Di seguito se ne descrivono alcune, nessuna delle quali è stata presa in considerazione nel confronto sperimentale di questa tesi.

\textbf{UWB (Ultra-WideBand).} È una tecnologia radio che trasmette impulsi molto brevi su una banda larga alcuni gigahertz, tipicamente fra 3 e 10 GHz. La brevità degli impulsi dà un'ottima risoluzione temporale, e quindi spaziale, che la rende adatta soprattutto alla misura di distanza fra dispositivi. Usata come radar è in grado di rilevare corpi fermi e di attraversare i materiali, con consumi contenuti \cite{rfwireless}. I moduli disponibili forniscono però dati grezzi che richiedono un'elaborazione dedicata, e il loro impiego principale resta la localizzazione, non il rilevamento di presenza in un volume.

\textbf{Ultrasuoni.} Un trasduttore emette un impulso acustico e misura il tempo di ritorno dell'eco. Sono economici e precisi sulla distanza, ma il segnale è assorbito dai tessuti e dipende fortemente dalla geometria della scena, per cui la presenza umana viene rilevata in modo poco affidabile.

\textbf{CO$_2$.} La concentrazione di anidride carbonica cresce con il numero di persone in un ambiente ed è un buon indicatore della sua occupazione complessiva. Il gas però si diffonde lentamente, con tempi di risposta dell'ordine dei minuti, e la misura non dà alcuna informazione su dove si trovino le persone.

\textbf{Visione.} Telecamere e termocamere sono la soluzione più ricca di informazione, perché permettono di contare, localizzare e riconoscere le persone. Richiedono però un'elaborazione pesante, costano di più e producono immagini, e questo pone vincoli di privacy che negli ambienti scolastici e domestici ne limitano fortemente l'impiego.

 
 \section{PIR e mmWave a confronto}
\label{sec:confronto}

 Le caratteristiche elencate nelle Sezioni~\ref{subsec:pir-usi} e~\ref{subsec:mmwave-usi} mostrano che PIR e radar mmWave non sono alternative equivalenti, ma rispondono a esigenze diverse. Le principali differenze sono riassunte nella Tabella~\ref{tab:pir-mmwave}.

\begin{table}[htbp]
\centering
\caption{Confronto fra PIR e radar mmWave.}
\label{tab:pir-mmwave}
\small
\renewcommand{\arraystretch}{1.25}
\begin{tabularx}{\linewidth}{@{}l>{\raggedright\arraybackslash}X>{\raggedright\arraybackslash}X@{}}
\toprule
 & \textbf{PIR} & \textbf{Radar mmWave a 24~GHz} \\
\midrule
Principio & variazione del flusso infrarosso del corpo & eco di un'onda radio modulata (FMCW) \\
Tipo & passivo & attivo \\
Persona in movimento & sì, se attraversa le zone della lente & sì \\
Persona immobile & no & sì, dai micro-movimenti del respiro \\
Distanza del bersaglio & no & sì \\
Dati prodotti & un bit, impulso di durata fissa & presenza, distanza, energia per gate \\
Ostacoli & bloccato da vetro e pannelli & attraversa legno, vetro e plastica, non il metallo \\
Consumo & decine di \textmu A & decine di mA \\
Costo & pochi euro & circa una decina di euro \\
Privacy & nessuna immagine & nessuna immagine, ma dati sul respiro \\
\bottomrule
\end{tabularx}
\end{table}

Il PIR è la scelta giusta quando l'obbiettivo è rilevare il passaggio o l'attività di una persona e il dispositivo deve funzionare a batteria per lungo tempo. Il suo limite è strutturale e non dipende dal modello o dalla configurazione. Per il principio descritto nella Sezione~\ref{subsec:pir-principio} non rileva un corpo immobile, e per la geometria della lente descritta nella Sezione~\ref{subsec:fresnel} rileva male anche chi si muove senza spostarsi, se non da vicino.

Il radar mmWave è preferibile quando serve confermare la presenza di una persona immobile, oppure sapere dove si trova e quanto si muove. Il suo limite resta il consumo. Queste capacità e questo limite sono verificati sperimentalmente nel Capitolo~\ref{chap:confronto} e analizzati nel Capitolo~\ref{chap:consumi}.
\chapter{Analisi dei componenti hardware e dei dataset prodotti}
\label{chap:moduli}

Questo capitolo descrive i moduli sensore impiegati nel lavoro sperimentale, analizzandone le specifiche tecniche, le interfacce di comunicazione e i dati che forniscono. L'analisi permette di comprendere quali grandezze ciascun modulo mette effettivamente a disposizione e con quali limiti, e costituisce la base su cui sono state costruite la piattaforma di acquisizione e le analisi dei capitoli successivi.

Il capitolo è organizzato come segue. La Sezione~\ref{sec:c3-1} presenta il radar \ldb{}, il modulo principale della tesi, con le specifiche, la piedinatura, gli strumenti di configurazione, il protocollo seriale, la engineering mode che sblocca i dati per gate e l'identità dell'esemplare in prova. La Sezione~\ref{sec:c3-2} presenta con la stessa struttura il secondo radar, l'\ldc{}. La Sezione~\ref{sec:c3-3} descrive il PIR \pir{} e la configurazione adottata. La Sezione~\ref{sec:c3-4} illustra la piattaforma di acquisizione, dal repository di riferimento ai firmware, al formato dei dati e alla catena di analisi. La Sezione~\ref{sec:c3-5} verifica i dati prodotti e descrive i comportamenti del modulo da tenere presenti nel leggere i risultati. La Sezione~\ref{sec:c3-6} spiega perché le soglie dell'\ldb{} sono rimaste di fabbrica mentre l'\ldc{} ha richiesto una taratura.

\FloatBarrier
\section{Il radar \ldb{}}
\label{sec:c3-1}

L'\ldb{} è un modulo radar a 24~GHz per il rilevamento della presenza umana prodotto da Shenzhen Hi-Link Electronic \cite{hilinkManualV104}. Si presenta come una piccola scheda che integra il chip radar FMCW, le antenne a patch per trasmissione e ricezione, un microcontrollore che elabora il segnale e, nella variante B, un ricetrasmettitore Bluetooth per la configurazione da smartphone. L'elaborazione del segnale avviene interamente a bordo. Il modulo emette il segnale, analizza l'eco e comunica all'esterno soltanto il risultato, cioè lo stato di presenza, la distanza e l'energia del bersaglio, attraverso un connettore a cinque piedini che porta una linea seriale e un'uscita digitale. È in grado di distinguere un bersaglio in movimento da uno stazionario e di mantenere la presenza anche su una persona immobile, sfruttando i micro-movimenti del corpo descritti nella Sezione~\ref{subsec:corpo-immobile}. Il modulo è mostrato nella Figura~\ref{fig:c3-foto2410}.

\begin{figure}[!htb]
\centering
\includegraphics[width=0.7\linewidth]{Foto_LD2410B}
\caption{Il modulo \ldb{}, retro (in alto) e fronte (in basso). Sul fronte le due antenne a patch alle estremità e il chip radar al centro; sul retro il connettore a cinque piedini, il microcontrollore e i quarzi. Immagine tratta da \cite{studiopietersLD2410}.}
\label{fig:c3-foto2410}
\end{figure}

\subsection{Specifiche tecniche}
\label{sec:c3-1-1}

Le specifiche dichiarate dal produttore sono riassunte nella Tabella~\ref{tab:c3-1} \cite{hilinkManualV104,hilinkManualV103,ld2410bDatasheet}.

\begin{table}[!htb]
\centering
\caption{Specifiche dell'\ldb{} \cite{hilinkManualV104,hilinkManualV103,ld2410bDatasheet}.}
\label{tab:c3-1}
\small
\begin{tabularx}{\linewidth}{@{}lX@{}}
\toprule
\textbf{Parametro} & \textbf{Valore} \\
\midrule
Frequenza di lavoro & 24–24,25~GHz, banda ISM \\
Tecnologia & FMCW (onda continua modulata in frequenza) \\
Portata dichiarata & 5–6~m \\
Campo visivo & \(\pm\)60\(^\circ\) \\
Alimentazione & 5~V, alimentatore con capacità \(>\)~200~mA \\
Corrente media & 80~mA \\
Interfaccia & UART 256000~baud, 8N1 \\
Bluetooth & integrato (variante B), password di default HiLink \\
Dimensioni & circa 35~mm \(\times\) 7~mm \\
Gate di distanza & 9 (0–8) \\
Risoluzione gate & 0,75~m oppure 0,20~m \\
\bottomrule
\end{tabularx}
\end{table}

\subsection{Piedinatura e collegamento all'ESP32}
\label{sec:c3-1-2}

Il modulo espone cinque piedini, definiti nella Tabella 1 del datasheet \cite{ld2410bDatasheet}. La Tabella~\ref{tab:c3-2} riporta la funzione di ciascuno e il collegamento all'ESP32 adottato in questo lavoro. Come d'uso per una linea seriale, i due fili sono incrociati, cioè il piedino di trasmissione del radar va al piedino di ricezione dell'ESP32 e viceversa.

\begin{table}[!htb]
\centering
\caption{Piedinatura dell'\ldb{} (Tabella 1 del datasheet \cite{ld2410bDatasheet}) e collegamento all'ESP32.}
\label{tab:c3-2}
\small
\begin{tabularx}{\linewidth}{@{}clXl@{}}
\toprule
\textbf{Pin} & \textbf{Nome} & \textbf{Funzione} & \textbf{Collegamento ESP32} \\
\midrule
1 & OUT & presenza digitale (alto = persona) & non collegato \\
2 & UART\_Tx & il radar trasmette & GPIO25 (RX2) \\
3 & UART\_Rx & il radar riceve & GPIO26 (TX2) \\
4 & GND & massa & GND \\
5 & VCC & alimentazione & VIN (5~V) \\
\bottomrule
\end{tabularx}
\end{table}

Il piedino OUT, che riporta lo stato di presenza come semplice livello logico, non è stato collegato perché la stessa informazione è già contenuta nel frame seriale, come verificato nella Sezione~\ref{sec:c3-5-4}.

La velocità di 256000~baud richiede una delle UART hardware dell'ESP32. Si è usata la seconda, sui pin GPIO25 e GPIO26, perché la prima è occupata dal collegamento USB con il PC.

\subsection{Strumenti di configurazione}
\label{sec:c3-1-3}

Il produttore distribuisce due strumenti di configurazione:

\begin{itemize}
  \item il LD2410 Tool per PC, che dialoga con il modulo via UART attraverso un adattatore USB-seriale \cite{hilinkDriveLD2410B}
  \item l'applicazione mobile HLKRadarTool, che si collega via Bluetooth con la password di default HiLink \cite{hilinkV244}.
\end{itemize}

Entrambi permettono di leggere e modificare i parametri (gate massimi, soglie per gate, timeout di presenza, risoluzione), di attivare l’engineering mode e di osservare in tempo reale i grafici dell'energia per gate e la distanza rilevata. L'applicazione mobile offre in più l'auto-calibrazione delle soglie introdotta dal firmware V2.44, non utilizzata in questo lavoro. Inizialmente il LD2410 Tool è stato utilizzato per una prima verifica del modulo e per confrontare i dati acquisiti. In seguito si è fatto fede alla lettura via UART, perché lo strumento mostra per le sensibilità il valore 100 su tutti i gate, mentre il comando di lettura del protocollo restituisce la scala di fabbrica riportata nella Tabella~\ref{tab:c3-6}, coerente con il comportamento osservato. Lo strumento non mostra la versione del firmware e non offre funzioni di aggiornamento. La Figura~\ref{fig:c3-tool2410} mostra la schermata dello strumento con il modulo in prova.

\begin{figure}[!htb]
\centering
\includegraphics[width=0.85\linewidth]{Tool_LD2410B}
\caption{Schermata del LD2410 Tool collegato all'esemplare in prova, in engineering mode. A sinistra i parametri letti dal modulo, con le sensibilità mostrate a 100 su tutti i gate; al centro le energie per gate della componente in movimento e di quella stazionaria (in colore) con le soglie di fabbrica sovrapposte (in verde); in basso la distanza rilevata negli ultimi due minuti.}
\label{fig:c3-tool2410}
\end{figure}

\subsection{Protocollo seriale}
\label{sec:c3-1-4}

Il protocollo seriale prevede due tipi di frame per lo scambio di informazioni fra il modulo e l'host \cite{hilinkProtoV107}:

\textbf{Frame di comando} (inviati dall'host al modulo): la struttura è riportata nella Tabella~\ref{tab:c3-3}. Ogni comando va inviato dentro la modalità di configurazione, durante la quale il modulo sospende l'invio dei dati. Il modulo risponde con un frame della stessa struttura, in cui la parola di comando ha il bit più significativo posto a uno ed è seguita da uno stato (zero = successo) e dagli eventuali dati di risposta. Tutti i campi a più byte sono little-endian.

\begin{table}[!htb]
\centering
\caption{Struttura del frame di comando, dall'host al modulo \cite{hilinkProtoV107}.}
\label{tab:c3-3}
\small
\begin{tabularx}{\linewidth}{@{}llX@{}}
\toprule
\textbf{Campo} & \textbf{Lunghezza} & \textbf{Valore} \\
\midrule
Intestazione & 4~byte & FD FC FB FA \\
Lunghezza dati & 2~byte & little-endian \\
Parola di comando & 2~byte & little-endian \\
Valore del comando & variabile & dipende dal comando \\
Chiusura & 4~byte & 04 03 02 01 \\
\bottomrule
\end{tabularx}
\end{table}

\textbf{Frame di dati} (inviati dal modulo all'host a cadenza fissa): la struttura è riportata nella Tabella~\ref{tab:c3-4}. Il byte di tipo distingue il frame normale (0x02) da quello di engineering mode (0x01), e il byte di stato è una maschera di bit in cui un bersaglio in movimento e uno stazionario possono essere presenti insieme. I campi aggiunti dall'engineering mode sono quelli che fanno la differenza fra un sensore binario e uno strumento di misura (Sezione~\ref{sec:c3-1-5}).

\begin{table}[!htb]
\centering
\caption{Struttura del frame di dati, dal modulo all'host \cite{hilinkProtoV107}.}
\label{tab:c3-4}
\small
\begin{tabularx}{\linewidth}{@{}lcX@{}}
\toprule
\textbf{Campo} & \textbf{Lunghezza} & \textbf{Valore} \\
\midrule
Intestazione & 4~byte & F4 F3 F2 F1 \\
Lunghezza dati & 2~byte & little-endian \\
Tipo di dato & 1~byte & 0x01 = engineering mode, 0x02 = normale \\
Testa & 1~byte & 0xAA \\
Stato del bersaglio & 1~byte & maschera: 0x01 in movimento, 0x02 stazionario \\
Distanza in movimento & 2~byte & cm, little-endian \\
Energia in movimento & 1~byte & 0–100 \\
Distanza stazionaria & 2~byte & cm, little-endian \\
Energia stazionaria & 1~byte & 0–100 \\
Distanza di rilevamento & 2~byte & cm, little-endian \\
\multicolumn{3}{@{}l@{}}{\emph{(solo in engineering mode)}} \\
Gate massimo movimento & 1~byte &  \\
Gate massimo stazionario & 1~byte &  \\
Energie per gate, movimento & 9~byte & gate 0–8, valori 0–100 \\
Energie per gate, stazionario & 9~byte & gate 0–8, valori 0–100 \\
Sensore di luce & 1~byte & 0–255 \\
Stato del pin OUT & 1~byte &  \\
\multicolumn{3}{@{}l@{}}{\emph{(fine della parte engineering)}} \\
Coda & 1~byte & 0x55 \\
Controllo & 1~byte & 0x00 \\
Chiusura & 4~byte & F8 F7 F6 F5 \\
\bottomrule
\end{tabularx}
\end{table}

I comandi utilizzati in questo lavoro sono riassunti nella Tabella~\ref{tab:c3-5}:

\begin{table}[!htb]
\centering
\caption{Comandi del protocollo seriale utilizzati in questo lavoro \cite{hilinkProtoV107}.}
\label{tab:c3-5}
\small
\begin{tabularx}{\linewidth}{@{}lX@{}}
\toprule
\textbf{Comando} & \textbf{Funzione} \\
\midrule
0x00FF & entra in modalità configurazione \\
0x00FE & esce dalla modalità configurazione \\
0x0062 & attiva engineering mode \\
0x0063 & disattiva engineering mode \\
0x00A0 & legge la versione firmware \\
0x00A5 & legge il MAC Bluetooth \\
0x0061 & legge i parametri correnti (soglie, gate, timeout) \\
0x0060 & imposta i parametri \\
0x00AA & imposta la risoluzione di distanza \\
0x000B & auto-calibrazione del rumore di fondo (firmware V2.44) \\
\bottomrule
\end{tabularx}
\end{table}

\subsection{Engineering mode e dati per gate}
\label{sec:c3-1-5}

In modalità normale il modulo riporta, per ogni frame, lo stato di presenza, la distanza e l'energia del bersaglio in movimento e stazionario, cioè una decisione già presa accompagnata da pochi valori aggregati. In engineering mode aggiunge le energie di ciascuno dei nove gate, separatamente per la componente in movimento e per quella stazionaria (diciotto valori aggiuntivi) più la lettura del fotosensore e lo stato del pin OUT.

Le serie temporali di questi canali sono la materia prima di tutte le analisi quantitative dei capitoli successivi, e per questo il firmware realizzato tiene l’engineering mode sempre attiva.

\subsection{Identità dell'esemplare in prova}
\label{sec:c3-1-6}

Per la riproducibilità dei risultati, la Tabella~\ref{tab:c3-6} riporta identità e parametri dell'esemplare effettivamente usato, letti via UART con i comandi di lettura della Tabella~\ref{tab:c3-5}.

\begin{table}[!htb]
\centering
\caption{Identità e parametri dell'esemplare di \ldb{} in prova, letti via UART.}
\label{tab:c3-6}
\small
\begin{tabularx}{\linewidth}{@{}lX@{}}
\toprule
\textbf{Dato} & \textbf{Valore letto} \\
\midrule
Versione firmware & 2.44.25070917 \\
Versione protocollo & 1 \\
MAC Bluetooth & 5E:E4:93:22:B9:3F \\
Baud confermato & 256000, 8N1 \\
Risoluzione gate & 75~cm; gate 0–8; portata configurabile 600~cm (8 gate \(\times\) 0,75~m \cite[§5.2]{hilinkManualV104}) \\
Timeout di presenza (no-one duration) & 5~s \\
Soglie di movimento (gate 0–8) & 50, 50, 40, 30, 20, 15, 15, 15, 15 \\
Soglie stazionarie (gate 0–8) & 0, 0, 40, 40, 30, 30, 20, 20, 20 \\
\bottomrule
\end{tabularx}
\end{table}

Le soglie coincidono valore per valore con quelle di fabbrica riportate nella Tabella 7 del protocollo \cite{hilinkProtoV107}. Le soglie stazionarie dei gate 0 e 1 sono indicate dal documento come non impostabili, con le conseguenze discusse nella Sezione~\ref{sec:c3-5-2}. Il firmware 2.44 è l'ultima versione rilasciata dal produttore, la stessa che introduce l'auto-calibrazione citata nella Sezione~\ref{sec:c3-1-3} \cite{hilinkV244}.

\FloatBarrier
\section{Il radar \ldc{}}
\label{sec:c3-2}

L'\ldc{} è un modulo radar a 24~GHz per il rilevamento della presenza umana, anch'esso prodotto da Shenzhen Hi-Link Electronic \cite{hilinkManualLD2420}. È una scheda quadrata che integra il chip radar FMCW, le antenne e un microcontrollore, ma non dispone di Bluetooth, quindi la configurazione avviene soltanto via seriale. L'elaborazione è a bordo e verso l'esterno il modulo espone, su un connettore a cinque piedini, una linea seriale e un'uscita digitale di presenza. Riporta lo stato di occupazione e la distanza del bersaglio in movimento. Sulla persona immobile, invece, mantiene la presenza sfruttando i micro-movimenti del corpo (Sezione~\ref{subsec:corpo-immobile}), ma non è in grado di misurarne la distanza, come dichiarato dal produttore \cite[§8]{hilinkManualLD2420}. Il modulo è mostrato nella Figura~\ref{fig:c3-foto2420}.

\begin{figure}[!htb]
\centering
\begin{minipage}[t]{0.42\linewidth}\centering\includegraphics[width=\linewidth]{LD2420_fronte}\end{minipage}\hfill
\begin{minipage}[t]{0.42\linewidth}\centering\includegraphics[width=\linewidth]{LD2420_retro}\end{minipage}
\caption{Il modulo \ldc{}, fronte e retro. Sul fronte le antenne a patch e il connettore J2 a cinque piedini usato per alimentazione e seriale; sul retro il microcontrollore e il connettore J1 dell'interfaccia di programmazione, non utilizzato. Immagini del produttore \cite[Tab.~3-1 e 3-2]{hilinkManualLD2420}.}
\label{fig:c3-foto2420}
\end{figure}

\subsection{Specifiche tecniche}
\label{sec:c3-2-1}

Le specifiche dichiarate dal produttore sono riassunte nella Tabella~\ref{tab:c3-7} \cite{hilinkManualLD2420}.

\begin{table}[!htb]
\centering
\caption{Specifiche dell'\ldc{} secondo il manuale del produttore \cite{hilinkManualLD2420}.}
\label{tab:c3-7}
\small
\begin{tabularx}{\linewidth}{@{}lX@{}}
\toprule
\textbf{Parametro} & \textbf{Valore} \\
\midrule
Frequenza di lavoro & 24–24,25~GHz, banda di sweep 0,25~GHz \\
Tecnologia & FMCW (onda continua modulata in frequenza) \\
Potenza irradiata (EIRP) & 11~dBm \\
Portata dichiarata & a parete: 8~m su bersaglio in movimento, 6~m su micro-movimento; a soffitto: 5~m e 4~m \\
Distanza minima & 0,2~m, nessuna zona morta \\
Accuratezza di distanza & \(\pm\)0,35~m (valore teorico, dipende dalla corporatura del bersaglio) \\
Campo visivo & \(\pm\)60\(^\circ\) (§1.1); \(\pm\)45\(^\circ\) in orizzontale e in elevazione (§5.2) \\
Cadenza dati & 10~Hz \\
Alimentazione & 3,0–3,6~V, tipica 3,3~V \\
Corrente media & 50~mA \\
Interfaccia & UART 115200~baud, 8N1 \\
Bluetooth & assente \\
Dimensioni & 20 \(\times\) 20~mm \\
Temperatura operativa & \(-\)40 / +85~\(^\circ\)C \\
Gate di distanza & 16 (0–15) \\
Risoluzione gate & 0,70~m, fissa \\
\bottomrule
\end{tabularx}
\end{table}

\textbf{Due incongruenze nella documentazione ufficiale:}

\begin{itemize}
  \item Apertura angolare: il manuale la riporta con valori diversi in due punti, \(\pm\)60\(^\circ\) fra le caratteristiche principali e \(\pm\)45\(^\circ\) nella sezione sul montaggio a parete \cite[§1.1 e §5.2]{hilinkManualLD2420}.
  \item Ritardo di scomparsa del bersaglio: il manuale indica l’intervallo 0--65\,535 \cite[§4.2.1]{hilinkManualLD2420}, mentre il documento di protocollo indica 0x00–0x0F \cite[Tab. 2]{hilinkProtoLD2420}. Proprio gli esempi contenuti nel protocollo, però, usano i valori 30 e 26, che cadono fuori da quell'intervallo.
\end{itemize}

\subsection{Piedinatura e collegamento all'ESP32}
\label{sec:c3-2-2}

Il modulo espone cinque piedini, definiti nella Tabella 3-2 del manuale \cite{hilinkManualLD2420}. La Tabella~\ref{tab:c3-8} riporta la funzione di ciascuno e il collegamento all'ESP32 adottato in questo lavoro.

\begin{table}[!htb]
\centering
\caption{Piedinatura del connettore J2 dell'\ldc{} \cite[Tab. 3-2]{hilinkManualLD2420} e collegamento all'ESP32.}
\label{tab:c3-8}
\small
\begin{tabularx}{\linewidth}{@{}llXl@{}}
\toprule
\textbf{Pin} & \textbf{Nome} & \textbf{Funzione} & \textbf{Collegamento ESP32} \\
\midrule
J2-1 & 3V3 & alimentazione 3,0–3,6~V & 3V3 \\
J2-2 & GND & massa & GND \\
J2-3 & OT1 & UART\_TX, 0–3,3~V (il radar trasmette) & GPIO16 (RX) \\
J2-4 & RX & UART\_RX, 0–3,3~V (il radar riceve) & GPIO17 (TX) \\
J2-5 & OT2 & uscita di presenza: alto occupato, basso libero & non collegato \\
\bottomrule
\end{tabularx}
\end{table}

La documentazione del componente ESPHome \cite{esphomeLD2420} segnala che nei moduli con firmware precedente alla versione 1.5.3 la trasmissione seriale avviene sul piedino OT2 anziché su OT1 e a 256000~baud anziché 115200. Il firmware dell'esemplare in dotazione è la 1.6.1, quindi più recente, e infatti il modulo trasmette su OT1 a 115200~baud, mentre OT2 riporta lo stato di presenza come livello logico, come da manuale.

\subsection{Strumenti di configurazione}
\label{sec:c3-2-3}

Il produttore distribuisce l'\ldc{}\_TOOL per PC \cite{hilinkDriveLD2420}, che dialoga con il modulo via UART attraverso un adattatore USB-seriale. Permette di leggere e scrivere i parametri (gate minimo e massimo, ritardo di scomparsa e le trentadue soglie per gate), di osservare in tempo reale l'energia di ciascun gate e la distanza rilevata, e di salvare la configurazione in un file XML. Le soglie sono mostrate in decibel, cioè come \(10\cdot\log_{10}\) del valore intero che il modulo usa internamente. La Figura~\ref{fig:c3-tool2420} mostra la schermata dei parametri con la configurazione di fabbrica dell'esemplare.

\begin{figure}[!htb]
\centering
\includegraphics[width=0.85\linewidth]{Tool_LD2420}
\caption{Schermata dei parametri dell'\ldc{}\_TOOL collegato all'esemplare in prova, con la configurazione di fabbrica letta dal modulo: gate massimo 12, ritardo di scomparsa 30~s, le sedici soglie di attivazione e le sedici di mantenimento in decibel, versione firmware 1.6.1.}
\label{fig:c3-tool2420}
\end{figure}

La funzione più importante è la scansione del rumore di fondo (bottom noise scan, \cite[§4.2.2–4.2.3]{hilinkManualLD2420}). A stanza vuota il tool registra per ogni gate il livello di energia del fondo e ne ricava le due soglie, una di attivazione (trigger, consigliata a circa 5 volte il rumore) e una di mantenimento (maintain, a 2–5 volte), che poi scrive nella memoria non volatile del modulo.

\subsection{Protocollo seriale}
\label{sec:c3-2-4}

Il documento di protocollo \cite{hilinkProtoLD2420} descrive esclusivamente i comandi di configurazione, con frame delimitati dall'intestazione FD FC FB FA e dalla chiusura 04 03 02 01 e con lunghezza e parola di comando a 2~byte little-endian. I comandi disponibili sono riassunti nella Tabella~\ref{tab:c3-9}, e i parametri si leggono e si scrivono per identificatore.

\begin{table}[!htb]
\centering
\caption{Comandi e identificatori di parametro del protocollo dell'\ldc{} \cite{hilinkProtoLD2420}. Il comando segnato con l'asterisco non compare nel protocollo ufficiale ed è ricostruito dalla comunità ESPHome \cite{esphomeLD2420src}.}
\label{tab:c3-9}
\small
\begin{tabularx}{\linewidth}{@{}lX@{}}
\toprule
\textbf{Comando} & \textbf{Funzione} \\
\midrule
0x00FF & entra in modalità comandi \\
0x00FE & esce dalla modalità comandi \\
0x0000 & legge la versione firmware \\
0x0068 & riavvia il modulo \\
0x0008 & legge parametri (per identificatore) \\
0x0007 & scrive parametri (per identificatore) \\
0x0012* & scrive un parametro di sistema: valore 0x0004 per l'energy mode, 0x0064 per la modalità di testo di fabbrica \\
id 0x00 / 0x01 & gate minimo / gate massimo (0–15) \\
id 0x04 & ritardo di scomparsa del bersaglio, in secondi \\
id 0x10–0x1F & soglie di attivazione (trigger) dei gate 0–15 \\
id 0x20–0x2F & soglie di mantenimento (maintain) dei gate 0–15 \\
\bottomrule
\end{tabularx}
\end{table}

\textbf{Dati trasmessi dal modulo:} il documento di protocollo descrive soltanto i comandi che il modulo riceve e non il formato con cui trasmette i risultati del rilevamento. Il manuale indica di aprire un terminale seriale a 115200~baud per leggerli \cite[§4.1]{hilinkManualLD2420}. In modalità di fabbrica il modulo trasmette righe di testo che riportano la presenza rilevata (ON), l'assenza di presenza (OFF) e la distanza del bersaglio (Range NN).

L'unità di misura del campo Range non è indicata in nessuno dei due documenti del produttore. Le prove a distanza nota descritte nel Capitolo~\ref{chap:confronto} hanno mostrato che il valore è espresso in centimetri quando un bersaglio è presente, mentre in assenza di bersaglio il campo assume un valore di riposo compreso fra 0 e 6, privo di significato come distanza. Il valore si aggiorna di rado e va quindi letto soltanto quando la presenza è alta.

\textbf{Persistenza dei parametri.} I parametri scritti con il comando 0x0007 sono accettati e riletti correttamente, ma il modulo li perde al ciclo di alimentazione, e il protocollo non prevede un comando di salvataggio. La configurazione di riferimento è quindi quella scritta in memoria non volatile dal tool (Sezione~\ref{sec:c3-2-3}).

\subsection{Energy mode e dati per gate}
\label{sec:c3-2-5}

L'\ldc{} può trasmettere, al posto della sola decisione di presenza, il profilo di energia per gate da cui quella decisione è ricavata. Questa modalità, detta energy mode, non è descritta in alcun documento ufficiale, anche se il tool del produttore la utilizza per tracciare i grafici dell'energia per gate. Il suo funzionamento è stato ricostruito dalla comunità ESPHome tramite reverse engineering, e il risultato è implementato nel componente ld2420 \cite{esphomeLD2420src}, che costituisce la fonte su cui si basa il logger realizzato per questo lavoro. Si attiva con il comando 0x0012 della Tabella~\ref{tab:c3-9}, dopodiché il modulo emette a 10~Hz frame binari da 45~byte con la struttura della Tabella~\ref{tab:c3-10}.

\begin{table}[!htb]
\centering
\caption{Struttura del frame di dati in energy mode, ricostruita dal componente ESPHome \cite{esphomeLD2420src}.}
\label{tab:c3-10}
\small
\begin{tabularx}{\linewidth}{@{}llX@{}}
\toprule
\textbf{Campo} & \textbf{Lunghezza} & \textbf{Valore} \\
\midrule
Intestazione & 4~byte & F4 F3 F2 F1 \\
Lunghezza dati & 2~byte & little-endian \\
Presenza & 1~byte & 0 = libero, 1 = occupato \\
Distanza & 2~byte & cm, little-endian \\
Energie per gate & 16 \(\times\) 2~byte & gate 0–15, interi 0--65\,535, little-endian \\
Chiusura & 4~byte & F8 F7 F6 F5 \\
\bottomrule
\end{tabularx}
\end{table}

Le energie sono interi a 16~bit e, convertite in decibel con la stessa relazione delle soglie, coincidono con i valori mostrati dal tool, a conferma che la struttura ricostruita è corretta.

\subsection{Identità dell'esemplare in prova}
\label{sec:c3-2-6}

La Tabella~\ref{tab:c3-11} riporta identità e parametri dell'esemplare effettivamente usato, letti via UART con i comandi della Tabella~\ref{tab:c3-9} e confrontati con il file XML salvato dal tool.

\begin{table}[!htb]
\centering
\caption{Identità e parametri dell'esemplare di \ldc{} in prova, letti via UART.}
\label{tab:c3-11}
\small
\begin{tabularx}{\linewidth}{@{}lX@{}}
\toprule
\textbf{Dato} & \textbf{Valore letto} \\
\midrule
Versione firmware & 1.6.1 (ultima rilasciata dal produttore) \\
Baud confermato & 115200, 8N1 \\
Gate minimo / massimo impostati & 0 / 12 (minimo / massimo), su un intervallo ammesso 0–15 (portata configurata 840~cm) \\
Ritardo di scomparsa & 30~s \\
Soglie trigger, gate 0–3 & 60\,000, 30\,000, 3000, 2000 (47,8 / 44,8 / 34,8 / 33,0~dB) \\
Soglie maintain, gate 9–15 & 100 (20~dB) \\
\bottomrule
\end{tabularx}
\end{table}

Il modulo era in configurazione di fabbrica \cite{hilinkProtoLD2420}. Il produttore non distribuisce aggiornamenti del firmware, per cui la versione riportata è quella con cui il modulo è stato prodotto.

Il radar in dotazione è con ogni probabilità difettoso, perché nelle prove a distanza nota rileva una persona in movimento fino a circa 1~m e la mantiene fino a 2~m, contro gli 8~m dichiarati. Le energie per gate indicano un trasmettitore molto più debole del previsto, mentre il ricevitore funziona. Soglie, configurazione dei gate, firmware, parser, orientamento e cablaggio sono stati esclusi con prove dedicate. Le conseguenze sul protocollo della campagna sono dichiarate nella Sezione~\ref{sec:protocollo-ld2420}. Una relazione di progetto svolta a Camerino nello stesso anno accademico, con altri esemplari degli stessi moduli~\cite{dellaquila2026}, aiuta a separare i limiti del pezzo da quelli del modello. Sull'\ldc{} vi si ritrovano la presenza inaffidabile e la necessità di tarare le soglie per gate, che sono quindi proprie del modello, mentre l'accuratezza di distanza vi è giudicata buona, coerente con un limite di portata proprio del nostro esemplare.

\FloatBarrier
\section{Il sensore PIR in dotazione: \pir{}}
\label{sec:c3-3}

Il modulo in dotazione è un \pir{}, il PIR più diffuso per microcontrollori. Sulla scheda si trova l'elemento piroelettrico, coperto dalla cupola bianca della lente di Fresnel (Sezione~\ref{subsec:fresnel}). Il circuito integrato BISS0001 \cite{biss0001} amplifica il segnale e lo trasforma in un'uscita digitale con soglia e tempo di ritenuta. Un regolatore lineare HT7133 porta l'alimentazione a 3,3~V e fissa a questo livello anche l'uscita. Due trimmer regolano la sensibilità, e quindi la portata, e il tempo di ritenuta dell'uscita, mentre un ponticello a tre pin seleziona la modalità di trigger, singolo o ripetibile. Il modulo è mostrato nella Figura~\ref{fig:c3-fotopir}.

\begin{figure}[!htb]
\centering
\includegraphics[width=0.5\linewidth]{Foto_HCSR501}
\caption{Il modulo \pir{}, fronte e retro. Sul fronte la cupola della lente di Fresnel; sul retro il circuito integrato BISS0001, i due trimmer di sensibilità e di ritenuta, il regolatore di tensione, il ponticello di trigger (in giallo) e i tre piedini di uscita. Immagine del distributore \cite{alibabaHCSR501}.}
\label{fig:c3-fotopir}
\end{figure}

\subsection{Specifiche tecniche}
\label{sec:c3-3-1}

Le specifiche dichiarate dal datasheet \cite{hcsr501} sono riassunte nella Tabella~\ref{tab:c3-12}.

\begin{table}[!htb]
\centering
\caption{Specifiche del PIR \pir{} in dotazione (datasheet \cite{hcsr501}).}
\label{tab:c3-12}
\small
\begin{tabularx}{\linewidth}{@{}lX@{}}
\toprule
\textbf{Parametro} & \textbf{Valore} \\
\midrule
Base & BISS0001 + elemento LHI778 \\
Alimentazione & 4,5–20~V DC (collegato a VIN 5~V dell'ESP32) \\
Livello logico di uscita & alto 3,3~V / basso 0~V, grazie al regolatore HT7133 a bordo: compatibile con i GPIO dell'ESP32 senza adattamento \\
Corrente di riposo & \(<\)~50~\textmu{}A \\
Portata & 3–7~m, regolabile con il trimmer di sensibilità \\
Angolo di copertura & cono \(<\)~100\(^\circ\) \\
Tempo di ritenuta & regolabile, 5--200~s dichiarati (circa 3,4~s misurati al minimo del trimmer) \\
Block time & 2,5~s dopo ogni rilascio \\
Trigger & jumper: L = singolo, H = ripetibile \\
Temperatura operativa & \(-\)15 / +70~\(^\circ\)C \\
Dimensioni & 32 \(\times\) 24~mm \\
\bottomrule
\end{tabularx}
\end{table}

\subsection{Piedinatura e collegamento all'ESP32}
\label{sec:c3-3-2}

Il modulo espone tre piedini, riportati nella Tabella~\ref{tab:c3-13} insieme al collegamento all'ESP32 adottato in questo lavoro.

\begin{table}[!htb]
\centering
\caption{Piedinatura dell'\pir{} (visto dal lato dei trimmer) e collegamento all'ESP32.}
\label{tab:c3-13}
\small
\begin{tabularx}{\linewidth}{@{}clXl@{}}
\toprule
\textbf{Pin} & \textbf{Nome} & \textbf{Funzione} & \textbf{Collegamento ESP32} \\
\midrule
1 & GND & massa & GND \\
2 & OUT & uscita di presenza, alto 3,3~V / basso 0~V & GPIO34 \\
3 & VCC & alimentazione 4,5–20~V & VIN (5~V) \\
\bottomrule
\end{tabularx}
\end{table}

\subsection{Configurazione adottata}
\label{sec:c3-3-3}

La configurazione del modulo è riassunta nell'elenco seguente, e i suoi effetti sono discussi nel Capitolo~\ref{chap:confronto}.

\begin{enumerate}
  \item \textbf{Ponticello di trigger:} le prime acquisizioni sono state svolte in posizione L, poi il ponticello è stato spostato su H ripetendo gli scenari interessati.
  \item \textbf{Trimmer del tempo di ritenuta al minimo:} la ritenuta misurata è di circa 3,4~s, invariata per tutta la campagna.
  \item \textbf{Trimmer di sensibilità a metà corsa:} è stato portato al massimo soltanto nelle prove di portata massima, per verificare il valore dichiarato dal datasheet.
\end{enumerate}

\FloatBarrier
\section{La piattaforma di acquisizione}
\label{sec:c3-4}

I tre sensori sono letti da un'unica piattaforma, costruita perché ogni misura dei capitoli successivi sia riproducibile. La Tabella~\ref{tab:c3-14} elenca i firmware per l'ESP32 e la Tabella~\ref{tab:c3-14b} il software per il PC, diviso in acquisizione, analisi e materiali derivati.

Le sezioni che seguono descrivono il punto di partenza (Sezione~\ref{sec:c3-4-1}), i firmware (Sezione~\ref{sec:c3-4-2}), il formato dei dati che li accomuna (Sezione~\ref{sec:c3-4-3}), lo script di acquisizione (Sezione~\ref{sec:c3-4-4}) e la catena di analisi (Sezione~\ref{sec:c3-4-5}). Tutto il software è nel repository della tesi insieme ai file CSV della campagna, che sono l'unico contenuto non ricostruibile.

\begin{table}[!htb]
\centering
\caption{Il firmware realizzato per l'ESP32. I percorsi sono relativi alla radice del repository.}
\label{tab:c3-14}
\small
\begin{tabularx}{\linewidth}{@{}>{\raggedright\arraybackslash}p{4.3cm}>{\raggedright\arraybackslash}X@{}}
\toprule
\textbf{Sketch} & \textbf{Contenuto} \\
\midrule
firmware/ld2410b\_\allowbreak{}logger & logger CSV a 5~Hz dell'\ldb{} con engineering mode e lettura del PIR (fasi 1--7 della campagna) \\
firmware/ld2420\_\allowbreak{}logger\_\allowbreak{}bin & logger dell'\ldc{} in energy mode, 16 energie a 16~bit, stesso schema CSV \\
firmware/ld2410b\_\allowbreak{}web & web UI della Sezione~\ref{sec:webui}, stesso CSV del logger sulla seriale (Fase 8 della campagna) \\
firmware/test00\_\allowbreak{}ld2410\_\allowbreak{}info,\allowbreak{} test04\_\allowbreak{}set\_\allowbreak{}gate & lettura dei parametri dell'\ldb{} via UART e configurazione del gate massimo \\
firmware/test06\_\allowbreak{}ld2420\_\allowbreak{}set\_\allowbreak{}gate & lettura e scrittura dei parametri dell'\ldc{} via UART (32 soglie, gate, ritardo) \\
firmware/test01\_\allowbreak{}ld2410\_\allowbreak{}base,\allowbreak{} test03\_\allowbreak{}pir\_\allowbreak{}base & sketch diagnostici, baud rate del radar e durata degli impulsi del PIR \\
\bottomrule
\end{tabularx}
\end{table}

\begin{table}[!htb]
\centering
\caption{Il software realizzato per il PC: acquisizione, analisi e materiali derivati. I percorsi sono relativi alla radice del repository.}
\label{tab:c3-14b}
\small
\begin{tabularx}{\linewidth}{@{}>{\raggedright\arraybackslash}p{4.3cm}>{\raggedright\arraybackslash}X@{}}
\toprule
\textbf{Componente} & \textbf{Contenuto} \\
\midrule
HLK-LD2410x/acquire.py & acquisizione dalla seriale con metadati sperimentali, adattato dal repository di riferimento \\
data & i CSV della campagna e il registro delle sessioni \\
analisi/verifica\_\allowbreak{}engineering.py & controllo di qualità di un CSV prima dell'uso \\
analisi/analizza\_\allowbreak{}test.py & metriche di confronto e aggregazione per scenario \\
analisi/analizza\_\allowbreak{}respiro.py & analisi in frequenza e ricerca del canale migliore \\
analisi/vitalita\_\allowbreak{}proto.py,\allowbreak{} verifica\_\allowbreak{}vitalita\_\allowbreak{}bordo.py & prototipo dell'indice di vitalità e verifica del porting (Capitolo~\ref{chap:vitalita}) \\
analisi/rigenera\_\allowbreak{}tutto.py e script collegati & figure, foglio Excel e pagina di riepilogo, rigenerabili dai CSV con un comando \\
\bottomrule
\end{tabularx}
\end{table}

\subsection{Il repository di riferimento}
\label{sec:c3-4-1}

Il lavoro parte dal repository del relatore \cite{repoCallisto}, che contiene un firmware PlatformIO per ESP32 basato sulla libreria ncmreynolds/ld2410 e uno script Python di acquisizione. Il firmware legge il radar via UART, gestisce opzionalmente un display OLED e un PIR e stampa righe CSV sulla seriale. Lo script acquire.py le salva su file e aggiunge i metadati sperimentali passati da riga di comando. Lo studio del codice ha evidenziato tre limiti per gli scopi di questa tesi:

\begin{enumerate}
  \item \textbf{Campionamento a 1~Hz}, troppo lento per misurare le latenze e insufficiente per l'analisi del respiro, perché il limite di Nyquist a 0,5~Hz coincide con il bordo superiore della banda respiratoria (Sezione~\ref{sec:respiro}).
  \item \textbf{Assenza dell'engineering mode}, quindi viene registrata la sola energia del bersaglio e non i diciotto valori per gate della Sezione~\ref{sec:c3-1-5}.
  \item \textbf{Piedinatura diversa} (radar TX su GPIO17 e RX su GPIO16), quindi i due firmware non sono interscambiabili senza ricablare.
\end{enumerate}

\subsection{I firmware}
\label{sec:c3-4-2}

Per ciascun radar sono stati scritti due programmi per l'ESP32, uno sketch di configurazione che legge e scrive i parametri del modulo e un logger che acquisisce i dati. I due si caricano in momenti diversi. Prima si carica lo sketch, per portare il radar nella configurazione voluta, poi il logger, che resta a bordo per tutta la sessione. Il logger legge il radar e il PIR e stampa sulla porta seriale USB una riga di testo per campione, nel formato della Sezione~\ref{sec:c3-4-3}. L'ESP32 non salva nulla, perché il file viene scritto dallo script di acquisizione sul PC (Sezione~\ref{sec:c3-4-4}). Tutti i programmi stampano all'avvio una riga con il numero di build, così l'esito di una prova non può essere attribuito a un binario diverso da quello caricato. Il collegamento dei tre sensori all'ESP32 è riassunto nella Figura~\ref{fig:c3-collegamenti}.

\begin{figure}[!htb]
\centering
\includegraphics[width=\linewidth]{Collegamenti}
\caption{Collegamento dei tre sensori all'ESP32. L'\ldb{} è alimentato a 5~V da VIN e dialoga sulla seconda UART hardware (GPIO25 e 26); l'\ldc{} è alimentato a 3,3~V e usa una terza UART sui GPIO16 e 17, mentre la sua uscita di presenza OT2 non è collegata perché la presenza arriva dal frame seriale; il PIR è alimentato a 5~V e la sua uscita a 3,3~V va a GPIO34. I due fili della seriale sono incrociati. Le Tabelle~\ref{tab:c3-2}, \ref{tab:c3-8} e \ref{tab:c3-13} riportano il dettaglio dei piedini.}
\label{fig:c3-collegamenti}
\end{figure}

\paragraph{Gli sketch di configurazione}

I due sketch, uno per modulo, hanno tre accorgimenti importanti:

\begin{itemize}
  \item Dopo ogni scrittura i parametri sono riletti per conferma.
  \item Lo sketch avvisa quando la configurazione corrente è diversa da quella di fabbrica, per non condurre acquisizioni con parametri modificati inconsapevolmente.
  \item Nello sketch dell'\ldb{} un comando riscrive i valori di fabbrica dell'esemplare (Tabella~\ref{tab:c3-6}), così da poter tornare in ogni momento alla configurazione di riferimento.
\end{itemize}

Sul \ldc{} lo sketch ha permesso di verificare che le 32 soglie in flash coincidessero con quelle mostrate in decibel dallo strumento del produttore (Sezione~\ref{sec:c3-2-3}).

\paragraph{Il logger dell'\ldb{}}

Usa la libreria MyLD2410 \cite{myld2410}. Tre scelte condizionano la qualità dei dati.

\begin{itemize}
  \item \textbf{Campionamento a 5~Hz a periodo fisso}: il ciclo principale calcola l'istante del campione successivo anziché attendere un intervallo dopo l'elaborazione. Per questo la cadenza misurata è di 200~ms esatti su tutti gli intervalli (Sezione~\ref{sec:c3-5}).
  \item \textbf{Engineering mode attivato all'avvio e verificato}: il firmware entra in configurazione, attiva la modalità, esce e controlla che i frame ricevuti siano del tipo esteso, segnalando l'esito sulla seriale. Un fallimento silenzioso produrrebbe file senza colonne per gate, scoperti solo in analisi.
  \item \textbf{Lettura del PIR}: avviene nello stesso frame, immediatamente prima di comporre la riga CSV, così che i due sensori condividano la base dei tempi per costruzione e non per allineamento a posteriori.
\end{itemize}

\paragraph{Il logger dell'\ldc{}}

Legge i frame da 45~byte dell'energy mode (Sezione~\ref{sec:c3-2-5}) a 10~Hz e li riduce a 5~Hz per coerenza con il resto della campagna. Gate massimo, gate minimo e ritardo di scomparsa possono essere scritti via UART a ogni accensione, ma la scrittura non è persistente (Sezione~\ref{sec:c3-2-4}). Il valore riletto finisce in un commento in testa al CSV, così ogni file dichiara la configurazione con cui è stato acquisito. Il formato delle colonne è nella Sezione~\ref{sec:c3-4-3}.

\paragraph{Il firmware della web UI}

Legge lo stesso radar e lo stesso PIR del logger dell'\ldb{} e scrive sulla seriale le sue stesse colonne, così lo script di acquisizione può girare in parallelo alla pagina. In più pubblica i dati in rete e calcola a bordo l'indice di vitalità (Sezione~\ref{sec:webui}).

\subsection{Il formato dei dati}
\label{sec:c3-4-3}

Tutti i file della tesi seguono lo stesso schema a strati (Tabella~\ref{tab:c3-15}). Le prime nove colonne coincidono con quelle del repository di riferimento e sono prodotte da tutti i firmware. Su queste lavora lo script di confronto del Capitolo~\ref{chap:confronto}, quindi le metriche di presenza, falsi positivi, latenze e distanza si calcolano allo stesso modo per i due radar. Seguono le colonne proprie di ciascun radar, con nomi distinti (menergy, senergy contro energy2420) perché le scale sono diverse, 0--100 contro 0--65\,535, e nessuno script deve confonderle. Le grandezze che un radar non possiede valgono zero. Il solo CSV esportato dalla web UI aggiunge in coda due colonne con il valore e la classe dell'indice di vitalità calcolato a bordo (Capitolo~\ref{chap:vitalita}), messe in fondo perché la catena di analisi non le richiede e le ignora.

\begin{table}[!htb]
\centering
\caption{Schema a strati dei file CSV della tesi. Le prime nove colonne sono comuni a tutti i firmware; le colonne proprie di ciascun radar hanno nomi distinti perché le scale sono diverse.}
\label{tab:c3-15}
\small
\begin{tabularx}{\linewidth}{@{}>{\raggedright\arraybackslash}p{2.1cm}>{\raggedright\arraybackslash}X>{\raggedright\arraybackslash}p{2.4cm}c@{}}
\toprule
\textbf{Strato} & \textbf{Colonne} & \textbf{Prodotte da} & \textbf{N.} \\
\midrule
Base comuni & timestamp\_\allowbreak{}ms, radar\_\allowbreak{}presence, moving\_\allowbreak{}target, stationary\_\allowbreak{}target, moving\_\allowbreak{}distance\_\allowbreak{}cm, stationary\_\allowbreak{}distance\_\allowbreak{}cm, moving\_\allowbreak{}energy, stationary\_\allowbreak{}energy, pir\_\allowbreak{}presence & tutti i firmware & 9 \\
LD2410B & menergy\_\allowbreak{}gate0..8, senergy\_\allowbreak{}gate0..8, light\_\allowbreak{}level, out\_\allowbreak{}level & logger LD2410B, web UI & 20 (29 in totale) \\
LD2420 & energy2420\_\allowbreak{}gate0..15, frames\_\allowbreak{}ok, dist\_\allowbreak{}raw\_\allowbreak{}cm & logger LD2420 & 18 (27 in totale) \\
Metadati & pc\_\allowbreak{}time\_\allowbreak{}s, group\_\allowbreak{}id, trial\_\allowbreak{}id, scenario, ground\_\allowbreak{}truth\_\allowbreak{}presence, ground\_\allowbreak{}truth\_\allowbreak{}state & acquire.py & 6 \\
Indice a bordo & vitality\_\allowbreak{}onboard, vitality\_\allowbreak{}class\_\allowbreak{}onboard & export della web UI (Sezione~\ref{sec:webui}) & 2 \\
\bottomrule
\end{tabularx}
\end{table}

Nei test dell'\ldc{} il PIR non è cablato. Un piedino non collegato leggerebbe zero per tutto il file, e gli script ne ricaverebbero un 100~\% di falsi negativi con una persona davanti. Per questo il firmware scrive il valore sentinella \(-\)1, che lo script di analisi riconosce. In quel caso lo script omette ogni metrica del PIR e lo segnala in una colonna di stato del foglio Excel.

L'intestazione è stampata dal firmware una sola volta all'avvio. Le righe che iniziano con il cancelletto sono commenti, che lo script di acquisizione scarta, e contengono il marcatore di build di ogni sketch e la configurazione dell'\ldc{} (Sezione~\ref{sec:c3-4-2}). Lo script di acquisizione ricostruisce un'intestazione mancante dal numero di colonne (Sezione~\ref{sec:c3-4-4}), ma a 27 colonne il logger dell'\ldc{} e l'engineering mode senza fotosensore non si distinguono, quindi l'intestazione stampata dal firmware è indispensabile.

\subsection{Lo script di acquisizione}
\label{sec:c3-4-4}

Lo script acquire.py legge la seriale e aggiunge a ogni riga le sei colonne di metadati sperimentali della Tabella~\ref{tab:c3-15}, dichiarati dall'operatore da riga di comando. Il file prende il nome dello scenario e del trial, e lo script richiede la sola libreria pyserial.

Lo script ricostruisce l'intestazione dal numero di colonne della prima riga di dati (9 per la modalità base, 29 per l'engineering mode completo), invece di attendere quella stampata dall'ESP32 all'avvio. Altrimenti, quando il reset all'apertura della porta seriale non scatta, il CSV resterebbe vuoto senza alcun avviso.

Per facilitare i test sono stati introdotti:

\begin{itemize}
  \item Un segnale acustico all'istante dell'evento per le misure di latenza (Sezione~\ref{sec:latenza})
  \item Annunci vocali a istanti prefissati per le prove di portata.
\end{itemize}

Entrambi contano dalla prima riga di dati ricevuta, evitando l'errore di 2--3~s che un cronometro esterno introdurrebbe per il reset dell'ESP32.

\subsection{La catena di analisi}
\label{sec:c3-4-5}

L'analisi è svolta da cinque script Python basati sulla sola libreria standard, con l'eccezione di NumPy per la parte in frequenza, più un gruppo di script che produce i materiali derivati e richiede matplotlib, openpyxl e pillow.

\paragraph{verifica\_engineering.py}

È il controllo di qualità di un CSV prima dell'uso. Verifica la cadenza reale e il jitter, presenza delle colonne per gate, percentuale di saturazione a 100 per canale, coerenza fra gate di picco e distanza riportata, rumore di fondo per gate a stanza vuota e transizioni di presenza. Va eseguito a ogni sessione ed è ciò che ha permesso di individuare i comportamenti descritti nella Sezione~\ref{sec:c3-5}.

\paragraph{analizza\_test.py}

Produce i valori su cui si basano i confronti del Capitolo~\ref{chap:confronto}. Su ogni trial calcola falsi positivi in eventi all'ora, falsi negativi in percentuale, latenza di rilevamento e di rilascio con la convenzione dell'evento a istante noto, e statistiche di distanza ed energia. Lo script riunisce le ripetizioni di ogni scenario e per ciascuna metrica riporta la media e la deviazione standard fra i trial. Prima del calcolo scarta i secondi iniziali in cui il soggetto raggiunge la posizione, con una durata che dipende dallo scenario (Sezione~\ref{sec:metriche-transitorio}). Lo script legge indifferentemente i file dei due radar e della web UI.

\paragraph{analizza\_respiro.py}

Esegue l'analisi in frequenza (Sezione~\ref{sec:respiro}) e richiede NumPy. Applica la FFT alla serie di energia, cerca il picco in banda 0,1--0,5~Hz, stima gli atti al minuto ed esporta lo spettro. La modalità di scansione prova tutti i canali di energia, scarta quelli saturi e quelli piatti, li ordina per rapporto segnale-rumore e riporta la mediana delle stime concordi. È nata dal primo pilota del respiro, in cui il canale predefinito, l'energia stazionaria aggregata, era saturo per tutto il file mentre il respiro era leggibile su un gate di movimento. Da allora la ricerca automatica del canale migliore, con criterio dichiarato, è parte del metodo.

\paragraph{vitalita\_proto.py}

È il prototipo dell'indice di vitalità (Capitolo~\ref{chap:vitalita}). Legge i CSV dell'\ldb{} e calcola l'indice campione per campione, sottraendo il fondo per gate, prendendo l'energia del gate in cui il radar riporta il bersaglio e filtrandola con due medie mobili esponenziali. Su un insieme di trial produce la distribuzione dell'indice per scenario e la matrice di confusione delle tre classi, con cui le soglie sono state tarate e poi validate su trial separati. Lo script verifica\_vitalita\_bordo.py ha verificato che il firmware calcoli lo stesso indice del prototipo, confrontando campione per campione le due colonne in coda al CSV della web UI con il ricalcolo offline.

\paragraph{Materiali derivati}

Un ultimo gruppo di script produce i materiali derivati, cioè rifà dai CSV, con un solo comando, le figure di questa tesi, il foglio Excel con una riga per trial valido e la pagina di riepilogo. Le figure importano le funzioni di analizza\_test.py anziché ricalcolare, così ogni numero nei grafici coincide per costruzione con quello nel testo. Il foglio Excel ha una lista separata con i file esclusi e la ragione dell'esclusione, perché nessun file di prova o di controllo entri nelle statistiche per una svista.

\FloatBarrier
\section{Verifica dei dati e comportamenti del modulo da conoscere}
\label{sec:c3-5}

Prima di avviare la campagna, i dati prodotti dal logger sono stati controllati su una sessione di prova di 134~s (673 campioni), con una persona che entra nella stanza, si muove, si ferma ed esce. Il test serve da un lato a verificare che la catena di acquisizione funzioni come previsto, dall'altro a imparare a leggere i dati del modulo, individuando i comportamenti da tenere presenti nell'interpretare i risultati del Capitolo~\ref{chap:confronto}.

La verifica della catena ha dato esito positivo su quattro punti:

\begin{itemize}
  \item Cadenza di 5,00~Hz esatti, con intervallo di 200~ms su tutti i 672 intervalli e jitter nullo, quindi la serie temporale è uniforme, requisito dell'analisi in frequenza.
  \item Engineering mode attivo, con tutte e diciotto le colonne per gate popolate.
  \item Mappa gate-distanza validata al 97,1~\%. È la percentuale di campioni in cui il gate di massima energia coincide, entro \(\pm\)1, con quello atteso dalla distanza riportata (distanza divisa per 0,75~m), su tutto il range 0--6~m. Su una seconda sessione indipendente la conferma è salita al 98,3~\% su 663 campioni.
  \item Zero falsi positivi nei 29~s di stanza vuota della sessione.
\end{itemize}

La Figura~\ref{fig:c3-1} mostra l'intera sessione. I due pannelli superiori riportano le energie dei nove gate, per il canale in movimento e per quello stazionario, mentre il pannello inferiore mostra la presenza dichiarata dai due sensori nello stesso istante. La figura rende evidente la differenza di informazione fra i due sensori e mostra già i tre comportamenti del radar descritti nei sottoparagrafi seguenti. Il primo è la saturazione dell'energia, visibile nelle fasce chiare a 100 del canale stazionario. Il secondo sono i gate stazionari 0 e 1 sempre nulli, cioè la fascia nera in basso nel pannello centrale. Il terzo è la coda di presenza dopo l'uscita, cioè la barra blu che si prolunga oltre l'ultimo movimento. Non sono difetti dell'esemplare ma proprietà del modulo, e condizionano il modo in cui i dati vanno letti in tutto il resto della tesi.

\begin{figure}[!htb]
\centering
\includegraphics[width=0.92\linewidth]{Energia_gate_2410B}
\caption{Engineering mode dell'\ldb{} su una sessione con persona che entra, si ferma e si muove. In alto i nove gate del canale in movimento, con sovrapposto il gate atteso dalla distanza riportata. Al centro i nove gate stazionari, con i gate 0 e 1 sempre nulli. In basso la presenza dichiarata dal radar e dal PIR, in modalità L, nello stesso istante.}
\label{fig:c3-1}
\end{figure}

\subsection{Saturazione dei canali di energia}
\label{sec:c3-5-1}

L'energia riportata dal modulo è un intero fra 0 e 100, che non è una potenza e non si converte in decibel. Per questo le attenuazioni misurate con questo valore non sono confrontabili numericamente con quelle dell'\ldc{}, che sono conteggi interni a 16~bit espressi in decibel dallo strumento del produttore, e neppure con i coefficienti di trasmissione dei materiali. Fra i due radar si confronta la graduatoria dei materiali, non il valore (Sezione~\ref{sec:ostacoli}). Con un bersaglio fermo a breve distanza il canale stazionario tocca il valore massimo nel 92~\% dei campioni e quello in movimento nel 41~\% (Figura~\ref{fig:c3-2}). Un canale che resta al fondo scala non registra più le variazioni del segnale, e ogni analisi costruita su quei dati perde di significato. Il problema non si corregge con la configurazione, perché le soglie agiscono soltanto sulla decisione di presenza e non sul valore riportato \cite[§1.2.2]{hilinkProtoV107}, e lo stesso vale per l'auto-calibrazione del firmware V2.44, che tara le soglie e non la scala. Il rimedio è quindi geometrico, e consiste nel tenere il soggetto ad almeno 1,5–2~m, nel disallineare il sensore oppure nell'attenuare il segnale. Gli esperimenti sul respiro e sull'indice di vitalità (Capitolo~\ref{chap:vitalita}) sono stati progettati con questo vincolo.

\begin{figure}[!htb]
\centering
\includegraphics[width=\linewidth]{Saturazione_segnale}
\caption{La saturazione dell'energia. L'indice è un intero 0--100 e sul bersaglio vicino il canale stazionario è a fondo scala nella quasi totalità dei campioni. }
\label{fig:c3-2}
\end{figure}

\subsection{Canali stazionari disattivati sotto 1,5~m}
\label{sec:c3-5-2}

Le energie stazionarie dei gate 0 e 1 sono risultate nulle in tutti i 673 campioni, ma non si tratta di un difetto dell'esemplare. L'esempio di frame riportato nel protocollo \cite[§2.3.2]{hilinkProtoV107} mostra gli stessi valori nulli, e il protocollo indica le soglie di quei due gate come non impostabili (la lettura restituisce 0, Tabella~\ref{tab:c3-6}). Il canale stazionario è quindi disattivato per costruzione entro 1,5~m. Il modulo continua comunque a rilevare un bersaglio fermo a quella distanza. Nei dati compaiono infatti distanze stazionarie fra 70 e 89~cm con energia aggregata pari a 100, e quello che manca è soltanto il dettaglio per gate. Il vincolo che ne deriva riguarda le analisi costruite sulle serie per gate, come il respiro e l'indice di vitalità, che per un soggetto entro 1,5~m devono usare i canali in movimento. È il caso dello scenario DIPME, dove la persona sotto il banco si trova proprio a quella distanza (Sezione~\ref{sec:sotto-banco}).

\subsection{Coda di presenza}
\label{sec:c3-5-3}

Dopo l'uscita del soggetto dalla stanza il radar ha mantenuto la presenza per circa 10~s, riportando in quel tratto un bersaglio fermo a circa 5,2~m con energia in calo. Il comportamento è previsto dal protocollo, che introduce un ritardo di scomparsa prima di dichiarare la stanza vuota \cite[§1.2.2]{hilinkProtoV107}, impostato a 5~s di fabbrica sull'esemplare. Sul piano del metodo questa coda va trattata in due modi:

\begin{itemize}
  \item All'inizio dei trial va scartata insieme al transitorio
  \item Nel confronto fra sensori va misurata come latenza di rilascio (Sezione~\ref{sec:latenza})
\end{itemize}

Non va invece contata come falso positivo, altrimenti al radar verrebbero attribuiti decine di errori inesistenti. Un bersaglio stazionario fittizio compare anche durante il movimento, perché con una sola persona in moto continuo il canale stazionario riporta un bersaglio saturo a una distanza sbagliata. Di conseguenza le sue grandezze non vanno lette finché il bersaglio si muove.

\subsection{Il pin OUT e il sensore di luce}
\label{sec:c3-5-4}

Due verifiche minori completano il quadro. Lo stato del pin OUT, letto dal frame UART, coincide con la presenza dichiarata dal radar in tutti i 1699 campioni di una sessione. Il piedino può comunque essere usato al posto della seriale se necessario. Il fotosensore funziona, con valori fra 21 e 29 di giorno e fra 0 e 1 di notte su una sessione di 6,5~h.

\FloatBarrier
\section{Soglie di fabbrica e necessità di taratura}
\label{sec:c3-6}

Sull'\ldb{} le soglie (Tabella~\ref{tab:c3-6}) sono state lasciate ai valori di fabbrica per tutta la campagna. Confrontate con il rumore di fondo misurato a stanza vuota, restano sopra il massimo del rumore su ogni gate, e in tutta la campagna il radar non ha prodotto alcun falso positivo (Sezione~\ref{sec:falsi-positivi}). Non c'era quindi motivo di ricalibrare, e tenere la configurazione di fabbrica ha due vantaggi:

\begin{itemize}
  \item I trial sono confrontabili fra loro, perché nessun parametro è cambiato a metà campagna
  \item Le prove sono riproducibili da chiunque disponga dello stesso modulo, senza dover replicare una taratura.
\end{itemize}

L'\ldc{} si è comportato in modo opposto. Con la configurazione di fabbrica (Tabella~\ref{tab:c3-11}) il gate massimo a 12 porta il campo a 840~cm, oltre le pareti della stanza di prova, e le soglie di mantenimento dei gate lontani sono così basse che il riflesso del muro tiene la presenza sempre attiva. A stanza vuota il modulo non ha mai dichiarato l'assenza. È stato quindi necessario ridurre il gate massimo a 6 e ricavare le soglie dal rumore di fondo della stanza. La procedura automatica del tool (Sezione~\ref{sec:c3-2-3}) produceva soglie sotto le code del rumore reale, e il modulo continuava a non rilasciare. Le soglie usate nella campagna sono quindi state calcolate a mano con la regola del manuale a partire dal fondo misurato a stanza vuota (Sezione~\ref{sec:protocollo-ld2420}).

\chapter{Metodologia di confronto}
\label{chap:metodologia}

Questo capitolo descrive come i tre sensori sono stati confrontati.

La Sezione~\ref{sec:metodo-approccio} distingue le due attività su cui poggia il confronto,
quella sperimentale e quella numerica. La Sezione~\ref{sec:protocollo} fissa il protocollo
con cui sono stati acquisiti i 444 trial del Capitolo~\ref{chap:confronto}, mentre la
Sezione~\ref{sec:metriche} definisce le metriche e le convenzioni con cui sono stati letti.
La Sezione~\ref{sec:protocollo-ld2420} dichiara le deroghe al protocollo adottate per il
secondo radar, e la Sezione~\ref{sec:webui} presenta l'interfaccia web realizzata a supporto
della campagna.

\section{Due strade complementari}
\label{sec:metodo-approccio}

Il confronto poggia su due attività distinte che si alimentano a vicenda.
\begin{itemize}
  \item \textbf{Sperimentale:} consiste nel disporre i sensori e il soggetto in geometrie
    definite e nel cambiare, da una prova all'altra, una sola variabile per volta, che può
    essere la distanza, la quantità o il tipo di movimento, l'angolo oppure il materiale
    interposto. Ogni configurazione viene ripetuta e ogni ripetizione produce un file.
  \item \textbf{Numerica:} i file formano un dataset omogeneo, raccolto nel formato unico
    della Sezione~\ref{sec:c3-4-3} e analizzato a posteriori dalla catena descritta nella
    Sezione~\ref{sec:c3-4-5}. Le metriche, quindi, non vengono lette durante la prova. Sono
    calcolate in seguito, dagli stessi script e con le stesse convenzioni per tutti i
    sensori. In questo modo un risultato si può ricalcolare e una figura si può rifare
    senza ripetere l'acquisizione, e ogni numero del Capitolo~\ref{chap:confronto} risale a
    un file e a una riga di comando precisi.
\end{itemize}

Durante la prova, quindi, l'operatore non ha bisogno di leggere i numeri, e l'interfaccia web serve a controllare in tempo reale che i sensori vedano ciò che devono vedere (Sezione~\ref{sec:webui-ruolo}).

\section{Protocollo sperimentale}
\label{sec:protocollo}
\label{sec:protocollo-configurazione}

Il protocollo è stato definito prima delle acquisizioni. Il PIR lavora nella configurazione a lui più favorevole, con il ponticello di ritrigger in posizione H, in tutte le serie in cui il ponticello conta (Sezione~\ref{sec:interpretazione-pir}), mentre
l'\ldb\ mantiene sempre le soglie di fabbrica (Sezione~\ref{sec:c3-6}). L'\ldc\ è stato
provato in un secondo tempo nelle stesse geometrie, con le deroghe dichiarate nella
Sezione~\ref{sec:protocollo-ld2420}. Ogni scenario è ripetuto più volte, e i valori
riportati nella forma $\mu \pm \sigma$ sono la media e la deviazione standard su quelle
ripetizioni.

\subsection{Setup}
\label{sec:protocollo-setup}

L'\ldb\ e l'\pir\ sono letti dallo stesso firmware, sulla stessa base dei tempi
(Sezione~\ref{sec:c3-4-2}). I due sensori sono fissati uno accanto all'altro su un
supporto di cartone appoggiato sul bordo di un tavolo, a circa 85~cm di altezza
(Figura~\ref{fig:setup-foto}), e puntano verso un'area di prova che ricade nel campo di
entrambi. Nella stanza non ci sono oggetti in movimento (nessun ventilatore, tende ferme,
porta chiusa). Le quattro geometrie di prova sono schematizzate in
Figura~\ref{fig:geometrie}.

\begin{figure}[!htbp]
\centering
\includegraphics[width=\linewidth]{Setup}
\caption{Il banco di prova. (a) Il PIR \pir\ e l'\ldb\ fissati sul supporto di cartone,
con l'ESP32 sulla breadboard accanto. (b) Lo stesso supporto con l'\ldc\ della campagna
ridotta.}
\label{fig:setup-foto}
\end{figure}

\begin{figure}[!htbp]
\centering
\includegraphics[width=\linewidth]{Geometrie}
\caption{Le geometrie della campagna, in scala. (a) La stanza, con i sensori sul tavolo e
le tacche a 1--5~m sull'asse, usata per distanza, dose-risposta, latenze, ostacoli e due
persone. (b) Lo scenario del progetto, con i sensori fissati sotto il piano e la persona
rannicchiata a 50--60~cm. (c) L'arco a 1~m per la copertura angolare. (d) Il corridoio
della prova di portata, con i due vani porta a 4 e 5~m attraverso cui si vedono i punti
più lontani.}
\label{fig:geometrie}
\end{figure}

\subsection{Scenari}
\label{sec:protocollo-scenari}

La Tabella~\ref{tab:scenari} raggruppa gli scenari per domanda. Di norma uno scenario è ripetuto
cinque volte, dieci per la latenza d'ingresso e tre quando serve una seconda persona o le
condizioni da coprire sono molte. Ogni scenario ha un nome, che è anche il nome dei suoi
file. Il respiro a ritmo imposto usa lo stesso setup e lo stesso formato dei file ed è trattato
nel Capitolo~\ref{chap:vitalita}.

\begin{table}[!htbp]
\centering
\caption{Gli scenari della campagna, con i trial per condizione, lo scarto iniziale
applicato da tutti gli script (Sezione~\ref{sec:metriche-transitorio}) e la sezione dei
risultati. La campagna ridotta sull'\ldc\ ripete entro 2~m gli scenari con l'asterisco
(Sezione~\ref{sec:protocollo-ld2420}), con uno scarto di 90~s nei trial da fermo.}
\label{tab:scenari}
\small
\begin{tabularx}{\linewidth}{@{}>{\raggedright\arraybackslash}p{0.27\linewidth}>{\raggedright\arraybackslash}X>{\raggedright\arraybackslash}p{0.13\linewidth}c>{\raggedright\arraybackslash}p{0.13\linewidth}@{}}
\toprule
\textbf{Scenario} & \textbf{Domanda} & \textbf{Trial} & \textbf{Scarto} & \textbf{Risultati} \\
\midrule
Stanza vuota, notte e giorno* & falsi positivi & 7,03~h & 60~s & Sezione~\ref{sec:falsi-positivi} \\
Persona ferma seduta a 2,3~m & falsi negativi & 5 & 30~s & Sezione~\ref{sec:persona-ferma} \\
Dose-risposta a 1~m (immobile, micro-movimenti, cammino sul posto)* & quantità di movimento & 5 per condizione & 20~s & Sezione~\ref{sec:dose-risposta} \\
Cammino sul posto a 1--5~m* & accuratezza della distanza, PIR sul posto & 5 per distanza & 20~s & Sezioni~\ref{sec:distanza} e~\ref{sec:dose-risposta} \\
Attraversamento a 1--5~m & tipo di movimento & 3 per distanza & 20~s & Sezione~\ref{sec:attraversamento} \\
Sotto il banco a 60~cm (immobile, micro-movimenti)* & scenario del progetto & 5 per condizione & 40~s & Sezione~\ref{sec:sotto-banco} \\
Ingresso e uscita dalla stanza* & latenze & 10 e 5 & evento a 30~s & Sezione~\ref{sec:latenza} \\
Arco a 1~m, sei azimut* & copertura angolare & 1--3 per azimut & 20~s & Sezione~\ref{sec:angolare} \\
Selettività con gate massimo 2* & vicino di banco & 3 per scenario & 120~s & Sezione~\ref{sec:angolare} \\
Due persone, tre geometrie* & separazione dei bersagli & 3 per geometria & 20~s & Sezione~\ref{sec:due-persone} \\
Pannelli interposti, a 3~m e a 1~m* & attenuazione & 3 per materiale & 20~s & Sezione~\ref{sec:ostacoli} \\
Corridoio, 5--8~m* & portata massima & 3--5 per distanza & 20~s & Sezione~\ref{sec:portata-massima} \\
\bottomrule
\end{tabularx}
\end{table}

\subsection{Ground truth}
\label{sec:protocollo-ripetizioni}
\label{sec:protocollo-groundtruth}

La presenza e lo stato reali del soggetto sono dichiarati dall'operatore nei metadati di
ogni acquisizione (Sezione~\ref{sec:c3-4-3}) e valgono per l'intero file. Per le latenze
questo non basta. Il soggetto entra o esce quindi a un istante annunciato dallo script di
acquisizione (Sezione~\ref{sec:c3-4-4}), e la latenza si misura a partire da
quell'istante.

\subsection{Registro}
\label{sec:registro}

Ogni sessione è annotata in un registro con data, ora, temperatura della stanza,
soggetto, alimentazione e note. Per ogni scenario, oltre alla distanza, sono annotati
anche il tipo di movimento e la postura, perché il PIR risponde a questi e non alla
distanza (Sezione~\ref{subsec:fresnel}). Le sessioni si sono svolte fra 25 e 29\,\textdegree C. Il registro
completo è nel repository del lavoro~\cite{repoTesi}.

\section{Metriche e convenzioni}
\label{sec:metriche}

\subsection{Metriche}
\label{sec:metriche-elenco}

Le metriche sono calcolate su ogni trial dallo script di analisi della
Sezione~\ref{sec:c3-4-5} e poi aggregate per scenario.
\begin{itemize}
  \item \textbf{Falsi negativi (\%):} campioni con presenza reale in cui il sensore riporta
    assenza. È la metrica dello scenario con la persona ferma, e il suo complemento, cioè
    la frazione di campioni con presenza riportata, è il tasso di rilevamento.
  \item \textbf{Falsi positivi (eventi/h):} transizioni da assenza a presenza per ogni ora
    di stanza vuota. Si contano gli eventi e non i campioni, perché un falso positivo di
    dieci secondi non è dieci volte più grave di uno di un secondo.
  \item \textbf{Latenza di rilevamento e di rilascio:} tempo che passa fra l'evento e la
    prima transizione stabile del sensore, all'ingresso e all'uscita. Per confrontare i
    sensori si usa la differenza appaiata sullo stesso trial, in cui il tragitto e il
    tempo di reazione si cancellano.
  \item \textbf{Accuratezza della distanza:} scostamento fra la distanza riportata dal
    radar e quella reale, tenendo separata la dispersione fra un trial e l'altro da quella
    entro il singolo trial. La linearità della risposta si valuta con una regressione lineare fra distanza riportata e reale. Il coefficiente di determinazione $R^2$ indica quanta parte della variazione della distanza riportata è spiegata dalla retta, e vale 1 quando i punti vi stanno esattamente sopra.
  \item \textbf{Energia:} valore normalizzato 0--100 che l'\ldb\ riporta per il bersaglio e
    per ogni gate (Sezione~\ref{sec:c3-5-1}). Fa da grandezza graduata del movimento e da
    testimone dello stimolo, perché due prove con la stessa energia hanno riprodotto lo
    stesso movimento.
\end{itemize}

\subsection{Zero eventi}
\label{sec:metriche-zero}

Quando in una finestra di osservazione non si verifica alcun evento, non si scrive zero
eventi all'ora ma si riporta il limite superiore dato dalla regola del tre, $3/T$ con $T$
in ore, che approssima il limite al 95\,\% di confidenza per un processo di Poisson senza
eventi~\cite{hanley1983rule}. Questo numero dipende dal tempo di osservazione, non dal sensore.

\subsection{Transitorio iniziale}
\label{sec:metriche-transitorio}

I primi secondi di ogni trial vengono scartati, perché contengono il posizionamento del
soggetto e la coda di presenza del radar (Sezione~\ref{sec:c3-5-3}). La durata dello
scarto dipende dallo scenario, è riportata nella Tabella~\ref{tab:scenari} ed è la stessa in
tutti gli script e in tutte le figure. Negli scenari ordinari vale 20~s. È più lunga sotto
il banco, dove il soggetto deve anche posizionarsi, a stanza vuota, dove l'operatore deve
prima uscire, e negli scenari di selettività, dove la coda del radar dopo il
posizionamento arriva a superare i 100~s (Figura~\ref{fig:transitorio} e
Sezione~\ref{sec:angolare}). Nei trial da fermo dell'\ldc, che rilascia in circa 55~s, lo
scarto sale a 90~s.

\begin{figure}[!htbp]
\centering
\includegraphics[width=\linewidth]{Transitorio}
\caption{Il transitorio iniziale nei tre trial con la persona ferma a 1~m a
90\textdegree\ (gate massimo 2). La presenza riportata è attiva dall'inizio del trial, si
spegne una sola volta e non si riaccende, perché è la coda del posizionamento e non un
rilevamento. Con lo scarto ordinario di 20~s entrerebbe nella statistica, con quello di
120~s resta fuori.}
\label{fig:transitorio}
\end{figure}

\subsection{Validità dei file}
\label{sec:metriche-validita}

Un file entra nelle statistiche solo se è un trial. I controlli del fondo, le prove
tecniche e i file marcati come non validi nel registro restano fuori da ogni tabella
grazie alla lista di esclusione della Sezione~\ref{sec:c3-4-5}. Un file non valido non
viene cancellato, ma resta nel dataset con un nome diverso. Il caso tipico è il radar
rimasto muto per un cavo mosso, che si riconosce da centinaia di frame identici.

\section{Protocollo per l'\ldc{}}
\label{sec:protocollo-ld2420}

L'esemplare di \ldc\ in dotazione ha una portata di circa 2~m (Sezione~\ref{sec:c3-2-6}),
e la campagna sul secondo radar si è svolta per intero entro quel limite. Sono stati
ripetuti tutti i test dell'\ldb\ che non richiedono distanze maggiori, per un totale di 161
trial validi, e sono rimaste fuori solo le prove che le richiedono. Entrano invariate le due
più importanti per il progetto, la dose-risposta a 1~m e la persona sotto il banco.

Vanno dichiarate quattro deroghe alla parità di configurazione. I dati sono letti in
energy mode (Sezione~\ref{sec:c3-2-5}), l'unica modalità che esponga le energie per gate.
Le soglie sono tarate sul fondo misurato e il gate massimo è fissato a 6, per le ragioni
della Sezione~\ref{sec:c3-6}. Inoltre, il gate~0 è portato sopra la coda del rumore, nella
configurazione v2 (Sezione~\ref{sec:falsi-positivi}). Infine, il PIR non è cablato nei
test del solo \ldc, e i suoi valori sono marcati come assenti (Sezione~\ref{sec:c3-4-3}).
Poiché l'esemplare in dotazione è con ogni probabilità difettoso, i risultati ottenuti descrivono quel singolo modulo e non vanno
generalizzati al modello \ldc.

Due accorgimenti operativi completano il protocollo. Dopo ogni accensione nessuno resta
entro 2~m dal modulo per 90~s, perché il fondo viene stimato all'avvio e un corpo vicino lo
falserebbe. Prima di ogni sessione, poi, un minuto di stanza vuota serve a verificare che
il fondo per gate sia quello atteso.

\section{L'interfaccia web di supporto}
\label{sec:webui}
\label{chap:webui}

L'ESP32 pubblica i dati e una pagina web li mostra in tempo reale, li salva insieme alle
statistiche di sessione e li esporta in CSV. Questa sezione ne riassume requisiti, progetto, realizzazione e verifica,
e chiarisce il ruolo che ha avuto nella campagna.

\subsection{Requisiti e architettura}
\label{sec:webui-requisiti}

Dall'obbiettivo discendono cinque requisiti funzionali.
\begin{itemize}
  \item \textbf{R1:} l'ESP32 deve rendere disponibili i dati in rete.
  \item \textbf{R2:} la visualizzazione deve seguire il sensore alla cadenza
    dell'acquisizione, cioè 5~Hz.
  \item \textbf{R3:} l'applicazione deve calcolare le statistiche di sessione.
  \item \textbf{R4:} l'applicazione deve esportare i dati in CSV.
  \item \textbf{R5:} l'applicazione deve mostrare l'indice di vitalità calcolato a bordo
    (Capitolo~\ref{chap:vitalita}).
\end{itemize}
A questi si aggiungono due vincoli, che nascono dal contesto.
\begin{itemize}
  \item \textbf{N1:} tutto deve funzionare senza appoggiarsi ad alcuna rete esistente,
    perché in emergenza sismica non c'è connettività, né verso Internet né verso una rete
    locale (Sezione~\ref{sec:dipme}).
  \item \textbf{N2:} il CSV esportato dal browser deve avere lo stesso formato di quello
    prodotto dalla catena seriale, così che in tutta la tesi esista un solo formato dati.
\end{itemize}
La Tabella~\ref{tab:webui-opzioni} riassume le tre architetture considerate.

\begin{table}[!htbp]
\centering
\caption{Opzioni architetturali valutate per la web UI. È stata scelta la A.}
\label{tab:webui-opzioni}
\small
\begin{tabularx}{\linewidth}{@{}p{0.34\linewidth}XX@{}}
\toprule
\textbf{Architettura} & \textbf{Vantaggi} & \textbf{Svantaggi} \\
\midrule
\textbf{A. Sito servito dall'ESP32} (web server e WebSocket a bordo) &
  nessuna infrastruttura, funziona offline, coerente con lo scenario di emergenza &
  RAM e flash limitate, pochi client simultanei \\
B. ESP32, broker MQTT, server con database &
  scalabile, storico illimitato, più vicino a una piattaforma reale &
  richiede broker e server sempre accesi, infrastruttura eccessiva \\
C. Piattaforma cloud (ThingSpeak, Blynk) &
  pronta all'uso &
  non funziona offline, cadenza limitata a circa 15~s, dati fuori controllo \\
\bottomrule
\end{tabularx}
\end{table}

È stata scelta l'opzione~A, perché un nodo che serve la propria dashboard senza alcuna
rete esterna è un caso concreto di scenario senza connessione, cioè proprio lo scenario
del progetto. L'opzione~C viola il vincolo N1, mentre la~B lo rispetta solo con un server
locale che in emergenza non ci sarebbe. Il volume dei dati è piccolo, una trentina di
valori a 5~Hz per circa 2~kB/s, e a bordo non serve persistenza, perché lo storico di
sessione si accumula nel browser ed è il browser a produrre il CSV. L'ESP32 crea quindi una
propria rete WiFi in modalità Access Point e non tenta mai di collegarsi ad altre.
L'opzione~B resta documentata come alternativa valutata e scartata.

\subsection{Progetto e realizzazione}
\label{sec:webui-progetto}
\label{sec:webui-piano}
\label{sec:webui-lezioni}
\label{sec:webui-riferimenti}

\paragraph{Firmware} Il firmware del nodo è scritto in C++ con il framework Arduino, come
il logger della Sezione~\ref{sec:c3-4-2}, di cui riprende per intero la lettura dei sensori
aggiungendo la parte di rete. Il server web e il canale WebSocket sono forniti dalla
libreria ESP Async WebServer, nella versione mantenuta dal progetto
ESP32Async~\cite{esp32async}, perché quella originale non compila più con le versioni
attuali del supporto ESP32 per Arduino. La libreria è asincrona, cioè gestisce le
richieste dei browser in un task separato, eseguito in parallelo al ciclo principale del
programma. In questo modo la lettura del radar prosegue regolarmente a 5~Hz anche mentre
un browser si collega o scarica la pagina. La pagina è incorporata nel firmware come dati
compressi salvati in flash, per circa 80~kB. Per farli stare serve lo schema di partizione
che riserva 3~MB all'applicazione, ed è l'unica impostazione da cambiare rispetto al
logger. Un piccolo server DNS a bordo, infine, risponde con l'indirizzo del nodo a
qualunque nome venga richiesto, così la dashboard si apre da sola appena ci si collega alla
rete del nodo, senza bisogno di conoscerne l'indirizzo IP.

\paragraph{Pagina} La pagina, cioè la parte che gira nel browser, è composta da tre file
(HTML, CSS e JavaScript) per circa 800 righe, con il codice scritto in JavaScript puro,
senza librerie di sviluppo come Angular o React e senza strumenti di compilazione, quindi
i file scritti sono esattamente quelli eseguiti dal browser. Il tempo reale è affidato al
canale WebSocket, disponibile in ogni browser senza librerie, e l'unica dipendenza è la
libreria di grafici Chart.js~\cite{chartjs}, servita dall'ESP32 insieme alla pagina. Come
riferimento di interfaccia è stato esaminato un configuratore web open source per
l'\ldb~\cite{ld2410configurator}, da cui è ripresa la rappresentazione a barre delle
energie per gate con la soglia sovrapposta. La pagina è divisa in quattro aree
(Figure~\ref{fig:dashboard} e~\ref{fig:dashboard-sessione}). La prima
mostra lo stato istantaneo dei due sensori. La seconda contiene la gauge dell'indice di
vitalità, con i colori del semaforo (Sezione~\ref{sec:vitalita-classi}). La terza
raccoglie i tre grafici degli ultimi 60~s, con l'energia, le barre per gate con le soglie
lette dal modulo all'avvio e la timeline della presenza di radar e PIR sulla stessa base
dei tempi. La quarta gestisce la sessione di registrazione, con gli stessi metadati
sperimentali della catena seriale. Le soglie, deliberatamente, non si possono modificare
dalla pagina, perché devono restare fisse per tutta la campagna
(Sezione~\ref{sec:protocollo-configurazione}).

\paragraph{Protocollo dati} Il nodo trasmette i dati ai browser collegati attraverso il
canale WebSocket, con un messaggio per ogni campione, cioè cinque volte al secondo. Ogni
messaggio è un oggetto JSON che contiene tutte le grandezze del CSV del logger più le due
calcolate a bordo (Tabella~\ref{tab:webui-messaggio}). Il nodo conserva inoltre uno storico
di 60~s e lo invia a ogni client che si collega, così il grafico è già pieno anche dopo un
aggiornamento della pagina. Le statistiche sono calcolate nel browser una volta al secondo.
Il CSV esportato riporta, nello stesso ordine, le colonne del logger e i metadati dello
script di acquisizione, più due colonne in coda con l'indice e la classe calcolati dal
firmware (Sezione~\ref{sec:c3-4-3}). Sono proprio questi i dati usati per verificare il
porting (Sezione~\ref{sec:vitalita-bordo}).

\begin{figure}[!htbp]
\centering
\includegraphics[width=\linewidth]{Dashboard}
\caption{La dashboard servita dal nodo, con la testata con lo stato di presenza, lo stato
live dei due sensori (A), l'indice di vitalità calcolato a bordo (B) e i tre grafici degli
ultimi 60~s (C). I dati sono quelli di una sessione registrata dall'ESP32 con una persona
immobile a 1~m, riprodotti a 5~Hz nella stessa pagina. Il radar resta attivo per tutto il
tempo, mentre il PIR non rileva nulla.}
\label{fig:dashboard}
\end{figure}

\begin{figure}[!htbp]
\centering
\includegraphics[width=\linewidth]{Dashboard_sessione}
\caption{La parte inferiore della dashboard in una sessione dal vivo di 60~s, con una
persona davanti al sensore. A sinistra il riquadro D con i metadati della sessione e le
statistiche calcolate nel browser, a destra la diagnostica dell'ESP32 con la build,
l'indirizzo, la memoria libera, i parametri letti dal radar e l'impronta della pagina.}
\label{fig:dashboard-sessione}
\end{figure}

\begin{table}[!htbp]
\centering
\caption{Campi del messaggio WebSocket inviato a ogni client collegato, uno per campione
(5~Hz), con un valore di esempio.}
\label{tab:webui-messaggio}
\small
\begin{tabularx}{\linewidth}{@{}>{\raggedright\arraybackslash}p{0.24\linewidth}>{\raggedright\arraybackslash}p{0.33\linewidth}X@{}}
\toprule
\textbf{Campo} & \textbf{Esempio} & \textbf{Significato} \\
\midrule
\texttt{t} & \texttt{123456} & istante del campione, in millisecondi dall'accensione del nodo \\
\texttt{radar\_ok} & \texttt{1} & 1 se nel periodo è arrivato un frame valido dal radar \\
\texttt{presence}, \texttt{moving}, \texttt{still} & \texttt{1, 1, 0} & presenza complessiva, bersaglio in movimento, bersaglio fermo \\
\texttt{mdist}, \texttt{sdist} & \texttt{145, 0} & distanza del bersaglio in movimento e di quello fermo, in cm \\
\texttt{menergy}, \texttt{senergy} & \texttt{72, 0} & energia del bersaglio in movimento e di quello fermo, 0--100 \\
\texttt{pir} & \texttt{1} & uscita del PIR, letta nello stesso campione \\
\texttt{gates\_m} & \texttt{[0,5,60,10,0,0,0,0,0]} & energia per gate 0--8, canale di movimento \\
\texttt{gates\_s} & \texttt{[0,0,45,8,0,0,0,0,0]} & energia per gate 0--8, canale stazionario \\
\texttt{light}, \texttt{out} & \texttt{24, 1} & livello del fotosensore e stato del pin OUT \\
\texttt{vitality} & \texttt{54} & indice di vitalità calcolato a bordo, 0--100 \\
\texttt{vitality\_class} & \texttt{vitalita\_moderata} & classe dell'indice di vitalità \\
\bottomrule
\end{tabularx}
\end{table}

\paragraph{Realizzazione} Lo sviluppo è stato organizzato in cinque passi, ciascuno con un
criterio di accettazione da verificare prima di passare al successivo.
\begin{enumerate}
  \item \textbf{Rete e pagina statica:} il passo era superato quando la pagina si apriva da
    telefono e da portatile.
  \item \textbf{Lettura dei sensori e WebSocket:} i valori in pagina dovevano seguire il
    sensore senza accumulare latenza.
  \item \textbf{Grafici e storico:} spostando il soggetto di 0,75~m, le barre per gate
    dovevano indicare il gate corretto.
  \item \textbf{Sessione ed esportazione:} il CSV del browser è stato verificato come
    descritto nella Sezione~\ref{sec:webui-verifica}.
  \item \textbf{Indice di vitalità a bordo:} a regime l'indice doveva coincidere entro
    $\pm 1$ con quello calcolato in Python (Sezione~\ref{sec:vitalita-bordo}).
\end{enumerate}
Un problema incontrato lungo la strada merita di essere segnalato a chi ripete il lavoro.
I callback del server asincrono girano su un task diverso dal ciclo principale, e una
stampa di log alla connessione di un client ha corrotto una riga CSV mentre veniva
scritta. Per questo gli eventi di rete passano ora da una coda e vengono stampati dal
ciclo principale fra una riga e l'altra.

\subsection{Verifica di equivalenza}
\label{sec:webui-verifica}

Il vincolo N2 lega l'interfaccia al resto della tesi, e la sua verifica è stata impostata
come un test di equivalenza. La stessa sessione è stata registrata contemporaneamente in
due modi (Figura~\ref{fig:percorsi-dati}), dal browser con la funzione di esportazione e
dalla seriale con lo script di acquisizione lanciato in parallelo sulla porta del nodo. I
due file sono stati poi confrontati campione per campione, allineandoli sul tempo
dell'ESP32. Sui 1346 campioni comuni, cioè sull'intera sovrapposizione temporale, tutte le
colonne coincidono in ogni valore, comprese le energie per gate e le distanze, e lo script
di analisi legge entrambi i file producendo le stesse metriche. Non si tratta quindi di due
formati compatibili, ma dello stesso dato raggiunto per due vie.

\begin{figure}[!htbp]
\centering
\includegraphics[width=\linewidth]{Percorsi_dati}
\caption{I due percorsi dei dati del firmware della web UI. Sopra il percorso seriale, in
cui lo script di acquisizione aggiunge i metadati sperimentali. Sotto la web UI, servita dal nodo in modalità Access Point, che esporta dal browser lo stesso CSV. I due file coincidono
campione per campione e sono letti dallo stesso script di analisi.}
\label{fig:percorsi-dati}
\end{figure}

\subsection{Ruolo nella campagna}
\label{sec:webui-ruolo}

Nella campagna l'interfaccia ha svolto due funzioni.
\begin{itemize}
  \item \textbf{Controllo in tempo reale:} durante il posizionamento dei sensori e del
    soggetto permette di verificare, prima di registrare, che ogni sensore veda ciò che
    deve vedere. Il caso più evidente è il corridoio, dove la timeline della presenza di
    radar e PIR mostra dal vivo il risultato del confronto.
  \item \textbf{Acquisizione della prova di portata:} la prova in corridoio
    (Sezione~\ref{sec:portata-massima}) è stata acquisita con il firmware della web UI, e
    grazie alla verifica di equivalenza quei dati sono entrati nella catena di analisi
    senza alcuna modifica.
\end{itemize}
Il calcolo delle metriche di confronto resta invece affidato alla catena di analisi sui file (Sezione~\ref{sec:metodo-approccio}).

\chapter{Confronto sperimentale PIR vs mmWave}
\label{chap:confronto}

Questo capitolo riporta i risultati della campagna sperimentale, svolta sui tre sensori, con 444 trial validi e 32,1 ore di acquisizione. Protocollo, metriche e convenzioni sono quelli del Capitolo~\ref{chap:metodologia}. Ogni risultato è accompagnato dalla sua dispersione su ripetizioni indipendenti. Il rilevamento del respiro poggia sulla stessa campagna ed è trattato nel Capitolo~\ref{chap:vitalita} insieme all'indice di vitalità.

Il capitolo è organizzato come segue. La Sezione~\ref{sec:falsi-positivi} misura i falsi positivi a stanza vuota. La Sezione~\ref{sec:persona-ferma} affronta la persona immobile e spiega che cosa misura davvero il bit del PIR, mentre la Sezione~\ref{sec:dose-risposta} mostra come la sua risposta dipenda dalla quantità e dal tipo di movimento più che dalla distanza. La Sezione~\ref{sec:sotto-banco} porta il confronto nello scenario del progetto, con la persona sotto il banco. Le Sezioni~\ref{sec:distanza} e~\ref{sec:latenza} caratterizzano l'accuratezza della distanza, l'energia del bersaglio e le latenze di rilevamento e di rilascio. La Sezione~\ref{sec:angolare} misura la copertura angolare e ne ricava le conseguenze sulla selettività fra banchi adiacenti, e la Sezione~\ref{sec:due-persone} esamina il caso di due persone davanti al radar. La Sezione~\ref{sec:ostacoli} quantifica l'attenuazione dei materiali interposti e la Sezione~\ref{sec:portata-massima} la portata massima dei tre sensori in corridoio.

\section{Stanza vuota: falsi positivi}
\label{sec:falsi-positivi}

L'acquisizione parte con l'operatore che esce subito dalla stanza, e sul resto del
file si contano le transizioni verso la presenza. Su 7,03~h di stanza vuota (6,55~h di notte e 0,48~h di giorno) né l'\ldb\ né il PIR hanno prodotto alcun
falso positivo. Con zero eventi osservati si riporta il limite superiore $3/T$ della
Sezione~\ref{sec:metriche-zero}, $\leq 0{,}43$ eventi/h per entrambi i sensori.

L'\ldc\ nella stessa stanza si comporta in modo molto diverso
(Figura~\ref{fig:falsi-positivi}). Con le soglie già tarate sul fondo misurato (Sezione~\ref{sec:c3-6}) produce 26,2 riaccensioni all'ora in una notte di 7~h
e dichiara la presenza nel 26\,\% del tempo senza nessuno nella stanza. Con la
successiva correzione del gate~0 scende a 5,4 eventi/h. Quest'ultimo valore viene però da un'ora di giorno e non da una notte.

\begin{figure}[htbp]
\centering
\includegraphics[width=\linewidth]{Falsi_positivi}
\caption{Falsi positivi a stanza vuota, di notte, in scala logaritmica. L'\ldb\ a
soglie di fabbrica non produce eventi in 6,6~h, quindi è riportato il limite superiore
$3/T$. L'\ldc\ con le soglie tarate sul fondo misurato produce 26 riaccensioni/h in
7~h di notte, che scendono a 6/h con il gate~0 portato sopra la coda del rumore,
misurati però in 1~h di giorno. I
conteggi della figura usano lo scarto di 60~s degli script e differiscono di poco da
quelli del testo, che usano la finestra del registro.}
\label{fig:falsi-positivi}

\end{figure}

\section{La persona ferma}
\label{sec:persona-ferma}

È lo scenario che risponde alla domanda di partenza. Il soggetto siede immobile a
2,30~m dai sensori, con la sola respirazione spontanea, per cinque trial da 302~s
utili a 27\,\textdegree C. I 7554 campioni con presenza reale danno il risultato della
Tabella~\ref{tab:fermo-seduto}.

\begin{table}[htbp]
\centering
\caption{Persona ferma seduta a 2,30~m, cinque trial da 302~s.}
\label{tab:fermo-seduto}

\small
\begin{tabular}{@{}lrr@{}}
\toprule
 & \textbf{PIR \pir} & \textbf{Radar \ldb} \\
\midrule
Falsi negativi & $99{,}78 \pm 0{,}49\,\%$ & $\mathbf{0{,}00 \pm 0{,}00\,\%}$ \\
Trial con falsi negativi al 100\,\% & 4 su 5 & 0 su 5 \\
Campioni con presenza rilevata & $\sim$17 su 7554 & 7554 su 7554 \\
\bottomrule
\end{tabular}
\end{table}

Con una persona viva a 2,3~m in campo libero, il PIR è rimasto muto per cinque minuti in quattro trial su cinque. Non è un esito inatteso, perché è ciò che il
principio del piroelettrico prevede (Sezione~\ref{subsec:pir-principio}), e qui ne viene misurata la dispersione su cinque ripetizioni. La
Figura~\ref{fig:timeline-sotto-banco} mostra l'andamento nel tempo di un trial in questa
geometria e di uno nella geometria del progetto DIPME (Sezione~\ref{sec:sotto-banco}).

\begin{figure}[htbp]
\centering
\includegraphics[width=\linewidth]{Timeline_sotto_banco}
\caption{Presenza dichiarata dai due sensori su una persona immobile, in due
geometrie. A sinistra il soggetto seduto a 2,3~m, a destra rannicchiato sotto il
banco a circa 60~cm, cinque minuti ciascuno. Il radar è continuo, il PIR produce al
più un impulso isolato. La terza riga è la distanza riportata dal radar sul bersaglio
fermo, con una deviazione standard media di 1,5~cm da seduto e molto più rumorosa sotto il piano.}
\label{fig:timeline-sotto-banco}

\end{figure}

\subsection{Che cosa misura il bit del PIR}
\label{sec:interpretazione-pir}
\label{sec:jumper}

La percentuale di campioni in cui il PIR è alto va letta sapendo come nasce. In
posizione L l'uscita è un monostabile a durata fissa. Su tutte le sessioni acquisite
in quella configurazione sono stati contati 195 impulsi con durata media
$3{,}46 \pm 0{,}09$~s, tutti fra 3,4 e 3,6~s e nessuno oltre 5~s, anche nei trial
in cui il soggetto si muoveva per cinque minuti (Figura~\ref{fig:impulsi-pir}). La
percentuale è quindi il numero di eventi moltiplicato per la ritenuta, cioè un
conteggio di eventi travestito da frazione di tempo. In posizione H gli impulsi si
concatenano finché il movimento continua, fino a coprire un intero trial di 62~s, e la
percentuale diventa allora una vera frazione di tempo con movimento rilevato. Il significato delle due modalità è descritto nella Sezione~\ref{subsec:pir-dati}.

\begin{figure}[htbp]
\centering
\includegraphics[width=\linewidth]{Impulsi_pir}
\caption{Struttura degli impulsi del PIR nelle due serie con movimento continuo, il cammino sul posto a 1~m e i movimenti sotto il banco, che in L contengono 181 dei 195 impulsi. In L la durata è fissa, in H il ritrigger concatena gli impulsi.}
\label{fig:impulsi-pir}

\end{figure}

Come anticipato nella Sezione~\ref{sec:c3-3-3}, la prima parte della campagna è stata acquisita in posizione L, e le serie del cammino sul posto a 1~m e dei micro-movimenti sotto il banco sono state ripetute in H e sono quelle riportate in questo capitolo. Anche se le altre serie non richiedevano una ripetizione, perché il ponticello agisce solo dopo un trigger e dove il PIR produce al più eventi isolati può solo allungarli di poco, sono stati comunque ripetuti in H, per verifica, il cammino a 2~m (0,00\,\% in entrambe le posizioni) e l'immobilità sotto il banco (0,10\,\% in L, 1,32\,\% in H).

\section{Il rilevamento al variare del movimento}
\label{sec:dose-risposta}
\label{sec:pir-movimento}

Le due prove di questa sezione variano una sola grandezza alla volta. La prima varia la quantità di movimento a distanza fissa, con una curva dose-risposta a 1~m, la seconda il tipo di movimento a parità di distanza, confrontando il cammino sul posto con l'attraversamento.

\subsection{Quantità di movimento}
\label{sec:dose-risposta-1m}

Il soggetto è in piedi a 1~m, la distanza più favorevole al PIR fra quelle provate, così che il confronto non sia costruito a suo sfavore. Il ponticello è in H, il montaggio è lo stesso in tutte le condizioni e la quantità di movimento assume tre livelli, con cinque trial ciascuno (Tabella~\ref{tab:dose-risposta} e Figura~\ref{fig:dose-risposta}).

\begin{table}[htbp]
\centering
\caption{Curva dose-risposta a 1~m, ponticello in H, cinque trial per condizione. Le due colonne
di destra sono grandezze del radar, riprese nel Capitolo~\ref{chap:vitalita}.}
\label{tab:dose-risposta}

\small
\begin{tabular}{@{}lrrrrr@{}}
\toprule
\textbf{Condizione} & \textbf{PIR} & \textbf{dev.\ rel.} & \textbf{Radar} &
\textbf{Energia mov.} & \textbf{Disp.\ dist.} \\
\midrule
Immobile          & $1{,}52 \pm 1{,}03\,\%$   & $68\,\%$ & $\mathbf{100\,\%}$ & 68,6 & 21,3~cm \\
Micro-movimenti   & $51{,}60 \pm 21{,}72\,\%$ & $42\,\%$ & $\mathbf{100\,\%}$ & 84,4 & 16,4~cm \\
Cammino sul posto & $85{,}20 \pm 11{,}52\,\%$ & $14\,\%$ & $\mathbf{100\,\%}$ & 99,3 & 10,4~cm \\
\bottomrule
\end{tabular}
\end{table}

\begin{figure}[htbp]
\centering
\includegraphics[width=\linewidth]{Dose_risposta}
\caption{La curva dose-risposta a 1~m. Il PIR passa dall'1,5 all'85\,\%,
con la dispersione massima nella zona intermedia. I due radar restano al 100\,\% in tutte e tre le condizioni. L'\ldc\ è stato provato nella stessa geometria in una sessione separata, e a stanza vuota non è riportato perché il suo valore dipende dalla configurazione (Sezione~\ref{sec:falsi-positivi}).}
\label{fig:dose-risposta}

\end{figure}

Il fallimento del PIR sulla persona ferma dipende quindi dall'immobilità del soggetto e non dalla distanza, dalla configurazione o dalla taratura, che fra le righe restano invariate. Da immobile il PIR produce due soli impulsi in 1010~s di acquisizione.

Per un'applicazione di soccorso il dato più rilevante è la dispersione della condizione intermedia, $\pm 21{,}72$ punti su una media di 51,60, perché due trial identici per l'operatore danno risultati molto diversi. Vicino alla propria soglia il PIR diventa imprevedibile e non permette di interpretare correttamente l'assenza di segnale.

Le ultime due colonne della
Tabella~\ref{tab:dose-risposta} sono un risultato a sé. L'energia del bersaglio in movimento
cresce con la quantità di movimento, mentre la dispersione della distanza cala, perché un'eco più forte dà una stima più stabile.
Sono due grandezze continue e monotone misurate a geometria costante, e su di esse poggia l'indice di vitalità del Capitolo~\ref{chap:vitalita}. L'\ldc, provato nella stessa geometria, rileva la persona nel 100\,\% del tempo in tutte e tre le condizioni (Figura~\ref{fig:dose-risposta}), ma la sua energia non cresce con il movimento.

\subsection{Tipo di movimento}
\label{sec:attraversamento}

Il cammino sul posto è stato ripetuto a 1, 2, 3, 4 e 5~m, con cinque trial per distanza (in H a 1 e 2~m, in L da 3 a 5~m, dove il ponticello non cambia nulla). A 1~m il PIR rileva l'85,2\,\% del
tempo, mentre da 2~m in poi resta a zero. Sono venti
trial consecutivi al 100\,\% di falsi negativi con un soggetto in movimento continuo
davanti al sensore, e il crollo fra 1 e 2~m non è graduale. Questo esito è in aperto
contrasto con la portata di 3--7~m dichiarata dal costruttore~\cite{hcsr501}, ma da solo non
distingue fra due spiegazioni, cioè il tipo di movimento oppure un esemplare guasto o
tarato male. Per decidere, il soggetto ha ripetuto la prova camminando in modo
continuo trasversalmente all'asse del sensore, a distanza costante. Il resto è rimasto invariato (trimmer di sensibilità nella stessa posizione), con il ponticello in H e tre trial per distanza. Il risultato è in Figura~\ref{fig:tipo-movimento}.

\begin{figure}[htbp]
\centering
\includegraphics[width=\linewidth]{Tipo_movimento}
\caption{Il PIR con due tipi di movimento. A parità di distanza e sensibilità, il cammino sul posto è rilevato solo a 1~m (85\,\%) e mai da 2~m in su.
L'attraversamento trasversale è rilevato al 99--100\,\% da 1 a 5~m. Il radar è al
100\,\% con entrambi i movimenti.}
\label{fig:tipo-movimento}

\end{figure}

Un sensore che rileva un attraversamento a 5~m nel 98,7\,\% dei campioni non è né
guasto né poco sensibile, e la portata dichiarata risulta confermata. A decidere è quindi il tipo di movimento, come prevede la geometria della lente di Fresnel (Sezione~\ref{subsec:fresnel}). Lo conferma una prova a 1~m con quattro movimenti diversi, per i quali l'energia in movimento del radar è rimasta fra 96,5 e 98,1, cioè la stessa quantità di movimento. Il PIR è passato dal 2\,\% con il movimento delle braccia al 100\,\% con l'oscillazione laterale del busto da seduto, che è una traversata in miniatura. I due sensori misurano quindi grandezze diverse, il radar la quantità di movimento e il PIR il passaggio da una zona all'altra della lente.

Per il progetto la conseguenza è diretta. Il movimento per cui il PIR è progettato,
attraversare un varco o percorrere un corridoio, è l'opposto di quello di una persona
intrappolata sotto un arredo, che, se si muove, lo fa sul posto. Sotto il banco, a 60~cm, il PIR vede ancora questi movimenti (Sezione~\ref{sec:sotto-banco}), ma da 2~m in su non più, e resta cieco sulla persona immobile a qualunque distanza.

\section{Lo scenario DIPME: la persona sotto il banco}
\label{sec:sotto-banco}
\label{sec:gate-sotto-banco}

È lo scenario che replica la condizione del progetto. I sensori sono fissati sotto il
piano di un tavolo e il soggetto è rannicchiato sotto di esso a 50--60~cm
(Figura~\ref{fig:geometrie}b). Sono stati acquisiti cinque trial da 302~s utili in due
condizioni, immobilità e presenza di micro-movimenti, entrambe con il ponticello del
PIR in H (Figura~\ref{fig:dipme-tre-sensori}).

\begin{figure}[htbp]
\centering
\includegraphics[width=\linewidth]{Dipme_tre_sensori}
\caption{I tre sensori con la persona immobile e con micro-movimenti, in piedi a 1~m e sotto il banco, cinque trial per condizione. L'\ldc\ è stato provato nella stessa geometria, con soglie tarate e scarto di 90~s.}
\label{fig:dipme-tre-sensori}

\end{figure}

Sotto il banco il contrasto è ancora più netto che a 1~m. Con i micro-movimenti il PIR
rileva il $96{,}52 \pm 3{,}60\,\%$ del tempo, da immobile l'$1{,}32 \pm 0{,}86\,\%$.
A 60~cm il PIR è un rilevatore di movimento quasi perfetto e resta cieco alla persona
ferma. L'\ldb\ è al 100\,\% in entrambe le condizioni su 7554 campioni per
condizione. L'\ldc\ è al 100\,\% con i micro-movimenti e al $98{,}4 \pm 3{,}6\,\%$
da immobile, con un solo rilascio di 24~s in un trial.

\subsection{Il radar rileva bene ma localizza male}
\label{sec:sotto-banco-localizzazione}

Sotto il banco la distanza riportata sulla persona immobile ha una deviazione standard entro il trial di $24 \pm 9$~cm, contro $1{,}48 \pm 0{,}77$~cm per il soggetto seduto immobile a 2,3~m (Figura~\ref{fig:timeline-sotto-banco}). Il bersaglio è più vicino e ugualmente fermo, eppure la dispersione è circa sedici volte più ampia. Non dipende da piccoli movimenti del soggetto, perché con i micro-movimenti la dispersione scende a 14~cm. Lo stesso accade a 1~m, dove la dispersione cala con il movimento (Sezione~\ref{sec:dose-risposta-1m}). Per il progetto conta poco, perché al soccorritore serve sapere se sotto quel banco c'è qualcuno e non a quanti centimetri si trova. In questa geometria, però, la distanza non va usata per localizzare con precisione la persona.

Gli sbalzi verso 200--300~cm visibili nella Figura~\ref{fig:timeline-sotto-banco} riguardano dallo 0,3 al 7,4\,\% dei campioni. Durante gli sbalzi il PIR resta a zero, quindi non sono movimenti del soggetto, ma il bersaglio stazionario fittizio già visto nella Sezione~\ref{sec:c3-5-3}, che il canale riporta a una distanza sbagliata.

La lettura per gate indica il perché. La presenza della persona immobile è retta soprattutto da gate più lontani di quello in cui si trova. Sull'\ldb\ la reggono i gate 2 e 3 del canale stazionario, i cui gate 0 e 1 sono nulli per costruzione (Sezione~\ref{sec:c3-5-2}), mentre il canale in movimento ha il massimo nel gate~1, dove la persona si trova. Sull'\ldc\ la regge il gate~2 (70--140~cm) invece del gate~1. La causa più probabile sono i cammini multipli sotto il piano, cioè echi che rimbalzano su pavimento e piano del banco e arrivano come da più lontano. Resta però un'ipotesi, perché i moduli forniscono solo l'energia già elaborata per gate e non il segnale da cui si potrebbero separare i diversi percorsi (Sezione~\ref{subsec:fmcw}).

\section{Accuratezza della distanza ed energia}
\label{sec:distanza}

\subsection{Accuratezza della distanza}
\label{sec:distanza-accuratezza}

Per misurare quanto è precisa la distanza riportata dall'\ldb, il soggetto ha camminato sul posto a 1, 2, 3, 4 e 5~m dal sensore, con cinque trial da 62~s per ogni distanza. Per ogni trial si confronta la distanza media riportata dal radar con quella segnata sul pavimento. La Tabella~\ref{tab:distanza} riporta i risultati distanza per distanza e la Figura~\ref{fig:distanza-regressione} li mostra in un grafico.

\begin{table}[htbp]
\centering
\caption{Accuratezza della distanza riportata dall'\ldb, cinque trial da 62~s per distanza. ``Misurata'' è la media dei cinque trial, ``errore'' la differenza dalla distanza reale, ``residuo'' lo scarto dalla retta di regressione e ``disp.\ entro trial'' quanto oscilla la distanza dentro un trial. L'energia è discussa nella Sezione~\ref{sec:energia-distanza}.}
\label{tab:distanza}

\small
\setlength{\tabcolsep}{4pt}
\begin{tabular}{@{}rrrrrrr@{}}
\toprule
\textbf{Reale} & \textbf{Misurata} & \textbf{Errore} & \textbf{Err.\ rel.} &
\textbf{Residuo} & \textbf{Disp.\ entro trial} & \textbf{Energia} \\
\midrule
1~m & $99{,}48 \pm 2{,}61$~cm  & $-0{,}52$~cm & $-0{,}5\,\%$ & $-3{,}0$~cm & $10{,}4$~cm & 99,0 \\
2~m & $211{,}22 \pm 0{,}77$~cm & $+11{,}22$~cm & $+5{,}6\,\%$ & $+4{,}9$~cm & $14{,}3$~cm & 85,1 \\
3~m & $309{,}76 \pm 1{,}13$~cm & $+9{,}76$~cm & $+3{,}3\,\%$ & $-0{,}4$~cm & $17{,}5$~cm & 54,4 \\
4~m & $411{,}84 \pm 2{,}76$~cm & $+11{,}84$~cm & $+3{,}0\,\%$ & $-2{,}1$~cm & $26{,}4$~cm & 35,1 \\
5~m & $518{,}20 \pm 2{,}13$~cm & $+18{,}20$~cm & $+3{,}6\,\%$ & $+0{,}5$~cm & $11{,}7$~cm & 27,6 \\
\bottomrule
\end{tabular}
\end{table}

\begin{figure}[htbp]
\centering
\includegraphics[width=\linewidth]{Distanza_regressione}
\caption{Accuratezza della distanza dell'\ldb, 25 trial fra 1 e 5~m, con la retta di regressione calcolata sulle cinque medie per distanza.}
\label{fig:distanza-regressione}

\end{figure}

A 1~m il radar sbaglia di mezzo centimetro, a 5~m di 18~cm. L'errore cresce con la distanza perché il radar legge sistematicamente un po' più lungo del vero. La retta che interpola i punti è
\begin{equation*}
d_{\text{misurata}} = 1{,}0381 \cdot d_{\text{reale}} - 1{,}32~\text{cm},
\qquad R^2 = 0{,}99965.
\end{equation*}
La pendenza 1,0381 indica che il radar allunga ogni distanza del 3,8\,\%, mentre il termine costante di $-1{,}32$~cm è trascurabile. I punti si scostano dalla retta al più di 4,9~cm (colonna ``residuo''). La retta è calcolata sulle cinque medie per distanza, e sui 25 trial singoli $R^2$ scende a 0,99949 mentre lo scostamento massimo sale a 6,0~cm. Il radar è quindi molto lineare, e il suo errore si corregge quasi del tutto dividendo la lettura per 1,0381. I dati non permettono di dire se il 3,8\,\% dipenda dal modulo o dal metro con cui sono state segnate le distanze, e per stabilirlo servirebbe un distanziometro indipendente.

La colonna ``disp.\ entro trial'' indica invece quanto oscilla la distanza riportata durante un singolo trial, fra 10 e 26~cm. Dipende dal movimento del busto (Figura~\ref{fig:energia-portata}) ma anche dalla forza dell'eco, tanto che la persona immobile a 1~m dà 21~cm (Tabella~\ref{tab:dose-risposta}). Il soggetto seduto a 2,3~m scende a 1,5~cm, ma sul canale stazionario (Sezione~\ref{sec:persona-ferma}).

\subsection{Energia del bersaglio}
\label{sec:energia-distanza}

\begin{figure}[htbp]
\centering
\includegraphics[width=\linewidth]{Energia_distanza}
\caption{Energia del bersaglio in movimento e dispersione della distanza entro il trial al variare della distanza, 25 trial.}
\label{fig:energia-portata}

\end{figure}

La stessa serie mostra come cambia con la distanza l'energia del bersaglio. Scende da 99 a 1~m a 28 a 5~m (Figura~\ref{fig:energia-portata}), perché una persona più lontana restituisce un'eco più debole, ma molto più lentamente di quanto farebbe la potenza ricevuta, che a 5~m sarebbe centinaia di volte più bassa. È la conferma che l'energia è un indice elaborato dal modulo e non una potenza (Sezione~\ref{sec:c3-5-1}).

Ne seguono due conseguenze pratiche. A 5~m l'energia vale ancora 27,6, quasi il doppio della soglia di 15 dei gate lontani, quindi il radar vedrebbe anche oltre i 5~m della stanza, come conferma la prova in corridoio (Sezione~\ref{sec:portata-massima}). Inoltre la stessa persona, con lo stesso movimento, dà un'energia diversa a seconda della distanza, quindi un indice costruito sull'energia va tarato sull'installazione (Capitolo~\ref{chap:vitalita}).

\subsection{Il secondo radar e il PIR}
\label{sec:distanza-ld2420}

Il PIR non compare in questa sezione perché non misura né una distanza né un'energia, e il suo dato è un bit.

L'\ldc\ riporta una distanza solo sul bersaglio in movimento e, sull'esemplare in dotazione, solo da vicino. Nella stessa prova ripetuta da 0,5 a 2~m, con venti trial, la distanza riportata vale $59 \pm 5$~cm a 0,5~m e $121 \pm 9$~cm a 1~m, cioè 9 e 21~cm più del vero, contro il mezzo centimetro dell'\ldb\ a 1~m. La lettura inoltre si aggiorna di rado, da 4 a 38 valori diversi al minuto, mentre quella dell'\ldb\ segue il movimento di continuo. A 1,5 e 2~m resta ferma su un solo valore per tutto il trial, quindi non è una misura e una retta di regressione come quella dell'\ldb\ non si può costruire. Sulla persona immobile il modulo dichiara la presenza al 100\,\%, ma come distanza riporta 2--5~cm, il valore che usa quando non ha un bersaglio in movimento. È quanto prevede il manuale, secondo cui il modulo non misura la distanza dei corpi fermi~\cite[§8]{hilinkManualLD2420}.

\section{Latenze di rilevamento e di rilascio}
\label{sec:latenza}
\label{sec:rilascio}

L'acquisizione parte con la stanza vuota oppure con il soggetto presente, che
rispettivamente entra o esce a $t = 30$~s, istante annunciato dallo script di
acquisizione (Sezione~\ref{sec:protocollo-groundtruth}). Sono stati acquisiti dieci
trial all'ingresso e cinque all'uscita, tutti validi, con nessun sensore attivo prima
dell'evento. I risultati sono in Tabella~\ref{tab:latenze} e Figura~\ref{fig:latenze}.

\begin{table}[htbp]
\centering
\caption{Latenze all'ingresso e tempo di rilascio all'uscita. \ldb\ e PIR letti
nello stesso trial (10 all'ingresso, 5 all'uscita), con la differenza appaiata
calcolata trial per trial e nessuna riaccensione dopo il rilascio. L'\ldc\ è misurato
in una sessione separata (5 trial all'ingresso, 3 all'uscita) con il ritardo di
scomparsa portato da 30 a 5~s come sull'\ldb. Il PIR non era cablato, quindi manca la
differenza appaiata. Tutti i valori sono in secondi.}
\label{tab:latenze}

\small
\setlength{\tabcolsep}{4pt}
\begin{tabularx}{\linewidth}{@{}>{\raggedright\arraybackslash\hsize=0.8\hsize}Xrrr>{\raggedright\arraybackslash\hsize=1.2\hsize}X@{}}
\toprule
\textbf{Grandezza} & \textbf{\ldb} & \textbf{PIR} & \textbf{\ldc} & \textbf{Nota} \\
\midrule
Latenza all'ingresso & $5{,}36 \pm 0{,}30$ & $5{,}96 \pm 0{,}31$ & $7{,}28 \pm 1{,}40$ &
  latenza operativa, non del solo sensore \\
\addlinespace
Differenza appaiata radar $-$ PIR & \multicolumn{2}{c}{$\mathbf{-0{,}60 \pm 0{,}35}$} & --- &
  il radar rileva per primo in 10 trial su 10 \\
\addlinespace
Tempo di rilascio all'uscita & $\mathbf{18{,}36 \pm 0{,}54}$ & $12{,}96 \pm 1{,}40$ & $7{,}7 \pm 0{,}5$ &
  \ldb\ più lento del PIR di $5{,}40 \pm 1{,}86$ (appaiata); \ldc\ con ritardo di 5~s, prima caduta \\
\bottomrule
\end{tabularx}
\end{table}

\begin{figure}[htbp]
\centering
\includegraphics[width=\linewidth]{Latenze}
\caption{A sinistra le latenze all'ingresso, appaiate per \ldb\ e PIR e da una sessione separata per l'\ldc. A destra il tempo di rilascio all'uscita, dove l'\ldb\ impiega 18~s a dichiarare la stanza vuota con un ritardo di scomparsa di 5~s.}
\label{fig:latenze}

\end{figure}

\paragraph{Ingresso.} Le latenze intorno a 5--6~s non misurano solo i sensori, perché comprendono anche il tragitto del soggetto dalla porta e il suo tempo di reazione all'annuncio. Questi due tempi però sono uguali per i due sensori, che osservano lo stesso ingresso nello stesso trial. Sottraendo trial per trial la latenza del PIR da quella del radar si cancellano, e resta solo la differenza fra i sensori (Sezione~\ref{sec:metriche-elenco}). Vale $-0{,}60 \pm 0{,}35$~s, cioè il radar rileva in media 0,6~s prima del PIR, con un errore standard della media di 0,11~s. Il radar arriva primo in tutti e dieci i trial, e se i due sensori fossero equivalenti la probabilità che accada per caso sarebbe di circa 2 su mille. A 5~Hz un campione dura 0,2~s e le singole differenze vanno da 0,2 a 1,4~s, quindi il singolo trial è poco preciso, e la conclusione poggia sul fatto che il radar arrivi primo ogni volta. Ci si
attendeva il contrario, perché attraversare una porta è lo stimolo trasversale per
cui la lente di Fresnel è progettata. Una spiegazione plausibile, non dimostrata, è
che il radar agganci il soggetto ancora in avvicinamento grazie ai 6~m di portata.

\paragraph{Uscita.} Qui il radar è nettamente più lento, di $5{,}40 \pm 1{,}86$~s nella differenza appaiata, con un errore standard di 0,83~s, cioè un effetto più di sei volte la sua incertezza già con cinque trial. Il
tempo di rilascio include il tempo di uscita dal campo, e il PIR può servire da
cronometro, perché la sua uscita ricade $3{,}46$~s dopo l'ultimo trigger. Il soggetto esce
quindi dal campo circa 9,5~s dopo l'annuncio, e la coda propria del radar vale circa
$18{,}36 - 9{,}5 \approx 8{,}9$~s. È una stima, perché i campi dei due sensori non
coincidono, ma concorda con altre due osservazioni indipendenti (10 e 12~s). Il
timeout di presenza letto dal modulo vale invece 5~s (Tabella~\ref{tab:c3-6}), quindi
il parametro configurato non descrive il comportamento osservato. Per il progetto
significa che una stanza appena liberata risulta occupata per quasi dieci secondi.

\paragraph{Il secondo radar.} Sullo stesso percorso l'\ldc\ rileva l'ingresso in
$7{,}28 \pm 1{,}40$~s su cinque trial, quasi due secondi dopo l'\ldb\ e con una
dispersione quattro volte maggiore. L'acquisizione coincide con il primo superamento della soglia di attivazione sul gate~2, cioè avviene entro circa 1~m, mentre l'\ldb\ verosimilmente aggancia il soggetto già in avvicinamento. Il rilascio va
letto insieme al suo parametro. Con il ritardo di scomparsa di fabbrica, 30~s, il
modulo restava in presenza per 87--120~s dopo l'uscita, perché basta un superamento
della soglia di mantenimento ogni mezzo minuto a riarmare il ritardo. Con il ritardo portato a 5~s, lo stesso valore dell'\ldb, il rilascio scende a $7{,}7 \pm 0{,}5$~s, e la prova mostra anche che il parametro è espresso in secondi. Il valore è la prima caduta della presenza, perché in due trial su tre il modulo si è poi riacceso per qualche secondo a stanza vuota, come nei falsi positivi della Sezione~\ref{sec:falsi-positivi}. A pari ritardo l'\ldc\ rilascia quindi in 7,7~s contro 18,4, ma la differenza viene dalla portata, come all'ingresso. Il soggetto esce prima dal campo più corto, e il modulo non è più rapido a dichiarare la stanza vuota.

\section{Copertura angolare e selettività fra banchi}
\label{sec:angolare}

Nell'architettura del progetto ogni arredo ha il proprio sensore, che deve rilevare chi sta sotto quel banco e non chi sta sotto quello accanto. Per restringere il campo l'\ldb\ offre una soluzione, il gate massimo. Con la regola documentata di 75~cm per gate, un gate massimo di 2 limita la portata a 150~cm (Tabella~\ref{tab:c3-6}). Il gate però taglia la distanza e non l'angolo. Le prove di questa sezione rispondono quindi a due domande. La prima è se il taglio a 150~cm funziona davvero, la seconda se entro quella distanza il radar vede una persona che sta di lato. La Sezione~\ref{sec:selettivita} prova il taglio con persone ferme, la~\ref{sec:angolare-copertura} misura fino a quale angolo il radar vede, e la~\ref{sec:selettivita-configurazione} mette insieme i due risultati.

\subsection{Selettività con gate massimo 2}
\label{sec:selettivita}

Sono stati provati quattro scenari con
soggetti sempre fermi, con tre trial da 182~s ciascuno (Figura~\ref{fig:selettivita}).
L'occupante a 1~m sull'asse è rilevato al 100\,\%, mentre una persona a 3~m sull'asse,
oltre il taglio, è rilevata allo 0,00\,\%. Anche una persona a 1~m a 90\textdegree\ è rilevata allo 0,00\,\% a regime, e il $24{,}8 \pm 19{,}6\,\%$ che si ottiene con lo scarto ordinario di 20~s è la coda del posizionamento della Figura~\ref{fig:transitorio}. L'\ldc\ nella stessa posizione dà lo stesso esito (0\,\% in tre trial), coerente con la sua copertura angolare (Sezione~\ref{sec:angolare-copertura}). Con il vicino accanto, l'occupante resta rilevato al 100\,\%, e distanza ed energia cambiano di circa 4~cm e 4 punti, entro la variabilità fra trial.

\begin{figure}[htbp]
\centering
\includegraphics[width=\linewidth]{Selettivita}
\caption{Selettività spaziale con gate massimo 2 (150~cm), tre trial per scenario.
Le barre chiare usano lo scarto ordinario dei primi 20~s, quelle scure lo scarto di
120~s (a regime).}
\label{fig:selettivita}

\end{figure}

\subsection{Copertura angolare}
\label{sec:angolare-copertura}

Il soggetto siede su un arco di raggio 1~m a sei
azimut (0, 45, 60, 75, 90 e 120\textdegree), con la sedia rivolta verso i sensori a
ogni posizione, e oscilla lateralmente il busto con ampiezza costante. Il gate massimo
vale 2 e i trial sono da uno a tre per angolo (Figura~\ref{fig:angolare-due-radar}). L'\ldb\ rileva al
100\,\% fino a 90\textdegree\ compresi, con l'energia che cala di due punti soli fra
l'asse e i 90\textdegree\ (da 96,5 a 94,3), e crolla fra 90 e 120\textdegree. La
copertura eccede quindi nettamente i $\pm 60$\textdegree\ del manuale~\cite{hilinkManualV104}. Il PIR collassa prima,
dal 100\,\% sull'asse al 12\,\% a 90\textdegree\ e a zero a 120\textdegree. Fra 45 e
75\textdegree\ la dispersione fra trial è però dell'ordine dell'effetto, e a
45\textdegree\ i tre trial danno 100, 100 e 33\,\%, per cui le sole affermazioni
sostenibili riguardano i tre estremi. L'irripetibilità è essa stessa un risultato, lo stesso
comportamento già visto nella zona intermedia della dose-risposta. L'\ldc\ a 1~m ha il fascio utile fino a
circa 75\textdegree\ ed è al fondo a 90\textdegree, più stretto dell'\ldb\ e più largo
dei $\pm 45$ o $\pm 60$\textdegree\ del suo manuale~\cite[§1.1 e §5.2]{hilinkManualLD2420}. Il confine è coerente anche con
il margine ridotto del trasmettitore dell'esemplare (Sezione~\ref{sec:c3-2-6}).

\begin{figure}[htbp]
\centering
\includegraphics[width=\linewidth]{Angolare}
\caption{Copertura angolare a 1~m, busto laterale da seduto, gate massimo 2. A: \ldb\ e PIR. B: esemplare di \ldc.}
\label{fig:angolare-due-radar}

\end{figure}

\subsection{Selettività per configurazione}
\label{sec:selettivita-configurazione}

Il vicino a 1~m a 90\textdegree\ compare in entrambe le prove, con lo stesso gate. Fermo (Sezione~\ref{sec:selettivita}) è rilevato allo 0\,\%, mentre in movimento (Sezione~\ref{sec:angolare-copertura}) è rilevato al 100\,\%. Quella che nella prima prova sembrava selettività era quindi la debolezza dell'eco di un corpo immobile. A 90\textdegree\ il diagramma di irradiazione non attenua abbastanza da escludere chi si muove, e un banco vuoto accanto a una persona in movimento può risultare occupato. La separazione fra arredi adiacenti va cercata nel montaggio fisico (orientamento verso
il basso, schermatura del lobo laterale) e non nei parametri del modulo. L'esemplare di \ldc, al fondo a 90\textdegree, non vedrebbe il vicino, ma il merito è del suo trasmettitore debole più che dell'antenna (Sezione~\ref{sec:c3-2-6}).

\section{Due persone davanti al radar}
\label{sec:due-persone}

L'\ldb\ misura la distanza dei bersagli ma non la loro direzione, quindi non può separare due persone in base a dove si trovano. Ha però due canali, uno per il bersaglio in movimento e uno per quello fermo, ciascuno con la propria distanza (Sezione~\ref{sec:c3-1-4}). La prova verifica se questi due canali bastano a riportare due persone insieme. La persona A sta ferma a 2~m e la persona B a 4~m, in tre geometrie con tre trial da 122~s utili ciascuna (Tabella~\ref{tab:due-persone}, Figura~\ref{fig:due-persone}). Nelle prime due le persone sono sfalsate di lato di 50~cm, perché A non copra B, e B cammina oppure sta ferma. Nella terza B cammina in fila dietro A. Da sola, ferma a 4~m, B è rilevata al 100\,\% (a $388{,}9 \pm 3{,}7$~cm), quindi quando nelle prove non compare è perché è coperta da A.

\begin{table}[htbp]
\centering
\caption{Capacità di separare due persone, su tre geometrie. ``B riportata'' è la
frazione di campioni in cui la persona lontana compare su almeno uno dei due canali,
riconosciuta dalla banda di distanza. Tre trial per geometria, 1833 campioni utili
ciascuna.}
\label{tab:due-persone}

\small
\begin{tabular}{@{}lrrrrr@{}}
\toprule
\textbf{Geometria} & \textbf{Entrambi i} & \textbf{B} & \textbf{Dist.} &
\textbf{Dist.} & \textbf{Presenza} \\
 & \textbf{canali attivi} & \textbf{riportata} & \textbf{staz.} & \textbf{mov.} &
\textbf{radar} \\
\midrule
B cammina, sfalsate di lato & $79{,}2\,\%$ & $\mathbf{73{,}8\,\%}$ & 212~cm & 389~cm & $100\,\%$ \\
Entrambe ferme, sfalsate    & $21{,}2\,\%$ & $\mathbf{19{,}5\,\%}$ & 220~cm & 375~cm & $100\,\%$ \\
B cammina, in fila dietro A & $42{,}4\,\%$ & $\mathbf{0{,}3\,\%}$  & 215~cm & 215~cm & $100\,\%$ \\
\bottomrule
\end{tabular}
\end{table}

\begin{figure}[htbp]
\centering
\includegraphics[width=\linewidth]{Due_persone}
\caption{Distanze riportate dai due canali dell'\ldb\ in un trial, con le due persone sfalsate di lato (a sinistra) e in fila sull'asse (a destra).}
\label{fig:due-persone}

\end{figure}

Con A ferma e B in movimento i due canali sono attivi insieme nel 79\,\% dei campioni, con distanze separate da 177~cm, e l'energia in movimento (38,4) è vicina a quella misurata a 4~m nella caratterizzazione della distanza (35,1), quindi il canale in movimento segue B. Con entrambe ferme le due persone si contendono il canale stazionario e vince la più vicina nel 93\,\% dei campioni, anche se B compare nel 19,5\,\% grazie ai suoi micro-movimenti involontari. In fila A nasconde del tutto B, e i due canali riportano la stessa persona (differenza mediana di 3~cm). Il modulo separa quindi due persone solo quando una si muove e l'altra è ferma, perché riporta un bersaglio per canale. La presenza resta al 100\,\% in ogni geometria, quindi il radar dice sempre che c'è qualcuno ma non quante persone ci sono. Per il progetto va bene, perché ogni arredo ha il proprio sensore e un solo occupante da rilevare, ma un sensore singolo non può contare le persone in un'aula.

\paragraph{Il secondo radar.} L'\ldc\ ha un solo canale di bersaglio (Sezione~\ref{sec:c3-2}), quindi non ha dove riportare la seconda persona, e la prova ripetuta con A ferma a 1~m e B a 2~m lo conferma. La presenza è al 100\,\% in tutte le geometrie. In fila B è invisibile come sull'\ldb, e con le persone sfalsate nemmeno le energie per gate le separano. La distanza riportata si aggancia a B quando B cammina (211--214~cm in due trial su tre) e resta su un valore fisso (268~cm) con entrambe ferme. Il secondo radar dice quindi che c'è qualcuno ma non quante persone, anche quando una delle due si muove.

\section{Penetrazione degli ostacoli}
\label{sec:ostacoli}

Il pannello è interposto a 20~cm dal sensore, perpendicolare all'asse e in piedi da solo, senza supporti. Il soggetto cammina sul posto, perché da fermo il canale stazionario satura e l'attenuazione non sarebbe misurabile. Le distanze sono due. L'\ldb\ lavora a 3~m, dove la sua energia sta circa a metà scala, lontana sia dalla saturazione sia dal fondo. Il PIR lavora a 1~m, perché a 3~m già senza ostacolo non vede il movimento sul posto. L'\ldc\ ha ripetuto la serie a 1~m, perché oltre i 2~m l'esemplare non rileva (Sezione~\ref{sec:portata-massima}). Per l'\ldc\ l'attenuazione è in decibel, $10\log_{10}(E_0/E)$ sull'energia del gate della persona, con $E_0$ senza ostacolo ed $E$ con l'ostacolo, la stessa conversione usata dal tool del produttore (Sezione~\ref{sec:c3-2-5}). Il produttore non documenta però che quelle energie siano potenze, quindi il valore in decibel resta un'ipotesi. Il riferimento senza ostacolo è la media di sei trial a inizio e fine sessione, e i risultati sono in Tabella~\ref{tab:ostacoli} e Figura~\ref{fig:attenuazione}.

\subsection{Attenuazione per materiale}
\label{sec:ostacoli-materiali}
\label{sec:ld2420-attenuazione}

\begin{table}[htbp]
\centering
\caption{Attenuazione per materiale. Per l'\ldb\ è la riduzione dell'energia in movimento rispetto al riferimento senza ostacolo (53,45), in unità dello strumento, per il PIR la riduzione del tempo con presenza rilevata. Il cartongesso di questi due sensori viene da una sessione successiva con la stessa geometria ed è calcolato sul riferimento di quella sessione (Sezione~\ref{sec:ostacoli-portata}). Per l'\ldc, con legno e metallo l'energia scende al fondo (11--13), al limite della dinamica dell'esemplare (circa 5~dB), e il valore è un minimo. Le coperture angolari diverse fanno sì che anche le altre attenuazioni possano essere sottostimate, e il confronto più pulito è legno contro cartone, che hanno la stessa copertura.}
\label{tab:ostacoli}

\small
\setlength{\tabcolsep}{4pt}
\begin{tabular}{@{}lllrrr@{}}
\toprule
\textbf{Materiale} & \textbf{Spessore} & \textbf{Copertura} & \textbf{LD2410B} & \textbf{LD2420} & \textbf{PIR} \\
 & & & \textbf{(3~m)} & \textbf{(1~m)} & \textbf{(1~m)} \\
\midrule
Nessuno (riferimento) & --- & --- & --- & --- & --- \\
Plastica (tanica cava) & $2 \times 5$~mm + aria & $\pm 42$\textdegree & $-7{,}5\,\%$ & 0,4~dB & $-100\,\%$ \\
Cartongesso & 10~mm & $\pm 63$\textdegree & $-26\,\%$ & \textbf{0,7~dB} & $-100\,\%$ \\
Cartone (scatola) & $2 \times 5$~mm + aria & $\pm 54$\textdegree & $-17{,}2\,\%$ & 1,0~dB & $-100\,\%$ \\
Vetroresina & 1~mm & $\pm 70$\textdegree & $-19{,}6\,\%$ & 1,0~dB & $-100\,\%$ \\
Vetro & 5~mm & $\pm 61$\textdegree & $-40{,}9\,\%$ & 2,4~dB & $-95\,\%$ \\
Legno & 10~mm & $\pm 54$\textdegree & $-41{,}6\,\%$ & $\geq 4{,}1$~dB & $-95\,\%$ \\
Metallo & 1~mm & $\pm 50$\textdegree & \textbf{blocca} & $\geq 5{,}5$~dB & $-100\,\%$ \\
\bottomrule
\end{tabular}
\end{table}

\begin{figure}[htbp]
\centering
\includegraphics[width=\linewidth]{Attenuazione_materiali}
\caption{Attenuazione per materiale sui due radar, con i valori della Tabella~\ref{tab:ostacoli}. A: \ldb\ a 3~m, in unità dello strumento. B: \ldc\ a 1~m, in decibel.}
\label{fig:attenuazione}

\end{figure}

Tutti i materiali dielettrici attenuano l'energia dell'\ldb\ fra il 7 e il 42\,\%, ma la presenza resta al 100\,\%. Il PIR invece, che a 1~m senza ostacolo rileva la persona fra il 77 e il 100\,\% del tempo, si azzera con ogni materiale. Il vetro è il caso più evidente, perché è trasparente alla luce ma opaco per il PIR. Il vetro comune non trasmette oltre i 4--5~\textmu m e il corpo emette intorno ai 10~\textmu m, mentre il radar lo attraversa. I due radar ordinano allo stesso modo i sei materiali della serie comune (plastica, cartone e vetroresina, vetro, legno, metallo), pur a distanze e in unità diverse, ed è questa graduatoria che si confronta, non i valori (Sezione~\ref{sec:c3-5-1}). Il cartongesso, materiale comune delle pareti interne, sull'\ldc\ attenua quanto plastica e cartone, sull'\ldb\ un po' di più, fra vetroresina e vetro, ma quel valore viene da un'altra sessione. In entrambi i casi resta ben sotto vetro e legno. Il metallo invece blocca del tutto l'\ldb, che riporta un bersaglio fisso a 30~cm, cioè il pannello, e non vede più il soggetto, quindi il segnale non aggira il pannello. Sull'\ldc\ il legno porta il gate~2 quasi sempre sotto la soglia di mantenimento (0--1\,\% dei campioni sopra, contro il 25--33\,\% senza ostacolo), e quindi anche sotto quella di attivazione. L'esemplare non acquisirebbe una persona a 1~m dietro un pannello di legno, che l'\ldb\ rileva al 100\,\% a 3~m.

\subsection{Portata residua con l'ostacolo}
\label{sec:ostacoli-portata}

L'\ldb\ ha ripetuto la geometria allontanando il soggetto da 1 a 5~m dietro il cartongesso e da 3 a 5~m dietro legno e vetro, con tre trial per punto (Figura~\ref{fig:portata-ostacoli}). Una prova a parte usa una porta interna chiusa, tamburata e senza vetro, a circa 3,8~m dal sensore, con il soggetto a 3, 4 e 5~m. Nessun ostacolo accorcia la portata entro i 5~m della stanza, e in tutti i trial con ostacolo la presenza è al 100\,\% e la distanza è corretta. A degradare è il canale in movimento. A 5~m la frazione di campioni in cui è attivo passa dal 78\,\% senza ostacolo al 45\,\% con il cartongesso, al 26\,\% con il vetro, al 19\,\% con il legno e al 3\,\% oltre la porta, mentre la presenza la tiene il canale stazionario. Con il soggetto a 3~m, ancora davanti alla porta, l'energia è pari al riferimento, e appena la porta si interpone, a 4~m, il canale in movimento scende al 20\,\%. Anche qui il PIR è a zero dietro ogni ostacolo, porta compresa.

\begin{figure}[htbp]
\centering
\includegraphics[width=\linewidth]{Portata_ostacoli}
\caption{\ldb, portata residua con i pannelli a 20~cm dal sensore e la porta a circa 3,8~m, tre trial per punto. A: energia in movimento. B: frazione di campioni con il canale in movimento attivo.}
\label{fig:portata-ostacoli}

\end{figure}

Per il progetto questo significa che il sensore può essere incassato nell'arredo,
dietro legno o cartongesso o dentro un guscio di vetroresina, mentre il PIR deve affacciarsi direttamente sull'ambiente. Il metallo resta l'unico vincolo assoluto, e nel banco
reale la lamiera di rinforzo sta dietro al modulo e non davanti.

\section{Portata massima in corridoio}
\label{sec:portata-massima}

La prova di portata si svolge in un corridoio largo 120~cm (Figura~\ref{fig:geometrie}d), con i sensori a 80~cm da terra e due vani porta a 4 e 5~m, attraverso i quali sono visti i punti a 6, 7 e 8~m. In venti minuti di corridoio vuoto l'\ldb\ e il PIR non danno alcun falso positivo, e i gate lontani dell'\ldb, dove cadono vani e muro di fondo, restano sempre sotto metà soglia (al massimo 7 contro 15). I risultati dei tre sensori sono in Figura~\ref{fig:portata-tre-sensori}.

\begin{figure}[htbp]
\centering
\includegraphics[width=\linewidth]{Portata_tre_sensori}
\caption{Portata dei tre sensori sulla stessa scala, con un punto per metro. L'\ldb\ e il PIR sono misurati in stanza fino a 5~m e in corridoio oltre, l'\ldc\ tutto in corridoio. Il PIR ha due serie, a metà corsa (la sensibilità della campagna) e al massimo. Per l'\ldc\ i triangoli pieni indicano l'acquisizione della presenza, quelli vuoti il solo mantenimento.}
\label{fig:portata-tre-sensori}

\end{figure}

\paragraph{\ldb.} Con un'oscillazione trasversale sul posto di circa 50~cm il radar
rileva al 100\,\% a 5 e 6~m e allo 0\,\% a 7~m in tre trial su tre. A 6~m la distanza
riportata ha massimo esattamente 600~cm in cinque trial su cinque, con dispersione di
2--5~cm contro i 10--11 di 5~m. Non si tratta di una misura più precisa, ma del tetto
degli otto gate da 75~cm configurati. La presenza a 6~m è retta dal canale stazionario
(energia 62--82), mentre a 7~m i gate 7 e 8 restano molto sotto le soglie. Il tetto di 600~cm dichiarato dal manuale~\cite[§5.2]{hilinkManualV104} risulta così misurato.

\paragraph{PIR.} Con il trimmer a metà corsa, la sensibilità usata in tutta la campagna, l'oscillazione trasversale è rilevata al 24--33\,\% anche a 7~m, dove il radar per configurazione non vede più nulla. L'attraversamento a circa 1~m/s, rilevato al 98,7\,\% a 5~m, scende invece a 0, 6 e 0\,\% nei tre trial a 6~m, e sale al 30\,\% se il transito è più veloce. Al limite della portata conta quindi anche la velocità del transito, oltre al tipo e all'ampiezza del movimento. Con il trimmer al massimo, in una serie a parte rispetto al resto della campagna, l'attraversamento è rilevato al 100, 100, $56{,}8 \pm 18{,}0$ e $14{,}7 \pm 10{,}0\,\%$ a 5, 6, 7 e 8~m. Il limite sta quindi fra 7 e 8~m, coerente con i 3--7~m del datasheet~\cite{hcsr501}, e supera il tetto del radar, ma solo per chi attraversa il campo.

\paragraph{\ldc.} Con il gate massimo portato a 12 (840~cm), il valore di fabbrica
che copre gli 8~m dichiarati~\cite{hilinkManualLD2420}, l'esemplare acquisisce a 1~m e mantiene a 2~m. Da 3~m in
su le sue energie sono il fondo del corridoio vuoto gate per gate. Il gate massimo non
ha quindi aggiunto nulla, e il limite è dell'esemplare e non della configurazione né
della stanza.

La prova serve a valutare un'alternativa all'architettura del progetto, cioè pochi sensori a parete che guardano l'intera stanza al posto di uno sotto ogni banco. Con 6~m di portata l'\ldb\ coprirebbe una stanza come quella di prova, ma un sensore a parete dice solo se nella stanza c'è qualcuno, non sotto quale banco né quante persone, perché chi sta dietro un altro corpo non viene separato (Sezione~\ref{sec:due-persone}). Per lo scenario del progetto resta quindi un sensore per arredo.

\chapter{L'indice di vitalità}
\label{chap:vitalita}

Questo capitolo presenta l'indice di vitalità, un numero da 0 a 100 che indica quanto si muove la persona rilevata dal radar. L'indice è costruito e tarato sui dati del Capitolo~\ref{chap:confronto}, calcolato a bordo dell'ESP32 e verificato contro lo stesso calcolo eseguito al computer.

Il capitolo è organizzato come segue. La Sezione~\ref{sec:vitalita-motivazione} spiega perché la vitalità è una domanda diversa dalla presenza e come l'indice la rende leggibile. La Sezione~\ref{sec:vitalita-ingressi} descrive i dati di ingresso e l'algoritmo, insieme alle due versioni precedenti abbandonate, e la Sezione~\ref{sec:vitalita-classi} definisce le tre classi. La Sezione~\ref{sec:vitalita-taratura} riporta la taratura dei parametri e la validazione su trial separati e su una geometria diversa. La Sezione~\ref{sec:vitalita-bordo} descrive il porting sull'ESP32 e la verifica del calcolo a bordo contro quello offline, e la Sezione~\ref{sec:vitalita-limiti} ne raccoglie i limiti, i casi limite e le possibili estensioni. La Sezione~\ref{sec:respiro} chiude con la stima del ritmo respiratorio e la sua verifica a ritmo imposto.

\section{Motivazione e presentazione}
\label{sec:vitalita-motivazione}
\label{sec:vitalita-presentazione}

Per i soccorritori contano due domande. La prima è se sotto il banco c'è qualcuno, e a questa risponde la presenza. La seconda è in che condizioni si trova la persona, e a questa risponde la vitalità. Il PIR non può rispondere alla seconda domanda, perché la sua uscita dice solo se c'è stato un movimento (Sezione~\ref{sec:interpretazione-pir}). Il radar invece fornisce l'energia riflessa per gate, che cresce con il movimento. Nella dose-risposta a 1~m vale 68,6 da immobile, 84,4 con i micro-movimenti e 99,3 in cammino (Sezione~\ref{sec:dose-risposta}). L'indice di vitalità riassume questa grandezza in un valore leggibile anche da chi non conosce il radar, con tre classi, vitalità bassa, moderata e alta.

\section{Dati di ingresso e algoritmo}
\label{sec:vitalita-ingressi}
\label{sec:vitalita-algoritmo}

L'algoritmo lavora sui campioni a 5~Hz prodotti dal logger e dal firmware della web UI (Sezione~\ref{sec:c3-4}). Sia $e_i[k]$ l'energia del canale in movimento nel gate $i$ al campione $k$, $f_i$ il rumore di fondo dello stesso gate, misurato a stanza vuota, e $d_k$ la distanza riportata, quella del bersaglio stazionario se disponibile, altrimenti quella del bersaglio in movimento. L'algoritmo sceglie il gate dalla distanza, ne calcola l'energia al netto del fondo e la segue con due medie esponenziali, che poi somma nell'indice.

\begin{align}
g^{*} &= \min\!\left(\left\lfloor d_k / 75~\text{cm} \right\rfloor,\; 8\right)
  && \text{gate attivo} \label{eq:gate-attivo}\\
\tilde e_k &= \max\!\left(0,\; \frac{e_{g^{*}}[k] - f_{g^{*}}}{100 - f_{g^{*}}} \cdot 100\right)
  && \text{energia netta, riscalata su 0--100} \label{eq:netta}\\
M_k &= \alpha_m\, \tilde e_k + (1-\alpha_m)\, M_{k-1}
  && \text{livello di movimento, } \alpha_m = 0{,}05 \label{eq:livello}\\
R_k &= \alpha_v\, \left|\tilde e_k - \tilde e_{k-1}\right| + (1-\alpha_v)\, R_{k-1}
  && \text{variabilità del ritorno, } \alpha_v = 0{,}01 \label{eq:variabilita}\\
V_k &= \mathrm{clamp}\!\left(M_k + \beta\, R_k,\; 0,\; 100\right), \quad \beta = 0{,}5
  && \text{indice di vitalità} \label{eq:vitalita}
\end{align}

Se il radar non riporta presenza l'indice è 0 e nessuna classe viene assegnata. Il gate attivo viene dalla distanza riportata e non dal massimo di energia, perché sotto il banco i cammini multipli portano energia nei gate 2--3 mentre la persona è a 60~cm (Sezione~\ref{sec:sotto-banco}). Il fondo va sottratto perché anche a stanza vuota l'energia non è zero, soprattutto nei gate più vicini al sensore (18 nel gate 0 e 13 nel gate 1, Tabella~\ref{tab:vitalita-gate}), e senza sottrazione questo rumore entrerebbe nell'indice come se fosse movimento. Calcolando l'indice anche quando il radar non riporta presenza, la stanza vuota darebbe 19,1 e la persona immobile a 2,3~m 17,6, cioè il solo rumore varrebbe più della persona. Con la sottrazione i due casi diventano 2,4 e 15,1.

Le due medie esponenziali seguono grandezze diverse. Il livello di movimento $M_k$ (Equazione~\ref{eq:livello}) è la media dell'energia netta, e dice quanto segnale torna dalla persona. La variabilità $R_k$ (Equazione~\ref{eq:variabilita}) è la media di quanto l'energia netta cambia da un campione all'altro, e dice quanto il ritorno oscilla. Il livello deve reagire in pochi secondi, e con $\alpha_m = 0{,}05$ a 5~Hz la sua costante di tempo è di circa 4~s. La variabilità è un segnale piccolo e rumoroso, e con $\alpha_v = 0{,}01$ viene mediata su circa 20~s. Si usano medie esponenziali perché ognuna si aggiorna dal proprio valore precedente, senza conservare i campioni passati. Il peso $\beta$ fissa quanto la variabilità conta nell'indice, e con $\beta = 0$ l'indice si ridurrebbe al solo livello. Con $\beta = 0{,}5$ i valori salgono di qualche punto, di più per i micro-movimenti a 1~m (da circa 59 a circa 71) e di meno per la persona immobile e per il cammino, quindi le tre condizioni risultano un po' più distanziate. La classe assegnata invece cambia poco, come mostra la taratura (Sezione~\ref{sec:vitalita-taratura}).

\subsection{Evoluzione dell'algoritmo}
\label{sec:vitalita-evoluzione}
L'algoritmo è arrivato alla forma attuale in tre versioni.

\begin{itemize}
    \item \textbf{v1.} Calcolava la variabilità sull'energia stazionaria, che però nei dati resta ferma a 100 quando c'è qualcuno e vale zero sotto 1,5~m (Sezione~\ref{sec:c3-5-2}). La variabilità era quindi sempre nulla.

    \item \textbf{v2.} Usava i canali in movimento, ma prendeva come livello il massimo fra l'energia del gate attivo e l'energia complessiva del bersaglio riportata dal radar. A 1~m funzionava, ma sotto il banco l'energia complessiva è quasi sempre satura a 100, quindi il massimo sceglieva quasi sempre quella al posto dell'energia del gate, e la persona immobile risultava più attiva di quella a 1~m.

    \item \textbf{v3.} È la versione descritta sopra, che usa solo l'energia del gate attivo. Lo scenario sotto il banco rientra in scala senza che quelli a 1~m cambino in modo sensibile (Tabella~\ref{tab:vitalita-v2v3}).
\end{itemize}

\begin{table}[htbp]
\centering
\caption{Indice medio per scenario nelle versioni v2 e v3, con la configurazione della taratura (gate fisso~2, fondo della stanza vuota di giorno, parametri tarati). La saturazione è la frazione di campioni con indice a 100.}
\label{tab:vitalita-v2v3}
\small
\begin{tabular}{@{}lrrc@{}}
\toprule
Scenario & v2 & v3 & Saturazione v2 $\to$ v3 \\
\midrule
Immobile, 1~m & 31,1 & 28,6 & 0 $\to$ 0\,\% \\
Micro-movimenti, 1~m & 77,3 & 71,0 & 4 $\to$ 0\,\% \\
Movimento pieno, 1~m & 99,6 & 99,2 & 92 $\to$ 71\,\% \\
\textbf{Immobile sotto il banco} & \textbf{51,1} & \textbf{25,7} & 1 $\to$ 0\,\% \\
\textbf{Movimenti sotto il banco} & 100,0 & \textbf{98,8} & \textbf{100 $\to$ 46\,\%} \\
\bottomrule
\end{tabular}
\end{table}

\section{Classificazione}
\label{sec:vitalita-classi}

L'indice di vitalità è diviso in tre classi, con le soglie a 45 e 95 fissate nella taratura a 1~m (Tabella~\ref{tab:vitalita-classi}). Sotto il banco la prima soglia sale a 70 (Sezione~\ref{sec:vitalita-taratura}). Chi si muove ampiamente è verosimilmente cosciente, chi compie solo micro-movimenti è vivo ma forse ferito o bloccato, e chi non si muove può essere incosciente. La priorità di soccorso va quindi al contrario dell'indice. La vitalità bassa è il caso da raggiungere per primo, mentre chi ha vitalità alta è probabilmente in grado di aspettare. Per questo nella dashboard (Sezione~\ref{sec:webui-progetto}) la classe bassa è rossa e quella alta verde.

\begin{table}[htbp]
\centering
\caption{Classi dell'indice di vitalità con le soglie tarate
(Sezione~\ref{sec:vitalita-taratura}) e loro significato operativo. }
\label{tab:vitalita-classi}
\small
\begin{tabularx}{\linewidth}{@{}llXl@{}}
\toprule
\textbf{Indice} & \textbf{Classe} & \textbf{Significato operativo} & \textbf{Priorità} \\
\midrule
--- & (nessuna presenza) & il sensore non rileva nessuno & --- \\
0--44 & vitalità bassa & presenza confermata, movimento minimo, possibile
  incoscienza & \textbf{la più alta} \\
45--94 & vitalità moderata & micro-movimenti, persona che si aggiusta & media \\
95--100 & vitalità alta & movimento ampio, persona reattiva & la più bassa \\
\bottomrule
\end{tabularx}
\end{table}

Le classi sono un supporto al triage e non una diagnosi medica. Vitalità bassa significa che il sensore rileva una presenza con variazioni minime, non che la persona sia in pericolo di vita. Se il radar non rileva nessuno l'indice non viene classificato, e l'assenza di presenza non è una quarta classe. Significa che il sensore non vede nessuno, non che sotto il banco non ci sia nessuno, perché la persona può essere fuori portata o schermata da metallo.

\section{Taratura e validazione}
\label{sec:vitalita-taratura}

L'algoritmo è stato tarato al computer sui file CSV e portato sull'ESP32 solo a parametri validati. I dati sono quelli della dose-risposta a 1~m (Sezione~\ref{sec:dose-risposta}), con cinque trial per ciascuno dei tre scenari (immobile, micro-movimenti, cammino sul posto), e delle due serie sotto il banco (Sezione~\ref{sec:sotto-banco}), usate come validazione su geometria diversa. I parametri ($\alpha_m$, $\alpha_v$, $\beta$ e le due soglie) sono stati scelti sui trial T01--T03 dei tre scenari a 1~m. La prima soglia separa la persona immobile dai micro-movimenti, la seconda i micro-movimenti dal cammino. La verifica è fatta sui trial T04--T05, che non sono serviti a scegliere i parametri e misurano quindi come si comporta l'indice su dati nuovi. Per ogni scenario si conta la percentuale di campioni che cadono nella classe attesa, e la media di queste percentuali è la recall media riportata nelle tabelle. Il transitorio iniziale è scartato come nel Capitolo~\ref{chap:confronto}. Nella taratura il gate attivo era fissato al gate~2 anziché ricavato dalla distanza, perché a 1~m la persona sta a cavallo dei gate 1 e 2 e il gate~2 ha il fondo più pulito. Senza la variabilità, e con le soglie rifissate, la recall media su tutti i trial a 1~m passa dal 92,2 al 91,6\,\%.

\begin{table}[htbp]
\centering
\caption{Prestazione dell'indice a tre classi con i parametri tarati e gate fisso 2.}
\label{tab:vitalita-validazione}
\small
\begin{tabular}{@{}lc@{}}
\toprule
Insieme & Recall media \\
\midrule
Taratura (T01--T03, 1~m) & $95{,}1\,\%$ \\
Validazione (T04--T05, 1~m) & $\mathbf{88{,}0\,\%}$ \\
Validazione su geometria diversa (sotto il banco) & $51{,}2\,\%$ \\
\midrule
Discriminazione immobile / in movimento, entrambe le geometrie & $\mathbf{96{,}1\,\%}$ \\
\bottomrule
\end{tabular}
\end{table}

\begin{table}[htbp]
\centering
\caption{Dettaglio per scenario con i parametri tarati, tutti i trial, gate fisso 2 e fondo della stanza vuota di giorno.
L'indice è la media dell'indice medio per trial. }
\label{tab:vitalita-scenari}
\small
\begin{tabular}{@{}lrlr@{}}
\toprule
Scenario & Indice & Classe attesa & Corretti \\
\midrule
Immobile a 1~m & 28,6 & bassa & $89{,}9\,\%$ \\
Micro-movimenti a 1~m & 71,0 & moderata & $90{,}5\,\%$ \\
Movimento pieno a 1~m & 99,2 & alta & $96{,}3\,\%$ \\
Immobile sotto il banco & 25,7 & bassa & $97{,}5\,\%$ \\
Movimenti sotto il banco & 98,8 & moderata & $4{,}8\,\%$ \\
\bottomrule
\end{tabular}
\end{table}

\begin{figure}[htbp]
\centering
\includegraphics[width=\linewidth]{Vitalita_scenari}
\caption{Distribuzione dell'indice di vitalità per scenario, con la configurazione del firmware (gate dalla distanza riportata, fondo notturno), transitorio scartato. Le fasce colorate sono le tre classi. La sesta scatola ripete la persona immobile sotto il banco con il gate fissato a 2, come nella taratura.}
\label{fig:vitalita-scenari}
\end{figure}

\paragraph{Sulla geometria di taratura.} Sui trial di validazione a 1~m la recall media è dell'88\,\% (Tabelle~\ref{tab:vitalita-validazione} e~\ref{tab:vitalita-scenari}). Le tre distribuzioni
si separano (Figura~\ref{fig:vitalita-scenari}), con mediane 23, 77 e 100. La
sovrapposizione è concentrata fra la persona immobile e quella che compie piccoli
movimenti.

\paragraph{Sulla geometria diversa.} Il 51,2\,\% sotto il banco è la media di due scenari molto diversi, la persona immobile riconosciuta nel 97,5\,\% dei campioni e i movimenti nel 4,8\,\%. Non è però una prova che le soglie si trasferiscano alla geometria del progetto. La ragione sta nel gate fissato a~2. Se a 1~m il gate~2 vede bene la persona, sotto il banco l'energia si concentra nel gate~1 (54 contro 21,7), e il gate~2 ne raccoglie soltanto un'eco indiretta (Tabella~\ref{tab:vitalita-gate}). Il 25,7 viene quindi da un gate che vede la persona solo di riflesso, e coincide con il valore dell'immobile a 1~m per caso. Nel gate~1 la stessa persona immobile ha energia 54, più vicina ai micro-movimenti a 1~m (66,5) che all'immobile (30,9), e con la configurazione del firmware (gate dalla distanza, fondo notturno) l'indice vale 45,4, in classe bassa solo nel 49\,\% dei campioni. Sotto il banco resta invece valida la distinzione fra immobile e in movimento, che è quella che conta per il triage. Su tutti e cinque gli scenari e su entrambe le geometrie ha una recall media del 96,1\,\%, con qualunque criterio di gate. L'unico scenario che fallisce sono i movimenti sotto il banco, classificati alta invece di moderata. Quella classe attesa però era stata assegnata per analogia con il caso a 1~m, e a 60~cm chi si aggiusta produce un ritorno molto forte, quindi non si può dire che l'indice sbagli.

\begin{table}[htbp]
\centering
\caption{Energia in movimento media grezza per gate nei cinque scenari.}
\label{tab:vitalita-gate}
\small
\begin{tabular}{@{}lrrrr@{}}
\toprule
Scenario & Gate 0 & Gate 1 & Gate 2 & Gate 3 \\
\midrule
Stanza vuota (fondo) & 17,6 & 13,2 & 4,4 & 2,9 \\
Immobile a 1~m & 18,6 & 30,9 & 25,3 & 7,4 \\
Micro-movimenti a 1~m & 23,6 & 66,5 & 60,5 & 14,3 \\
\textbf{Immobile sotto il banco (60~cm)} & 33,1 & \textbf{54,0} & 21,7 & 9,5 \\
Movimenti sotto il banco & 96,5 & 99,8 & 95,1 & 43,9 \\
\bottomrule
\end{tabular}
\end{table}

\paragraph{Taratura nella geometria d'uso.} Le soglie tarate a 1~m non vanno quindi bene sotto il banco. Con la configurazione del firmware la persona immobile sta a cavallo della prima soglia (Figura~\ref{fig:vitalita-scenari}), e il gate fisso non risolve il problema ma lo nasconde. La taratura è stata perciò ripetuta con lo stesso metodo sui dati sotto il banco, scegliendo la soglia sui trial T01--T03 e verificandola sui T04--T05. La prima soglia sale da 45 a 70. Sui trial di verifica la persona immobile cade così in classe bassa nel 96,9\,\% dei campioni, contro il 75,1\,\% con la soglia tarata a 1~m, e nessun campione con movimento finisce in quella classe. Le soglie vanno dunque tarate nella geometria in cui il sensore è montato. La Sezione~\ref{sec:vitalita-limiti} propone di renderla automatica.

\section{Porting a bordo e verifica}
\label{sec:vitalita-bordo}

L'indice è calcolato a bordo del nodo e mostrato nella dashboard (Sezione~\ref{sec:webui}). Le Equazioni~\ref{eq:gate-attivo}--\ref{eq:vitalita} sono implementate in 74 righe di C++ senza dipendenze, con le stesse costanti del prototipo raccolte in un file di configurazione, e costano poche operazioni in virgola mobile per campione. A differenza della taratura, il firmware ricava il gate dalla distanza riportata e usa il fondo della stanza vuota notturna di 6,5~ore. A 1~m i risultati coincidono quasi con quelli della taratura (28,4, 73,9 e 99,4 contro 28,6, 71,0 e 99,2), e cambiare il fondo sposta l'indice di meno di un punto.

Il porting è stato verificato confrontando l'indice calcolato sull'ESP32 con quello ricalcolato in Python sullo stesso CSV, con il criterio che coincidessero entro $\pm 1$. Il firmware scrive indice e classe in due colonne aggiuntive del CSV esportato dal browser (Sezione~\ref{sec:webui-progetto}), e uno script di confronto rilegge il file con il prototipo campione per campione. Le sessioni sono sei, da 90~s, con soggetto immobile, con micro-movimenti e in movimento, a circa 1,9~m e a 1~m. A regime i due indici coincidono entro $\pm 1$ nel 98,7--100\,\% dei campioni (Figura~\ref{fig:vitalita-bordo}). Nei primi 60~s circa divergono, perché le medie esponenziali del firmware sono aggiornate dall'accensione del nodo, mentre quelle del prototipo partono dal primo campione del file. La differenza dipende quindi solo dal punto di partenza, non dal calcolo.

\begin{figure}[htbp]
\centering
\includegraphics[width=\linewidth]{Vitalita_bordo}
\caption{Verifica del porting, con l'indice calcolato a bordo dell'ESP32 e il ricalcolo offline sullo stesso CSV esportato dalla web UI, persona immobile a 1~m.}
\label{fig:vitalita-bordo}
\end{figure}

Queste sessioni mostrano il funzionamento a bordo, ma non sono una seconda validazione delle soglie (Tabella~\ref{tab:vitalita-live}). Tre valori cadono fuori dalla classe attesa. A 1,9~m il movimento resta appena sotto la soglia 95 tarata a 1~m, e a 1~m la persona immobile cade appena sopra la soglia 45. Sempre a 1~m i micro-movimenti sono classificati alta, ma in quella sessione l'energia del gate~1 era 97,9, contro 66--68 nei micro-movimenti della taratura. Lo stimolo era quindi un movimento pieno, e l'indice l'ha letto correttamente. È la condizione ``micro-movimenti'' a non riprodursi fra sessioni, perché dipende da quanto si muove la persona.

\begin{table}[htbp]
\centering
\caption{Indice medio nelle sessioni dal vivo, con la classe assegnata.}
\label{tab:vitalita-live}
\small
\begin{tabular}{@{}lccc@{}}
\toprule
 & Immobile & Micro-movimenti & Movimento \\
\midrule
1,9~m & 19,9 (bassa) & 68,7 (moderata) & 94,0 (moderata) \\
1~m & 47,7 (moderata) & 99,4 (alta) & 99,1 (alta) \\
\bottomrule
\end{tabular}
\end{table}

\section{Limiti, casi limite ed estensioni}
\label{sec:vitalita-limiti}

L'indice non si può calcolare con l'\ldc. Sull'esemplare in dotazione l'energia non cresce con il movimento (Sezione~\ref{sec:dose-risposta}), e manca quindi la grandezza continua su cui l'indice si basa. L'indice è definito solo sull'\ldb, e il limite va attribuito all'esemplare, non al modello.

I comportamenti dell'indice nei casi limite
sono raccolti in Tabella~\ref{tab:vitalita-casilimite}.

\begin{table}[htbp]
\centering
\caption{Casi limite dell'indice di vitalità e comportamento previsto o osservato.}
\label{tab:vitalita-casilimite}
\small
\begin{tabularx}{\linewidth}{@{}XXX@{}}
\toprule
\textbf{Situazione} & \textbf{Comportamento} & \textbf{Note} \\
\midrule
La persona esce dal campo & indice a 0 e nessuna classe entro la coda di presenza
  del radar (circa 9--12~s stimati, Sezione~\ref{sec:latenza}) & l'indice dipende dalla presenza. Le medie continuano ad aggiornarsi \\
La persona si addormenta & da alta a bassa in pochi secondi & il livello decade con la sua costante di circa 4~s \\
Ventilatore o tenda in movimento & possibile falsa classe moderata & mitigabile con le soglie per gate del radar \\
Due persone & indice riferito al gate della distanza riportata & coerente con il
  limite di un bersaglio per canale (Sezione~\ref{sec:due-persone}) \\
Persona oltre 4--5~m & energia debole, indice sottostimato & la portata utile
  dell'indice non è stata caratterizzata \\
Persona a meno di 1,5~m & nessun effetto, l'indice usa i soli canali in movimento &
  perché i gate stazionari sotto 1,5~m sono nulli (Sezione~\ref{sec:c3-5-2}) \\
\bottomrule
\end{tabularx}
\end{table}

I limiti dell'indice:

\begin{itemize}
    \item le soglie sono specifiche dell'installazione;

    \item la validazione poggia su un solo soggetto e due trial per scenario;

    \item la classe intermedia non ha un riferimento indipendente, perché nessuno strumento ha misurato quanto si muoveva la persona durante i micro-movimenti;

    \item non è stato misurato fino a che distanza l'indice classifica correttamente, perché è tarato a 1~m e sotto il banco.
\end{itemize}

Tre estensioni che migliorerebbero l'indice:

\begin{itemize}
    \item \textbf{Auto-taratura all'installazione.} Registrerebbe il fondo a banco vuoto, la persona immobile e la persona in movimento sotto l'arredo, e ne ricaverebbe le soglie come fatto a mano nella Sezione~\ref{sec:vitalita-taratura}.

    \item \textbf{Classi più stabili.} La classe cambierebbe solo dopo qualche secondo oltre la soglia, così che l'etichetta non sfarfalli quando l'indice oscilla intorno a una soglia, come nella sessione a 1,9~m.

    \item \textbf{Indicatore di confidenza.} Affiancato all'indice e calcolato dalla stabilità del gate attivo, direbbe quanto fidarsi del valore.
\end{itemize}

Resta infine un aspetto di trattamento dei dati. Per il monitoraggio in ambienti scolastici il progetto indica le cautele del GDPR, fino a una valutazione d'impatto quando sono coinvolti minori~\cite{callisto2026dipme}. L'indice di vitalità non aggiunge un problema nuovo, perché non è una misura fisiologica, e rientra nella stessa valutazione che riguarda già il dato di presenza.

\section{Il rilevamento del respiro}
\label{sec:respiro}

Il respiro è il micro-movimento che permette al radar di rilevare anche la persona immobile, e l'energia per gate permette di stimarne il ritmo. Il principio è quello della Sezione~\ref{subsec:corpo-immobile}. Si campiona
a 5~Hz un canale di energia per gate in engineering mode, si applica la FFT alla serie
temporale e si cerca il picco in banda 0,1--0,5~Hz. La frequenza del picco per 60 dà
gli atti al minuto. L'\ldb\ espone venti serie di energia e non tutte sono
utilizzabili. Alcune sono sature (Sezione~\ref{sec:c3-5-1}), altre nulle
(Sezione~\ref{sec:c3-5-2}), altre piatte perché lontane dal bersaglio. Lo script di
analisi esamina tutti i canali, scarta saturi e piatti con criteri dichiarati, ordina
i rimanenti per rapporto segnale-rumore e riporta la mediana delle stime concordi~\cite{repoTesi}.

\subsection{Prove senza ground truth}
\label{sec:respiro-senza-gt}

Su un soggetto fermo a 40~cm sette canali in
movimento indipendenti concordano su 0,264~Hz, cioè 15,8 atti al minuto, con rapporto
segnale-rumore fino a 12,8. Sui cinque trial del soggetto seduto a 2,30~m
(Sezione~\ref{sec:persona-ferma}) i canali in movimento dei gate 2--8 concordano su 18--20
atti al minuto, e a quella distanza non sono saturi (Figura~\ref{fig:respiro}). Sotto il
banco (Sezione~\ref{sec:sotto-banco}) il respiro è estraibile in quattro trial su cinque, con
$21{,}3 \pm 2{,}3$ atti al minuto.
In tutte e tre le geometrie i canali stazionari danno 6--9 atti al minuto, troppo
pochi per un adulto a riposo, e sono verosimilmente un artefatto del filtraggio
interno. Queste sono stime concordi fra canali e non misure con un errore, perché il
ritmo reale non è noto.

\begin{figure}[htbp]
\centering
\includegraphics[width=\linewidth]{Respiro}
\caption{Energia per gate di un soggetto immobile a 2,3~m e suo spettro.}
\label{fig:respiro}
\end{figure}

\subsection{La verifica a ritmo imposto}
\label{sec:respiro-metronomo}

Per verificare la stima serve conoscere il ritmo reale, e il modo più semplice è imporlo.
Il soggetto è seduto a 2~m rivolto al sensore, perché a 1~m i canali in movimento
saturano, e respira seguendo un metronomo impostato al doppio della frequenza respiratoria, un colpo per l'inspirazione e uno per l'espirazione. Sono stati provati tre ritmi, con cinque trial ciascuno da 162~s utili,
più tre trial di stanza vuota come controllo negativo
(Tabella~\ref{tab:respiro-metronomo}).

Quando l'analisi trova il ritmo giusto la stima è accurata, con errore medio di 0,22 atti al minuto sui tredici trial concordi e massimo 0,4. Nei due trial restanti l'analisi ha restituito esattamente il doppio del ritmo imposto, e il motivo è fisico. L'energia per gate misura quanto si sposta il torace, non in che direzione. In ogni respiro il torace si sposta due volte, in inspirazione e in espirazione, quindi l'energia ha due picchi per respiro e nello spettro compare, oltre al ritmo vero, anche il suo doppio. Quando il picco al doppio è più alto, l'analisi sceglie quello. Dallo spettro di un solo canale non si può distinguere un respiro a 30 atti al minuto da uno a 15 con due picchi per ciclo.

\begin{table}[htbp]
\centering
\caption{Stima respiratoria contro ritmo imposto, \ldb\ a 2~m. ``Trial concordi''
conta quelli in cui l'analisi ha restituito il ritmo imposto e non il suo doppio. }
\label{tab:respiro-metronomo}
\small
\begin{tabular}{@{}lccc@{}}
\toprule
Ritmo imposto & Trial concordi & Stima & Errore \\
\midrule
10 atti/min & 4/5 & $10{,}00 \pm 0{,}00$ & $0{,}00$ \\
15 atti/min & 4/5 & $14{,}95 \pm 0{,}30$ & $-0{,}05$ \\
20 atti/min & 5/5 & $19{,}64 \pm 0{,}09$ & $-0{,}36$ \\
\midrule
Stanza vuota & --- & \multicolumn{2}{c}{\textbf{nessuna stima in 3 trial su 3}} \\
\bottomrule
\end{tabular}
\end{table}

Il confronto con il ritmo imposto ha anche permesso di correggere lo script. Prima scartava stime valide quando canali in movimento e stazionari concordavano, e quando trovava sia il ritmo sia il suo doppio poteva scegliere il doppio. Ora sceglie il più basso dei due. È stata inoltre aggiunta una soglia di plausibilità a 8 atti al minuto, che rende il metodo cieco a una respirazione molto lenta.

Restano due limiti. Il metronomo approfondisce il respiro, quindi la validazione riguarda la frequenza e non l'ampiezza del segnale in condizioni realistiche. Inoltre sull'esemplare di \ldc\ lo stesso metodo recupera la frequenza in un trial su tre, e il respiro resta quindi una misura dell'\ldb.

\chapter{Consumo energetico}
\label{chap:consumi}

Questo capitolo stima il consumo energetico dei sensori e del nodo che li ospita. L'analisi è a titolo informativo e usa i valori dichiarati dai produttori.

Il capitolo è organizzato come segue. La Sezione~\ref{sec:consumi-dati} riprende i consumi dei sensori dal Capitolo~\ref{chap:moduli} e vi aggiunge quelli del microcontrollore. La Sezione~\ref{sec:consumi-nodi} li somma nelle configurazioni del nodo, e la Sezione~\ref{sec:autonomia} ne calcola l'autonomia a batteria. La Sezione~\ref{sec:consumi-vincoli} ne trae i vincoli per DIPME.

\section{Dati di partenza}
\label{sec:consumi-dati}

Le correnti dei sensori sono riportate nelle specifiche del Capitolo~\ref{chap:moduli}, e per i conti di questo capitolo basta la corrente media di ciascuno. L'\ldb\ assorbe in media 80~mA a 5~V, pari a circa 400~mW, e l'\ldc\ 50~mA a 3,3~V, circa 165~mW. L'\pir\ a riposo resta sotto i 50~\textmu A, cioè circa 0,3~mW. Per l'ESP32-WROOM-32 le correnti nelle diverse modalità operative sono riportate nella Tabella~\ref{tab:consumi-esp32}, dal datasheet Espressif~\cite{esp32datasheet}.

\begin{table}[htbp]
\centering
\caption{Consumo dell'ESP32-WROOM-32 per modalità operativa.}
\label{tab:consumi-esp32}
\small
\begin{tabular}{@{}lr@{}}
\toprule
\textbf{Modalità} & \textbf{Corrente} \\
\midrule
Attivo, WiFi in trasmissione (picco) & 160--260~mA \\
Attivo, WiFi connesso in ascolto & 95--100~mA \\
Modem-sleep (CPU attiva a 80--240~MHz, WiFi spento) & 20--68~mA \\
Light-sleep & 0,8~mA \\
Deep-sleep (solo dominio RTC) & 0,01~mA \\
\bottomrule
\end{tabular}
\end{table}

Il nodo del progetto DIPME usa invece un microcontrollore STM32WLE con radio LoRa e un PIR digitale Panasonic EKMB139~\cite{callisto2026dipme}, componenti a bassissimo consumo pensati per la batteria. I valori di questo capitolo si riferiscono al kit degli esperimenti, ma anche con componenti diversi è il radar a determinare il consumo finale del nodo.

\section{Il nodo completo}
\label{sec:consumi-nodi}

Sommando microcontrollore e sensore si ottiene il consumo del nodo completo (Tabella~\ref{tab:consumi-nodi}). Le prime tre righe sono le configurazioni usate negli esperimenti, le ultime due quelle con il solo PIR, con il microcontrollore acceso oppure in deep-sleep.

\begin{table}[htbp]
\centering
\caption{Consumo stimato del nodo completo nelle configurazioni di interesse, somma
dei valori di datasheet.}
\label{tab:consumi-nodi}
\small
\begin{tabular}{@{}lrr@{}}
\toprule
\textbf{Configurazione} & \textbf{Corrente} & \textbf{Potenza a 5~V} \\
\midrule
ESP32 + \ldb\ + WiFi (web UI) & 175--340~mA & 0,9--1,7~W \\
ESP32 + \ldb, WiFi spento (logger seriale) & 100--150~mA & 0,5--0,75~W \\
ESP32 + \ldc, WiFi spento & 70--120~mA & 0,35--0,6~W \\
ESP32 + PIR, WiFi spento & 20--70~mA & 0,1--0,35~W \\
ESP32 in deep-sleep + PIR, sveglia su interrupt & $\mathbf{0{,}06}$~\textbf{mA} & 0,3~mW \\
\bottomrule
\end{tabular}
\end{table}

L'ultima riga è il punto chiave, perché l'uscita digitale del PIR può risvegliare
l'ESP32 dal deep-sleep con un interrupt su GPIO. Il nodo può quindi restare in attesa a
costo praticamente nullo, con il PIR nel ruolo di guardiano.

\section{Autonomia a batteria}
\label{sec:autonomia}

Una cella litio-ione 18650 da 3000~mAh, con un regolatore all'85\,\% di rendimento, fornisce circa 2550~mAh utili, che divisi per la corrente media danno le autonomie della Tabella~\ref{tab:autonomia} e della Figura~\ref{fig:consumi}. Il calcolo non include l'autoscarica della cella.

\begin{table}[htbp]
\centering
\caption{Autonomia stimata con una cella 18650 da 3000~mAh.}
\label{tab:autonomia}
\small
\begin{tabular}{@{}lrr@{}}
\toprule
\textbf{Scenario} & \textbf{Corrente media} & \textbf{Autonomia} \\
\midrule
ESP32 + \ldb\ + WiFi (web UI) & 200~mA & 13~h \\
ESP32 + \ldb, WiFi spento & 130~mA & 20~h \\
ESP32 + PIR, WiFi spento & 25~mA & 4 giorni \\
ESP32 in deep-sleep + PIR & 0,06~mA & 4,9 anni \\
\bottomrule
\end{tabular}
\end{table}

\begin{figure}[htbp]
\centering
\includegraphics[width=0.9\linewidth]{Consumi}
\caption{Consumi dei componenti e autonomia dei nodi, in scala logaritmica.}
\label{fig:consumi}
\end{figure}

\section{Cosa indicano i consumi}
\label{sec:consumi-vincoli}

Con il radar sempre acceso la batteria dura meno di un giorno (Tabella~\ref{tab:autonomia}), mentre un nodo in deep-sleep svegliato dal PIR arriva a 4,9 anni, un limite superiore perché durante le lezioni il PIR lo sveglia di continuo. Accendendo il radar un minuto ogni cinque, con l'ESP32 in light-sleep negli altri quattro, la corrente media scende a circa 27~mA e l'autonomia supera i tre giorni, con una latenza fino a quattro minuti. Su questi valori si basa l'architettura delle conclusioni (Capitolo~\ref{chap:conclusioni}), in cui il PIR decide quando accendere il radar.

\chapter{Conclusioni e sviluppi futuri}
\label{chap:conclusioni}

La tesi ha confrontato il sensore PIR con due moduli radar mmWave a 24~GHz per il rilevamento della presenza di una persona, prendendo come scenario applicativo quello del progetto DIPME. Lo studio dei sensori e della loro documentazione ha chiarito il principio di funzionamento e i dati che ciascuno fornisce, dal singolo bit del PIR alla distanza e all'energia per gate del radar. Per il confronto è stata realizzata una piattaforma di acquisizione su ESP32, con un'interfaccia web servita dal nodo stesso che mostra i dati in tempo reale e li esporta in CSV.

La campagna sperimentale ha confermato che il PIR, per principio, non vede una persona immobile, nemmeno sotto il banco. A decidere se il PIR vede è il movimento, non la distanza. L'\ldb\ invece rileva la persona immobile in tutti gli scenari, attraversa legno, vetro, cartongesso e plastica, che bloccano il PIR, e non ha prodotto falsi positivi a stanza vuota. L'\ldc\ funziona nello stesso scenario, ma non localizza la persona ferma e richiede una taratura per installazione, ed è quindi meno adatto. Sulle grandezze del radar è stato poi costruito un indice di vitalità, calcolato a bordo del nodo, che distingue la persona immobile da quella in movimento, con soglie che dipendono dalla geometria in cui sono tarate.

Il radar consuma però circa mille volte più del PIR e, sempre acceso, esaurisce la batteria in meno di un giorno. La risposta alla domanda di partenza è quindi che il radar rileva la persona con maggiore affidabilità, ma deve affiancare il PIR, non sostituirlo. In tempo di pace il nodo resta in deep-sleep e il PIR lo sveglia quando rileva una presenza, come già avviene nel progetto DIPME~\cite{callisto2026dipme}. In emergenza si accende il radar, che rileva anche la persona ferma, e alternandolo con un ciclo di lavoro l'autonomia supera i tre giorni. In concreto, al dispositivo del progetto va aggiunto un sensore radar come l'\ldb, tarato all'installazione sotto l'arredo.

\section{Sviluppi futuri}
\label{sec:conclusioni-sviluppi}

Restano tuttavia diverse aree da esplorare. La prima è l'integrazione con la rete LoRa. Nel regime di emergenza il bit di presenza decide la cadenza dei messaggi del nodo~\cite{callisto2026dipme}. Il primo passo è far decidere quella cadenza al radar invece che al PIR, e aggiungere ai messaggi l'indice di vitalità, che occupa un solo byte, li renderebbe più informativi.

Un'altra riguarda la prova nel dimostratore reale, dove restano da verificare i falsi positivi con arredi metallici, il comportamento con più persone, il montaggio rispetto alla lamiera forata del banco, la separazione fra banchi adiacenti e il rilevamento immediato della persona quando il radar viene riacceso.

Infine, il firmware dell'\ldb\ tara già le proprie soglie di presenza sul rumore di fondo, su comando e con il campo libero~\cite{hilinkV244}. Una procedura analoga eseguita dal nodo all'installazione, con il banco vuoto e poi occupato, adatterebbe anche il fondo per gate e le soglie dell'indice di vitalità, che oggi dipendono dalla geometria in cui sono state tarate.

In sintesi, il radar mmWave colma il limite del PIR nello scenario per cui il progetto DIPME è nato, e la piattaforma realizzata è pronta a portare il confronto nelle aule reali.



\appendix
%\input{schema_elettrico}
%\input{Appendix1}
%\input{Appendix2}
%\nocite{*}
\printbibliography
\printindex

\chapter*{Ringraziamenti}
\label{chap:Ringraziamenti}

Al termine di questo percorso desidero ringraziare le persone che, in modi diversi, mi hanno accompagnato fino a qui.

Il primo ringraziamento va al mio relatore, il professor Massimo Callisto De Donato, per avermi proposto un lavoro legato a un progetto reale e per la disponibilità dimostrata durante la tesi.

Ringrazio la mia famiglia, che mi ha sempre sostenuto e incoraggiato. Senza di voi non sarei arrivato fin qui.

Grazie ai miei amici e colleghi, che hanno reso questi anni più leggeri e mi hanno regalato momenti che porterò sempre con me.

A tutti voi, grazie.


\end{document}